% Authority slice: EGA II ega2-1-fr.tex lines 1-4040, 173428 LF UTF-8 bytes, SHA-256 E6911F259382EBEFA590F25970FD3F47297E1EC78D8C4703AD1627DA0BA4C720.
\setcounter{section}{0}
\section{아핀 사상}
\label{section:II.1-ko}

이 절의 결과 대부분은 제I장 \S~1의 ``전역적'' 대응물이다. 따라서
본질적으로 새로운 것은 아니며, 이후를 위한 편리한 언어를 제공할 뿐이다.

\setcounter{subsection}{0}
\subsection{$S$-준스킴과 $\mathcal{O}_S$-대수}
\label{subsection:II.1.1-ko}

\begin{env}[1.1.1]
\label{II.1.1.1-ko}
$S$를 준스킴, $X$를 $S$-준스킴, $f:X\to S$를 그 구조 사상이라 하자.
(\hyperref[0.4.2.4-ko]{0, 4.2.4})에 따라 직상
$f_*(\mathcal{O}_X)$는 $\mathcal{O}_S$-대수이다. 혼동의 우려가 없으면
이를
\oldpage[II]{6}
$\mathcal{A}(X)$로 나타내기로 한다. $U$가 $S$의 열린집합이면
\[
  \mathcal{A}(f^{-1}(U))=\mathcal{A}(X)|U.
\]
마찬가지로 임의의 $\mathcal{O}_X$-가군 $\mathcal{F}$(각각
$\mathcal{O}_X$-대수 $\mathcal{B}$)에 대하여, 직상
$f_*(\mathcal{F})$(각각 $f_*(\mathcal{B})$)를
$\mathcal{A}(\mathcal{F})$(각각 $\mathcal{A}(\mathcal{B})$)로
나타내기로 한다. 이는 단순히 $\mathcal{O}_S$-가군(각각
$\mathcal{O}_S$-대수)일 뿐만 아니라 $\mathcal{A}(X)$-가군(각각
$\mathcal{A}(X)$-대수)이다.
\end{env}

\begin{env}[1.1.2]
\label{II.1.1.2-ko}
$Y$를 또 하나의 $S$-준스킴, $g:Y\to S$를 그 구조 사상,
$h:X\to Y$를 $S$-사상이라 하자. 그러면 다음 도식은 가환한다.
\[
  \xymatrix{
    X\ar[rr]^h\ar[rd]_f && Y\ar[ld]^g\\
    & S
  }
  \tag{1.1.2.1}\label{II.1.1.2.1-ko}
\]

정의상 $h=(\psi,\theta)$이고, 여기서
$\theta:\mathcal{O}_Y\to h_*(\mathcal{O}_X)=\psi_*(\mathcal{O}_X)$는
환의 층의 준동형이다. 이로부터
(\hyperref[0.4.2.2-ko]{0, 4.2.2})에 따라 다음
$\mathcal{O}_S$-대수 준동형을 얻는다.
\[
  g_*(\theta):g_*(\mathcal{O}_Y)\to
  g_*(h_*(\mathcal{O}_X))=f_*(\mathcal{O}_X),
\]
다시 말해 $\mathcal{O}_S$-대수 준동형
$\mathcal{A}(Y)\to\mathcal{A}(X)$를 얻으며, 이를
$\mathcal{A}(h)$로도 나타낸다. $h':Y\to Z$가 또 하나의
$S$-사상이면
$\mathcal{A}(h'\circ h)=\mathcal{A}(h)\circ\mathcal{A}(h')$임은
곧바로 알 수 있다. 따라서 $S$-준스킴의 범주 위에 정의되어
$\mathcal{O}_S$-대수의 범주에 값을 갖는 \emph{반변 함자}
$\mathcal{A}(X)$를 정의하였다.

이제 $\mathcal{F}$를 $\mathcal{O}_X$-가군, $\mathcal{G}$를
$\mathcal{O}_Y$-가군이라 하고, $u:\mathcal{G}\to\mathcal{F}$를
$h$-사상이라 하자. 다시 말해
(\hyperref[0.4.4.1-ko]{0, 4.4.1})에 따라 $u$는
$\mathcal{O}_Y$-가군 준동형
$\mathcal{G}\to h_*(\mathcal{F})$이다. 그러면
\[
  g_*(u):g_*(\mathcal{G})\to g_*(h_*(\mathcal{F}))=f_*(\mathcal{F})
\]
는 $\mathcal{O}_S$-가군 준동형
$\mathcal{A}(\mathcal{G})\to\mathcal{A}(\mathcal{F})$이며, 이를
$\mathcal{A}(u)$로 나타낸다. 더욱이 쌍
$(\mathcal{A}(h),\mathcal{A}(u))$는
$\mathcal{A}(Y)$-가군 $\mathcal{A}(\mathcal{G})$에서
$\mathcal{A}(X)$-가군 $\mathcal{A}(\mathcal{F})$로 가는
\emph{디-준동형}을 이룬다.
\end{env}

\begin{env}[1.1.3]
\label{II.1.1.3-ko}
$S$를 고정하면, $X$가 $S$-준스킴이고 $\mathcal{F}$가
$\mathcal{O}_X$-가군인 쌍 $(X,\mathcal{F})$들이 하나의
\emph{범주}를 이룬다고 볼 수 있다. 여기서
$(X,\mathcal{F})$에서 $(Y,\mathcal{G})$로 가는 \emph{사상}은 쌍
$(h,u)$로 정의한다. 이때 $h:X\to Y$는 $S$-사상이고
$u:\mathcal{G}\to\mathcal{F}$는 $h$-사상이다. 그러면
$(\mathcal{A}(X),\mathcal{A}(\mathcal{F}))$는
$\mathcal{O}_S$-대수와 그 대수 위의 가군으로 이루어진 쌍들을
대상으로 하고 디-준동형들을 사상으로 하는 범주에 값을 갖는
\emph{반변 함자}라고 할 수 있다.
\end{env}

\subsection{준스킴 위에서 아핀인 준스킴}
\label{subsection:II.1.2-ko}

\begin{env}[1.2.1]
\label{II.1.2.1-ko}
$X$를 $S$-준스킴, $f:X\to S$를 그 구조 사상이라 하자.
$S$의 아핀 열린집합들로 이루어진 덮개 $(S_\alpha)$가 존재하여,
모든 $\alpha$에 대하여 열린집합 $f^{-1}(S_\alpha)$ 위에 $X$가
유도하는 준스킴이 아핀이면, $X$가 \emph{$S$ 위에서 아핀}이라고
한다.
\end{env}

\begin{env}[1.2.2]
\label{II.1.2.2-ko}
$S$의 모든 닫힌 부분준스킴은 $S$ 위에서 아핀인
$S$-준스킴이다
(\hyperref[I.4.2.3-ko]{I, 4.2.3}과
\hyperref[I.4.2.4-ko]{4.2.4}).
\end{env}

\begin{env}[1.2.3]
\label{II.1.2.3-ko}
$S$ 위에서 아핀인 준스킴 $X$가 반드시 아핀 스킴인 것은 아니다. 예
$X=S$가 이를 보여 준다(\hyperref[II.1.2.2-ko]{1.2.2}). 한편 아핀 스킴
$X$가 $S$-준스킴이라 해도, $X$가 반드시 $S$ 위에서 아핀인
\oldpage[II]{7}
것은 아니다(\hyperref[II.1.3.3-ko]{1.3.3} 참조). 그러나 $S$가
\emph{스킴}이면, 아핀 스킴인 모든 $S$-준스킴은 $S$ 위에서 아핀임을
상기하자(\hyperref[I.5.5.10-ko]{I, 5.5.10}).
\end{env}

\begin{env}[1.2.4]
\label{II.1.2.4-ko}
$S$ 위에서 아핀인 모든 $S$-준스킴은 $S$ 위에서 분리되어 있다
(다시 말해 $S$-스킴이다).
\end{env}

이는 (\hyperref[I.5.5.5-ko]{I, 5.5.5})와
(\hyperref[I.5.5.8-ko]{I, 5.5.8})에서 곧바로 따른다.

\begin{env}[1.2.5]
\label{II.1.2.5-ko}
$X$를 $S$ 위에서 아핀인 $S$-스킴, $f:X\to S$를 그 구조 사상이라
하자. 모든 열린집합 $U\subset S$에 대하여, $f^{-1}(U)$는 $U$ 위에서
아핀이다.
\end{env}

정의 (\hyperref[II.1.2.1-ko]{1.2.1})에 따라,
$S=\operatorname{Spec}(A)$와 $X=\operatorname{Spec}(B)$가 아핀인
경우로 환원된다. 이때 $f=({}^a\varphi,\widetilde{\varphi})$이고,
여기서 $\varphi:A\to B$는 환 준동형이다. $g\in A$인 $D(g)$들이
$S$의 기저를 이루므로, $U=D(g)$인 경우로 환원된다. 그런데 이때
$f^{-1}(U)=D(\varphi(g))$임을 알고 있으므로
(\hyperref[I.1.2.2.2-ko]{I, 1.2.2.2}), 명제가 따른다.

\begin{env}[1.2.6]
\label{II.1.2.6-ko}
$X$를 $S$ 위에서 아핀인 $S$-스킴, $f:X\to S$를 그 구조 사상이라
하자. 임의의 준연접 $\mathcal{O}_X$-가군 $\mathcal{F}$에 대하여,
$f_*(\mathcal{F})$는 준연접 $\mathcal{O}_S$-가군이다.
\end{env}

(\hyperref[II.1.2.4-ko]{1.2.4})를 고려하면, 이는
(\hyperref[I.9.2.2-ko]{I, 9.2.2, a)})에서 따른다.

특히 $\mathcal{O}_S$-대수
$\mathcal{A}(X)=f_*(\mathcal{O}_X)$는 \emph{준연접}이다.

\begin{env}[1.2.7]
\label{II.1.2.7-ko}
$X$를 $S$ 위에서 아핀인 $S$-스킴이라 하자. 임의의 $S$-준스킴
$Y$에 대하여, 집합 $\operatorname{Hom}_S(Y,X)$에서 집합
$\operatorname{Hom}(\mathcal{A}(X),\mathcal{A}(Y))$로 가는 사상
$h\mapsto\mathcal{A}(h)$
(\hyperref[II.1.1.2-ko]{1.1.2})은 전단사이다.
\end{env}

$f:X\to S$와 $g:Y\to S$를 구조 사상이라 하자. 먼저
$S=\operatorname{Spec}(A)$와 $X=\operatorname{Spec}(B)$가
아핀이라고 하자. 임의의 $\mathcal{O}_S$-대수 준동형
$\omega:f_*(\mathcal{O}_X)\to g_*(\mathcal{O}_Y)$에 대하여,
$\mathcal{A}(h)=\omega$인 $S$-사상 $h:Y\to X$가 유일하게 존재함을
보여야 한다. 정의에 따라, 모든 열린집합 $U\subset S$에 대하여
$\omega$는 다음 $\Gamma(U,\mathcal{O}_S)$-대수 준동형을 정한다.
\[
  \omega_U=\Gamma(U,\omega):
  \Gamma(f^{-1}(U),\mathcal{O}_X)\to
  \Gamma(g^{-1}(U),\mathcal{O}_Y)
\]
특히 $U=S$일 때, 이는 $\Gamma(S,\mathcal{O}_S)$-대수 준동형
$\varphi:\Gamma(X,\mathcal{O}_X)\to\Gamma(Y,\mathcal{O}_Y)$를
준다. $X$가 아핀이므로, 이 준동형에는 유일하게 정해지는
$S$-사상 $h:Y\to X$가 대응한다
(\hyperref[I.2.2.4-ko]{I, 2.2.4}). 이제
$\mathcal{A}(h)=\omega$임을 보여야 한다. 다시 말해, $S$의 한 기저에
속하는 모든 열린집합 $U$에 대하여 $\omega_U$가 $h$의 제한인
$S$-사상 $g^{-1}(U)\to f^{-1}(U)$에 대응하는 대수 준동형
$\varphi_U$와 일치함을 보이면 된다. $\lambda\in A$인
$U=D(\lambda)$의 경우만 살피면 충분하다.\footnote{\textbf{번역 주.}
인쇄 원문은 $\lambda\in S$라고 쓰지만, $S=\operatorname{Spec}(A)$이고
$\rho:A\to B$이며 뒤에서 $\rho(\lambda)$를 사용하므로
EGA2-CANON-DISP-0001에 따라 $\lambda\in A$로 교정했다.} 이때
$f=({}^a\rho,\widetilde{\rho})$이고 $\rho:A\to B$가 환 준동형이면,
$f^{-1}(U)=D(\mu)$이며 여기서 $\mu=\rho(\lambda)$이다. 또한
$\Gamma(f^{-1}(U),\mathcal{O}_X)$는 분수환 $B_\mu$이다. 그런데 다음
도식은 가환한다.
\[
  \xymatrix{
    B\ar[r]^\varphi\ar[d] & \Gamma(Y,\mathcal{O}_Y)\ar[d]\\
    B_\mu\ar[r]_{\varphi_U} & \Gamma(g^{-1}(U),\mathcal{O}_Y)
  }
\]
$\varphi_U$를 $\omega_U$로 바꾼 유사한 도식도 가환한다. 따라서
분수환의 보편 성질
(\hyperref[0.1.2.4-ko]{0, 1.2.4})에 의해
$\varphi_U=\omega_U$이다.

이제 일반적인 경우를 보자. $(S_\alpha)$를 $S$의 아핀 열린집합들로
이루어진 덮개라 하자.
\oldpage[II]{8}
여기서 각 $f^{-1}(S_\alpha)$는 아핀이라고 하자. 그러면 임의의
$\mathcal{O}_S$-대수 준동형
$\omega:\mathcal{A}(X)\to\mathcal{A}(Y)$는 제한에 의해 다음
$\mathcal{O}_{S_\alpha}$-대수 준동형들의 족을 정한다.
\[
  \omega_\alpha:
  \mathcal{A}(f^{-1}(S_\alpha))\to
  \mathcal{A}(g^{-1}(S_\alpha))
\]
따라서 앞에서 보인 바에 의해 $S_\alpha$-사상들의 족
$h_\alpha:g^{-1}(S_\alpha)\to f^{-1}(S_\alpha)$가 정해진다. 이제
$S_\alpha\cap S_\beta$의 한 기저에 속하는 모든 아핀 열린집합 $U$에
대하여, $h_\alpha$와 $h_\beta$를 $g^{-1}(U)$에 제한한 것들이
일치함을 보이면 된다. 이는 명백하다. 실제로 앞에서 보인 바에 의해
이 두 제한은 모두 $\omega$의 제한인 준동형
$\mathcal{A}(X)|U\to\mathcal{A}(Y)|U$에 대응한다.

\begin{env}[1.2.8]
\label{II.1.2.8-ko}
$X$와 $Y$를 $S$ 위에서 아핀인 두 $S$-스킴이라 하자. $S$-사상
$h:Y\to X$가 동형사상이기 위한 필요충분조건은
$\mathcal{A}(h):\mathcal{A}(X)\to\mathcal{A}(Y)$가 동형사상인 것이다.
\end{env}

이는 (\hyperref[II.1.2.7-ko]{1.2.7})과 $\mathcal{A}(X)$의 함자성에서
곧바로 따른다.

\subsection{$\mathcal{O}_S$-대수에 대응하는 $S$ 위에서 아핀인 준스킴}
\label{subsection:II.1.3-ko}

\begin{env}[1.3.1]
\label{II.1.3.1-ko}
$S$를 준스킴이라 하자. 임의의 준연접 $\mathcal{O}_S$-대수
$\mathcal{B}$에 대하여, $\mathcal{A}(X)=\mathcal{B}$를 만족하는
$S$ 위에서 아핀인 준스킴 $X$가 존재하며, $X$는 유일한
$S$-동형까지 유일하게 결정된다.
\end{env}

유일성은 (\hyperref[II.1.2.8-ko]{1.2.8})에서 곧바로 따른다.
$X$의 존재를 증명하자. 임의의 아핀 열린집합 $U\subset S$에 대하여
$X_U$를 준스킴
$\operatorname{Spec}(\Gamma(U,\mathcal{B}))$라 하자.
$\Gamma(U,\mathcal{B})$는 $\Gamma(U,\mathcal{O}_S)$-대수이므로
$X_U$는 $S$-준스킴이다
(\hyperref[I.1.6.1-ko]{I, 1.6.1}). 또한 $\mathcal{B}$가
준연접이므로 $\mathcal{O}_S$-대수 $\mathcal{A}(X_U)$는
$\mathcal{B}|U$와 표준적으로 동일시된다
(\hyperref[I.1.3.7-ko]{I, 1.3.7},
\hyperref[I.1.3.13-ko]{1.3.13}과
\hyperref[I.1.6.3-ko]{1.6.3}).

$V$를 $S$의 두 번째 아핀 열린집합이라 하고, $f_U$가 구조 사상
$X_U\to S$를 나타낼 때 $X_{U,V}$를 $X_U$가
$f_U^{-1}(U\cap V)$ 위에 유도하는 준스킴이라 하자.
$X_{U,V}$와 $X_{V,U}$는 $U\cap V$ 위에서 아핀이고
(\hyperref[II.1.2.5-ko]{1.2.5}), 정의에 따라
$\mathcal{A}(X_{U,V})$와 $\mathcal{A}(X_{V,U})$는
$\mathcal{B}|(U\cap V)$와 표준적으로 동일시된다. 그러므로
(\hyperref[II.1.2.8-ko]{1.2.8})에 따라 표준 $S$-동형사상
$\theta_{U,V}:X_{V,U}\to X_{U,V}$가 존재한다. 더욱이 $W$를
$S$의 세 번째 아핀 열린집합이라 하고, 구조 사상들에 의해
$X_V$, $X_W$, $X_W$ 안에서 각각 얻는 $U\cap V\cap W$의
역상들로 $\theta_{U,V}$, $\theta_{V,W}$, $\theta_{U,W}$를
제한한 사상을 차례로 $\theta'_{U,V}$, $\theta'_{V,W}$,
$\theta'_{U,W}$라 하면
$\theta'_{U,V}\circ\theta'_{V,W}=\theta'_{U,W}$이다.

따라서 준스킴 $X$, 아핀 열린집합들로 이루어진 $X$의 덮개
$(T_U)$, 그리고 각 $U$에 대한 동형사상
$\varphi_U:X_U\to T_U$가 존재하여, $\varphi_U$는
$f_U^{-1}(U\cap V)$를 $T_U\cap T_V$ 위로 보내고
$\theta_{U,V}=\varphi_U^{-1}\circ\varphi_V$가 성립한다
(\hyperref[I.2.3.1-ko]{I, 2.3.1}). 사상
$g_U=f_U\circ\varphi_U^{-1}$는 $T_U$에 $S$-준스킴 구조를 주며,
$g_U$와 $g_V$는 $T_U\cap T_V$에서 일치한다. 따라서 $X$는
$S$-준스킴이다. 또한 정의상 $X$가 $S$ 위에서 아핀이고
$\mathcal{A}(T_U)=\mathcal{B}|U$임은 명백하므로
$\mathcal{A}(X)=\mathcal{B}$이다.

이렇게 정의한 $S$-스킴 $X$가
\emph{$\mathcal{O}_S$-대수 $\mathcal{B}$에 대응한다}고 하거나
$X$를 \emph{$\mathcal{B}$의 스펙트럼}이라 하며, 이를
$\operatorname{Spec}(\mathcal{B})$로 나타낸다.

\begin{env}[1.3.2]
\label{II.1.3.2-ko}
$X$를 $S$ 위에서 아핀인 준스킴, $f:X\to S$를 그 구조 사상이라
하자. 임의의 아핀 열린집합 $U\subset S$에 대하여
$f^{-1}(U)$ 위에 유도된 준스킴은 환
$\Gamma(U,\mathcal{A}(X))$의 아핀 스킴이다.
\end{env}

\oldpage[II]{9}
(\hyperref[II.1.2.6-ko]{1.2.6})과
(\hyperref[II.1.3.1-ko]{1.3.1})에 따라 $X$가 어떤
$\mathcal{O}_S$-대수에 대응한다고 가정해도 되므로, 이 따름정리는
(\hyperref[II.1.3.1-ko]{1.3.1})에서 서술한 $X$의 구성에서 따른다.

\begin{env}[1.3.3]
\label{II.1.3.3-ko}
$S$를 점 $0$을 이중화한 체 $K$ 위의 아핀 평면이라 하자
(\hyperref[I.5.5.11-ko]{I, 5.5.11}). 그 표기에서 $S$는 두 아핀
열린집합 $Y_1$, $Y_2$의 합집합이다
(\hyperref[I.5.5.11-ko]{I, 5.5.11}). $f$가 열린 몰입 사상
$Y_1\to S$이면
$f^{-1}(Y_2)=Y_1\cap Y_2$는 $Y_1$ 안의 아핀 열린집합이 아니다
(\emph{loc. cit.}). 따라서 이로써 아핀 스킴이지만 $S$ 위에서
아핀이 아닌 예를 얻는다.
\end{env}

\begin{env}[1.3.4]
\label{II.1.3.4-ko}
$S$를 아핀 스킴이라 하자. $S$-준스킴 $X$가 $S$ 위에서
아핀이기 위한 필요충분조건은 $X$가 아핀 스킴인 것이다.
\end{env}

\begin{env}[1.3.5]
\label{II.1.3.5-ko}
$S$를 준스킴, $X$를 $S$ 위에서 아핀인 준스킴, $Y$를
$X$-준스킴이라 하자. $Y$가 $S$ 위에서 아핀이기 위한
필요충분조건은 $Y$가 $X$ 위에서 아핀인 것이다.
\end{env}

곧바로 $S$가 아핀 스킴인 경우로 환원된다. 그러면 $X$도 아핀
스킴이다(\hyperref[II.1.3.4-ko]{1.3.4}). 이때 명제의 두 조건은
모두 $Y$가 아핀 스킴이라는 것을 나타낸다
(\hyperref[II.1.3.4-ko]{1.3.4}).

\begin{env}[1.3.6]
\label{II.1.3.6-ko}
$X$를 $S$ 위에서 아핀인 준스킴이라 하자.
(\hyperref[II.1.3.5-ko]{1.3.5})에 따라, $X$ 위에서 아핀인
준스킴 $Y$를 정하는 것은 $S$ 위에서 아핀인 준스킴 $Y$와
$S$-사상 $g:Y\to X$를 정하는 것과 같다. 다시 말해
(\hyperref[II.1.3.1-ko]{1.3.1}과
\hyperref[II.1.2.7-ko]{1.2.7}), 이는 준연접
$\mathcal{O}_S$-대수 $\mathcal{B}$와 $\mathcal{O}_S$-대수 준동형
$\mathcal{A}(X)\to\mathcal{B}$를 정하는 것과 같다(이를
$\mathcal{B}$에 $\mathcal{A}(X)$-대수 구조를 정하는 것으로 볼 수
있다). $f:X\to S$가 구조 사상이면, 이때
$\mathcal{B}=f_*(g_*(\mathcal{O}_Y))$이다.
\end{env}

\begin{env}[1.3.7]
\label{II.1.3.7-ko}
$X$를 $S$ 위에서 아핀인 준스킴이라 하자. $X$가 $S$ 위에서
유한형이기 위한 필요충분조건은 준연접 $\mathcal{O}_S$-대수
$\mathcal{A}(X)$가 유한 생성인 것이다
(\hyperref[I.9.6.2-ko]{I, 9.6.2}).
\end{env}

정의 (\hyperref[I.9.6.2-ko]{I, 9.6.2})에 따라 $S$가 아핀인
경우로 환원된다. 그러면 $X$는 아핀 스킴이고
(\hyperref[II.1.3.4-ko]{1.3.4}), $S=\operatorname{Spec}(A)$,
$X=\operatorname{Spec}(B)$이면 $\mathcal{A}(X)$는
$\mathcal{O}_S$-대수 $\widetilde{B}$이다. 그런데
$\Gamma(S,\widetilde{B})=B$이므로,\footnote{\textbf{번역 주.} 인쇄
원문은 이 증명에서 정의되지 않은 $U$를 쓰지만,
EGA2-CANON-DISP-0030에 따라 $S$로 교정했다.} 이 따름정리는
(\hyperref[I.9.6.2-ko]{I, 9.6.2})와
(\hyperref[I.6.3.3-ko]{I, 6.3.3})에서 따른다.

\begin{env}[1.3.8]
\label{II.1.3.8-ko}
$X$를 $S$ 위에서 아핀인 준스킴이라 하자. $X$가 축소이기 위한
필요충분조건은 준연접 $\mathcal{O}_S$-대수 $\mathcal{A}(X)$가
축소인 것이다
(\hyperref[0.4.1.4-ko]{0, 4.1.4}).
\end{env}

실제로 이 문제는 명백히 $S$ 위에서 국소적이다.
(\hyperref[II.1.3.2-ko]{1.3.2})에 따라, 이 따름정리는
(\hyperref[I.5.1.1-ko]{I, 5.1.1})과
(\hyperref[I.5.1.4-ko]{I, 5.1.4})에서 따른다.

\subsection{$S$ 위에서 아핀인 준스킴 위의 준연접층}
\label{subsection:II.1.4-ko}

\begin{env}[1.4.1]
\label{II.1.4.1-ko}
$X$를 $S$ 위에서 아핀인 준스킴, $Y$를 $S$-준스킴,
$\mathcal{F}$(각각 $\mathcal{G}$)를 준연접
$\mathcal{O}_X$-가군(각각 $\mathcal{O}_Y$-가군)이라 하자. 그러면
$(h,u)\mapsto(\mathcal{A}(h),\mathcal{A}(u))$로 주어지는, 쌍
$(Y,\mathcal{G})$에서 $(X,\mathcal{F})$로 가는 사상들의 집합에서 쌍
$(\mathcal{A}(X),\mathcal{A}(\mathcal{F}))$에서
$(\mathcal{A}(Y),\mathcal{A}(\mathcal{G}))$로 가는 디-준동형들의
집합으로 가는 사상
(\hyperref[II.1.1.2-ko]{1.1.2}와
\hyperref[II.1.1.3-ko]{1.1.3})은 전단사이다.
\end{env}

증명은 (\hyperref[I.2.2.5-ko]{I, 2.2.5})와
(\hyperref[I.2.2.4-ko]{I, 2.2.4})를 사용하여
(\hyperref[II.1.2.7-ko]{1.2.7})의 증명과 정확히 같은 절차를 따르며,
세부사항은 독자에게 맡긴다.

\begin{env}[1.4.2]
\label{II.1.4.2-ko}
(\hyperref[II.1.4.1-ko]{1.4.1})의 가정들에 더하여 $Y$가 $S$ 위에서
아핀이라고 가정하면, $(h,u)$가 동형이기 위한 필요충분조건은
$(\mathcal{A}(h),\mathcal{A}(u))$가 디-동형인 것이다.
\end{env}

\oldpage[II]{10}
\begin{env}[1.4.3]
\label{II.1.4.3-ko}
준연접 $\mathcal{O}_S$-대수 $\mathcal{B}$와 준연접
$\mathcal{B}$-가군 $\mathcal{M}$으로 이루어진 모든 쌍
$(\mathcal{B},\mathcal{M})$에 대하여
($\mathcal{M}$을 $\mathcal{O}_S$-가군으로 보든
$\mathcal{B}$-가군으로 보든 준연접성은 같은 조건이다
(\hyperref[I.9.6.1-ko]{I, 9.6.1})), $S$ 위에서 아핀인 준스킴 $X$와
준연접 $\mathcal{O}_X$-가군 $\mathcal{F}$로 이루어진 쌍
$(X,\mathcal{F})$가 존재하여
$\mathcal{A}(X)=\mathcal{B}$이고
$\mathcal{A}(\mathcal{F})=\mathcal{M}$이다. 더욱이 이 쌍은 유일한
동형까지 유일하게 결정된다.
\end{env}

유일성은 (\hyperref[II.1.4.1-ko]{1.4.1})과
(\hyperref[II.1.4.2-ko]{1.4.2})에서 따른다. 존재성은
(\hyperref[II.1.3.1-ko]{1.3.1})에서와 같이 증명되며, 여기서도
세부사항은 독자에게 맡긴다. $\mathcal{O}_X$-가군 $\mathcal{F}$를
$\widetilde{\mathcal{M}}$로 나타내고, 이 가군을 준연접
$\mathcal{B}$-가군 $\mathcal{M}$에 \emph{대응하는} 가군이라고 한다.
모든 아핀 열린집합 $U\subset S$에 대하여
$\widetilde{\mathcal{M}}|p^{-1}(U)$는
($p$가 구조 사상 $X\to S$를 나타낼 때)
$(\Gamma(U,\mathcal{M}))^{\sim}$와 표준적으로 동형이다.

\begin{env}[1.4.4]
\label{II.1.4.4-ko}
준연접 $\mathcal{B}$-가군들의 범주에서
$\widetilde{\mathcal{M}}$는 $\mathcal{M}$에 관한 공변 가법 완전
함자이고, 귀납극한과 직합과 가환한다.
\end{env}

실제로 곧 $S$가 아핀인 경우로 환원되며, 그 경우 따름정리는
(\hyperref[I.1.3.5-ko]{I, 1.3.5}),
(\hyperref[I.1.3.9-ko]{I, 1.3.9})와
(\hyperref[I.1.3.11-ko]{I, 1.3.11})에서 따른다.

\begin{env}[1.4.5]
\label{II.1.4.5-ko}
(\hyperref[II.1.4.3-ko]{1.4.3})의 가정 아래에서
$\widetilde{\mathcal{M}}$가 유한 생성 $\mathcal{O}_X$-가군이기 위한
필요충분조건은 $\mathcal{M}$이 유한 생성 $\mathcal{B}$-가군인 것이다.
\end{env}

곧 $S=\operatorname{Spec}(A)$가 아핀 스킴인 경우로 환원된다. 이때
$\mathcal{B}=\widetilde{B}$이고, 여기서 $B$는 $A$-대수이다.\footnote{
\textbf{번역 주.} 인쇄 원문은 여기서 $B$가 유한 생성 $A$-대수라고
하고 (I, 9.6.2)를 인용하지만, 이는 명제의 가정에 없으므로
EGA2-CANON-DISP-0031에 따라 그 구절과 인용을 삭제했다.} 또한
$\mathcal{M}=\widetilde{M}$이고, 여기서 $M$은 $B$-가군이다
(\hyperref[I.1.3.13-ko]{I, 1.3.13}). \emph{준스킴 $X$ 위에서}
$\mathcal{O}_X$는 환 $B$에 대응하고
$\widetilde{\mathcal{M}}$는 $B$-가군 $M$에 대응한다. 따라서
$\widetilde{\mathcal{M}}$가 유한 생성이기 위한 필요충분조건은 $M$이
유한 생성인 것이다(\hyperref[I.1.3.13-ko]{I, 1.3.13}). 이로써
주장이 증명된다.

\begin{env}[1.4.6]
\label{II.1.4.6-ko}
$Y$를 $S$ 위에서 아핀인 준스킴, $X$와 $X'$을 $Y$ 위에서 아핀인
두 준스킴이라 하자(따라서 이들은 $S$ 위에서도 아핀이다
(\hyperref[II.1.3.5-ko]{1.3.5})). 다음과 같이 놓자.
$\mathcal{B}=\mathcal{A}(Y)$, $\mathcal{A}=\mathcal{A}(X)$,
$\mathcal{A}'=\mathcal{A}(X')$. 그러면 $X\times_Y X'$은 $Y$ 위에서
아핀이고(따라서 $S$ 위에서도 아핀이다),
$\mathcal{A}(X\times_Y X')$는
$\mathcal{A}\otimes_{\mathcal{B}}\mathcal{A}'$과 표준적으로
동일시된다.
\end{env}

실제로 (\hyperref[I.9.6.1-ko]{I, 9.6.1})에 따라
$\mathcal{A}\otimes_{\mathcal{B}}\mathcal{A}'$은 준연접
$\mathcal{B}$-대수이고, 따라서 준연접 $\mathcal{O}_S$-대수이기도
하다(\hyperref[I.9.6.1-ko]{I, 9.6.1}).
$Z$를 $\mathcal{A}\otimes_{\mathcal{B}}\mathcal{A}'$의 스펙트럼이라
하자(\hyperref[II.1.3.1-ko]{1.3.1}). 표준
$\mathcal{B}$-대수 준동형
$\mathcal{A}\to\mathcal{A}\otimes_{\mathcal{B}}\mathcal{A}'$과
$\mathcal{A}'\to\mathcal{A}\otimes_{\mathcal{B}}\mathcal{A}'$은
(\hyperref[II.1.2.7-ko]{1.2.7})에 따라 $Y$-사상 $p:Z\to X$와
$p':Z\to X'$에 대응한다. 삼중항 $(Z,p,p')$가 곱
$X\times_Y X'$임을 보이기 위해, $S$가 환 $C$의 아핀 스킴인
경우로 환원할 수 있다(\hyperref[I.3.2.6.4-ko]{I, 3.2.6.4}).
그러면 $Y$, $X$, $X'$은 아핀 스킴이고
(\hyperref[II.1.3.4-ko]{1.3.4}), 그 환 $B$, $A$, $A'$은
$C$-대수이며
$\mathcal{B}=\widetilde{B}$, $\mathcal{A}=\widetilde{A}$,
$\mathcal{A}'=\widetilde{A'}$이다.
(\hyperref[I.1.3.13-ko]{I, 1.3.13})에 따라
$\mathcal{A}\otimes_{\mathcal{B}}\mathcal{A}'$은
$\mathcal{O}_S$-대수 $\widetilde{A\otimes_B A'}$와 표준적으로
동일시된다. 따라서 환 $A(Z)$는 $A\otimes_B A'$과 동일시되고,
사상 $p$, $p'$은 표준 준동형
$A\to A\otimes_B A'$, $A'\to A\otimes_B A'$에 대응한다.
이제 명제는 (\hyperref[I.3.2.2-ko]{I, 3.2.2})에서 따른다.

\begin{env}[1.4.7]
\label{II.1.4.7-ko}
$\mathcal{F}$(각각 $\mathcal{F}'$)를 준연접
$\mathcal{O}_X$-가군(각각 $\mathcal{O}_{X'}$-가군)이라 하자. 그러면
$\mathcal{A}(\mathcal{F}\otimes_Y\mathcal{F}')$는
$\mathcal{A}(\mathcal{F})\otimes_{\mathcal{A}(Y)}
\mathcal{A}(\mathcal{F}')$과 표준적으로 동일시된다.
\end{env}

$\mathcal{F}\otimes_Y\mathcal{F}'$은 $X\times_Y X'$ 위에서
준연접이다(\hyperref[I.9.1.2-ko]{I, 9.1.2}).
$g:Y\to S$, $f:X\to Y$, $f':X'\to Y$를 구조 사상이라 하자.
그러면 구조 사상
\oldpage[II]{11}
$h:Z\to S$는 $g\circ f\circ p$ 및 $g\circ f'\circ p'$과 같다.
표준 준동형
\[
  \mathcal{A}(\mathcal{F})\otimes_{\mathcal{A}(Y)}\mathcal{A}(\mathcal{F}')
  \longrightarrow \mathcal{A}(\mathcal{F}\otimes_Y\mathcal{F}')
\]
을 다음과 같이 정의한다. 임의의 열린집합 $U\subset S$에 대하여
표준 준동형
\[
  \begin{aligned}
  \Gamma(f^{-1}(g^{-1}(U)),\mathcal{F})
  &\longrightarrow \Gamma(h^{-1}(U),p^*(\mathcal{F}))\quad\text{및}\quad
  \Gamma(f^{\prime-1}(g^{-1}(U)),\mathcal{F}')
  \longrightarrow \Gamma(h^{-1}(U),p^{\prime*}(\mathcal{F}'))
  \end{aligned}
\]
들이 있다(\hyperref[0.4.4.3-ko]{0, 4.4.3}). 따라서 표준 준동형
\[
  \begin{aligned}
  &\Gamma(f^{-1}(g^{-1}(U)),\mathcal{F})
    \otimes_{\Gamma(g^{-1}(U),\mathcal{O}_Y)}
    \Gamma(f^{\prime-1}(g^{-1}(U)),\mathcal{F}')\longrightarrow\\
  &\hspace{17em}
    \Gamma(h^{-1}(U),p^*(\mathcal{F}))
    \otimes_{\Gamma(h^{-1}(U),\mathcal{O}_Z)}
    \Gamma(h^{-1}(U),p^{\prime*}(\mathcal{F}'))
  \end{aligned}
\]
을 얻는다.

이것이 실제로 $\mathcal{A}(Z)$-가군의 동형임을 보이기 위해,
$S$(따라서 $X$, $X'$, $Y$, $X\times_Y X'$도)가 아핀 스킴인
경우로 환원할 수 있다. 이때
(\hyperref[II.1.4.6-ko]{1.4.6})의 표기를 사용하여
$\mathcal{F}=\widetilde{M}$, $\mathcal{F}'=\widetilde{M'}$라 쓸 수
있으며, $M$(각각 $M'$)은 $A$-가군(각각 $A'$-가군)이다. 그러면
$\mathcal{F}\otimes_Y\mathcal{F}'$은 $X\times_Y X'$ 위에서
$(A\otimes_B A')$-가군 $M\otimes_B M'$에 대응하는 층과
동일시된다(\hyperref[I.9.1.3-ko]{I, 9.1.3}). 따라서 따름정리는
$\mathcal{O}_S$-가군 $\widetilde{M\otimes_B M'}$와
$\widetilde{M}\otimes_{\widetilde{B}}\widetilde{M'}$의 표준적
동일시에서 따른다. 여기서 $M$, $M'$, $B$는 $C$-가군으로 본다
(\hyperref[I.1.3.13-ko]{I, 1.3.13}과
\hyperref[I.1.6.3-ko]{I, 1.6.3}).

특히 (\hyperref[II.1.4.7-ko]{1.4.7})을 $X=Y$이고
$\mathcal{F}'=\mathcal{O}_{X'}$인 경우에 적용하면,
$\mathcal{A}'$-가군 $\mathcal{A}(f^{\prime*}(\mathcal{F}))$가
$\mathcal{A}(\mathcal{F})\otimes_{\mathcal{B}}\mathcal{A}'$과
동일시됨을 알 수 있다.

\begin{env}[1.4.8]
\label{II.1.4.8-ko}
특히 $X=X'=Y$일 때($X$는 $S$ 위에서 아핀이다),
$\mathcal{F}$와 $\mathcal{G}$가 두 준연접
$\mathcal{O}_X$-가군이면
\[
  \mathcal{A}(\mathcal{F}\otimes_{\mathcal{O}_X}\mathcal{G})
  =\mathcal{A}(\mathcal{F})\otimes_{\mathcal{A}(X)}\mathcal{A}(\mathcal{G})
  \tag{1.4.8.1}\label{II.1.4.8.1-ko}
\]
이라는 표준적이고 함자적인 동형이 있다. 더욱이 $\mathcal{F}$가
유한 표시를 가지면
(\hyperref[I.1.6.3-ko]{I, 1.6.3}과
\hyperref[I.1.3.12-ko]{I, 1.3.12})에 따라
\[
  \mathcal{A}(\mathcal{H}om_X(\mathcal{F},\mathcal{G}))
  =\mathcal{H}om_{\mathcal{A}(X)}
    (\mathcal{A}(\mathcal{F}),\mathcal{A}(\mathcal{G}))
  \tag{1.4.8.2}\label{II.1.4.8.2-ko}
\]
이라는 표준적 동형이 있다.
\end{env}

\begin{env}[1.4.9]
\label{II.1.4.9-ko}
$X$, $X'$가 $S$ 위에서 아핀인 두 준스킴이면, 합
$X\sqcup X'$도 $S$ 위에서 아핀이다. 두 아핀 스킴의 합은 아핀
스킴이기 때문이다.
\end{env}

\begin{env}[1.4.10]
\label{II.1.4.10-ko}
$S$를 준스킴, $\mathcal{B}$를 준연접 $\mathcal{O}_S$-대수라 하고,
$X=\operatorname{Spec}(\mathcal{B})$라 하자. $\mathcal{J}$를
$\mathcal{B}$의 임의의 준연접 아이디얼층이라 하자. 그러면
$\widetilde{\mathcal{J}}$는 $\mathcal{O}_X$의 준연접 아이디얼층이다.
$Y$를 $X$에서 $\widetilde{\mathcal{J}}$로 정의되는 닫힌 부분준스킴이라
하면, 이는 $\operatorname{Spec}(\mathcal{B}/\mathcal{J})$와 표준적으로
동형이다.
\end{env}

실제로 (\hyperref[I.4.2.3-ko]{I, 4.2.3})에서 $Y$가 $S$ 위에서
아핀임이 곧바로 따른다. 따라서 (\hyperref[II.1.3.1-ko]{1.3.1})에 의해
$S$가 아핀인 경우로 환원되고, 이 경우 명제는
(\hyperref[I.4.1.2-ko]{I, 4.1.2})에서 곧바로 따른다.

(\hyperref[II.1.4.10-ko]{1.4.10})의 결과는 다음과 같이도 표현할 수
있다. $h:\mathcal{B}\to\mathcal{B}'$가 준연접 $\mathcal{O}_S$-대수들의
\emph{전사} 준동형이면,
${}^a h:\operatorname{Spec}(\mathcal{B}')\to
\operatorname{Spec}(\mathcal{B})$는 \emph{닫힌 몰입 사상}이다.

\oldpage[II]{12}
\begin{env}[1.4.11]
\label{II.1.4.11-ko}
$S$를 준스킴, $\mathcal{B}$를 준연접 $\mathcal{O}_S$-대수라 하고,
$X=\operatorname{Spec}(\mathcal{B})$라 하자. $\mathcal{K}$를
$\mathcal{O}_S$의 임의의 준연접 아이디얼층이라 하자. 그러면
($f$가 구조 사상 $X\to S$를 나타낼 때)
$f^*(\mathcal{K})\mathcal{O}_X=\widetilde{\mathcal{K}\mathcal{B}}$가
표준 동형까지 성립한다.
\end{env}

이 문제는 $S$ 위에서 국소적이므로 $S=\operatorname{Spec}(A)$가 아핀인
경우로 환원된다. 이 경우 명제는
(\hyperref[I.1.6.9-ko]{I, 1.6.9})에 다름 아니다.

\subsection{밑준스킴의 변경}
\label{II.1.5-ko}

\begin{env}[1.5.1]
\label{II.1.5.1-ko}
$X$를 $S$ 위에서 아핀인 준스킴이라 하자. 밑준스킴의 임의의 확장
$g:S'\to S$에 대하여 $X'=X_{(S')}=X\times_S S'$은 $S'$ 위에서
아핀이다.
\end{env}

$f'$을 사영 $X'\to S'$이라 하자. $g(U')$가 $S$의 어떤 아핀
열린집합 $U$에 포함되는 $S'$의 모든 아핀 열린집합 $U'$에 대하여
$f^{\prime-1}(U')$가 아핀 열린집합임을 보이면 충분하다
(\hyperref[II.1.2.1-ko]{1.2.1}). 따라서 $S$와 $S'$가 아핀인 경우로
한정할 수 있고, 이 경우 $X'$이 아핀 스킴임을 보이면 충분하다
(\hyperref[II.1.3.4-ko]{1.3.4}). 그런데 이때
(\hyperref[II.1.3.4-ko]{1.3.4}) $X$는 아핀 스킴이다. $S$, $S'$, $X$의
환을 각각 $A$, $A'$, $B$라 하면, $X'$은 환 $A'\otimes_A B$의 아핀
스킴임을 안다(\hyperref[I.3.2.2-ko]{I, 3.2.2}).

\begin{env}[1.5.2]
\label{II.1.5.2-ko}
(\hyperref[II.1.5.1-ko]{1.5.1})의 가정 아래에서 $f:X\to S$를 구조
사상, $f':X'\to S'$과 $g':X'\to X$를 사영이라 하자. 그러면 다음
도식은 가환한다.
\[
  \xymatrix{
    X\ar[d] & X'\ar[l]_{g'}\ar[d]^{f'}\\
    S & S'\ar[l]_g
  }
\]
모든 준연접 $\mathcal{O}_X$-가군 $\mathcal{F}$에 대하여 다음 표준
$\mathcal{O}_{S'}$-가군 동형이 존재한다.
\[
  u:g^*(f_*(\mathcal{F}))\xrightarrow{\sim}
  f'_*(g^{\prime*}(\mathcal{F})).
  \tag{1.5.2.1}\label{II.1.5.2.1-ko}
\]
특히 $\mathcal{A}(X')$에서 $g^*(\mathcal{A}(X))$로 가는 표준 동형이
존재한다.
\end{env}

$u$를 정의하려면 준동형
\[
  v:f_*(\mathcal{F})\longrightarrow
  g_*(f'_*(g^{\prime*}(\mathcal{F})))
  =f_*(g'_*(g^{\prime*}(\mathcal{F})))
\]
를 정의하고 $u=v^\sharp$로 두면 충분하다
(\hyperref[0.4.4.3-ko]{0, 4.4.3}). $v=f_*(\rho)$로 두겠다. 여기서
$\rho$는 표준 준동형
$\mathcal{F}\to g'_*(g^{\prime*}(\mathcal{F}))$이다
(\hyperref[0.4.4.3-ko]{0, 4.4.3}). $u$가 동형임을 보이기 위해서는
$S$와 $S'$, 따라서 $X$와 $X'$도 아핀인 경우로 한정할 수 있다.
(\hyperref[II.1.5.1-ko]{1.5.1})의 표기를 사용하면 이때
$\mathcal{F}=\widetilde{M}$이고, 여기서 $M$은 $B$-가군이다. 그러면
$g^*(f_*(\mathcal{F}))$와 $f'_*(g^{\prime*}(\mathcal{F}))$는 둘 다
$A'$-가군 $A'\otimes_A M$에 대응하는 $\mathcal{O}_{S'}$-가군과 같다
(여기서 $M$은 $A$-가군으로 본다). 또한 $u$는 항등사상에 대응하는
준동형임을 곧바로 알 수 있다
(\hyperref[I.1.6.3-ko]{I, 1.6.3},
\hyperref[I.1.6.5-ko]{1.6.5}와 \hyperref[I.1.6.7-ko]{1.6.7}).

\begin{env}[1.5.3]
\label{II.1.5.3-ko}
$X$가 더 이상 $S$ 위에서 아핀이라고 가정되지 않을 때에는
(\hyperref[II.1.5.2-ko]{1.5.2})가 여전히 성립한다고 생각해서는 안
된다. 이는 심지어 $S'=\operatorname{Spec}(k(s))$ ($s\in S$)이고
$S'\to S$가 표준 사상인 경우에도 마찬가지이다
(\hyperref[I.2.4.5-ko]{I, 2.4.5}) --- 이 경우 $X'$은 바로
\emph{섬유} $f^{-1}(s)$이다(\hyperref[I.3.6.2-ko]{I, 3.6.2}). 다시
말해 $X$가 $S$ 위에서 아핀이 아닐 때에는
\oldpage[II]{13}
``준연접층의 직상'' 연산은 ``섬유 취하기'' 연산과 교환하지 않는다.
그러나 제III장(\agref{III.4.2.4-ko}{III, 4.2.4})에서 $f$가 고유
사상이고(\agref{II.5.4-ko}{5.4}) $S$가 뇌터일 때 $X$ 위의
\emph{연접층}에 대하여 성립하는, 이 방향의 ``점근적'' 성격의 결과를
보게 될 것이다.
\end{env}

\begin{env}[1.5.4]
\label{II.1.5.4-ko}
$S$ 위에서 아핀인 모든 준스킴 $X$와 모든 $s\in S$에 대하여, 섬유
$f^{-1}(s)$는 아핀 스킴이다(여기서 $f$는 구조 사상 $X\to S$를
나타낸다).
\end{env}

(\hyperref[II.1.5.1-ko]{1.5.1})을 $S'=\operatorname{Spec}(k(s))$에
적용하고 (\hyperref[II.1.3.4-ko]{1.3.4})을 사용하면 충분하다.

\begin{corollary}[1.5.5]
\label{II.1.5.5-ko}
$X$를 $S$-준스킴, $S'$을 $S$ 위에서 아핀인 준스킴이라 하자. 그러면
$X'=X_{(S')}$은 $X$ 위에서 아핀인 준스킴이다. 또한 $f:X\to S$가
구조 사상이면 다음 표준 $\mathcal{O}_X$-대수 동형이 존재한다.
$\mathcal{A}(X')\xrightarrow{\sim}f^*(\mathcal{A}(S'))$.
그리고 모든 준연접 $\mathcal{A}(S')$-가군 $\mathcal{M}$에 대하여 다음
표준 디-동형이 존재한다.
$f^*(\mathcal{M})\xrightarrow{\sim}
\mathcal{A}(f^{\prime*}(\widetilde{\mathcal{M}}))$.
여기서 $f'=f_{(S')}$은 구조 사상 $X'\to S'$을 나타낸다.
\end{corollary}

(\hyperref[II.1.5.1-ko]{1.5.1})과
(\hyperref[II.1.5.2-ko]{1.5.2})에서 $X$와 $S'$의 역할을 서로 바꾸면
충분하다.

\begin{env}[1.5.6]
\label{II.1.5.6-ko}
이제 $S$, $S'$을 두 준스킴, $q:S'\to S$를 사상,
$\mathcal{B}$(각각 $\mathcal{B}'$)를 준연접
$\mathcal{O}_S$-대수(각각 준연접 $\mathcal{O}_{S'}$-대수),
$u:\mathcal{B}\to\mathcal{B}'$를 $q$-사상, 즉 $\mathcal{O}_S$-대수
준동형 $\mathcal{B}\to q_*(\mathcal{B}')$이라 하자.
$X=\operatorname{Spec}(\mathcal{B})$,
$X'=\operatorname{Spec}(\mathcal{B}')$라 하면, 다음 사상을 표준적으로
얻는다.
\[
  v=\operatorname{Spec}(u):X'\to X
\]
이때 도식
\[
  \xymatrix{
    X'\ar[r]^v\ar[d] & X\ar[d]\\
    S'\ar[r]_q & S
  }
  \tag{1.5.6.1}\label{II.1.5.6.1-ko}
\]
은 가환한다(수직 화살표들은 구조 사상이다). 실제로 $u$를 주는 것은
준연접 $\mathcal{O}_{S'}$-대수들의 준동형
$u^\sharp:q^*(\mathcal{B})\to\mathcal{B}'$를 주는 것과 동치이다
(\hyperref[0.4.4.3-ko]{0, 4.4.3}). 따라서 이는
$\mathcal{A}(w)=u^\sharp$인 표준 $S'$-사상
\[
  w:\operatorname{Spec}(\mathcal{B}')\to
  \operatorname{Spec}(q^*(\mathcal{B}))
\]
을 정의한다(\hyperref[II.1.2.7-ko]{1.2.7}). 한편
(\hyperref[II.1.5.2-ko]{1.5.2})에 따라
$\operatorname{Spec}(q^*(\mathcal{B}))$는 $X\times_S S'$과 표준적으로
동일시된다. 사상 $v$는 $w$에 이어 첫째 사영을 취한 합성
$X'\xrightarrow{w}X\times_S S'\xrightarrow{p_1}X$이고,
(\hyperref[II.1.5.6.1-ko]{1.5.6.1})의 가환성은 정의에서 따른다.
$U$(각각 $U'$)를 $q(U')\subset U$인 $S$(각각 $S'$)의 아핀
열린집합이라 하고,
$A=\Gamma(U,\mathcal{O}_S)$,
$A'=\Gamma(U',\mathcal{O}_{S'})$를 그 환들,
$B=\Gamma(U,\mathcal{B})$,
$B'=\Gamma(U',\mathcal{B}')$라 하자. $u$의 제한인
$(q|U')$-사상 $\mathcal{B}|U\to\mathcal{B}'|U'$은 대수의
디-준동형 $B\to B'$에 대응한다. $V$와 $V'$을 각각 구조 사상
$X\to S$와 $X'\to S'$에 의한 $U$와 $U'$의 역상이라 하면, $v$의 제한인 사상
$V'\to V$는 (\hyperref[I.1.7.3-ko]{I, 1.7.3})에 따라 앞의
디-준동형에 대응한다.
\end{env}

\begin{env}[1.5.7]
\label{II.1.5.7-ko}
(\hyperref[II.1.5.6-ko]{1.5.6})과 같은 가정 아래에서 $\mathcal{M}$을
준연접 $\mathcal{B}$-가군이라 하자. 그러면 다음 표준
$\mathcal{O}_{X'}$-가군 동형이 존재한다.
\[
  v^*(\widetilde{\mathcal{M}})\xrightarrow{\sim}
  \bigl(q^*(\mathcal{M})\otimes_{q^*(\mathcal{B})}\mathcal{B}'\bigr)^{\sim}
  \tag{1.5.7.1}\label{II.1.5.7.1-ko}
\]

\oldpage[II]{14}
실제로 (\hyperref[II.1.5.2.1-ko]{1.5.2.1})의 표준 동형은
$p_1^*(\widetilde{\mathcal{M}})$에서
$\operatorname{Spec}(q^*(\mathcal{B}))$ 위의
$q^*(\mathcal{B})$-가군 $q^*(\mathcal{M})$에 대응하는 층으로 가는
표준 동형을 준다. 이제 (\hyperref[II.1.4.7-ko]{1.4.7})을 적용하면
충분하다.
\end{env}

\subsection{아핀 사상}
\label{II.1.6-ko}

\begin{env}[1.6.1]
\label{II.1.6.1-ko}
준스킴의 사상 $f:X\to Y$를 구조 사상으로 보았을 때 $X$가 $Y$ 위에서
아핀인 준스킴이면, $f$를 \emph{아핀} 사상이라 한다. 한 준스킴 위에서
아핀인 준스킴들의 성질은 이 용어로 다음과 같이 표현된다.
\end{env}

\begin{proposition}[1.6.2]
\label{II.1.6.2-ko}
\begin{enumerate}
  \item[\rm(i)] 닫힌 몰입 사상은 아핀 사상이다.
  \item[\rm(ii)] 두 아핀 사상의 합성은 아핀 사상이다.
  \item[\rm(iii)] $f:X\to Y$가 아핀 $S$-사상이면, 밑준스킴의 모든 확장
    $S'\to S$에 대하여 $f_{(S')}:X_{(S')}\to Y_{(S')}$도 아핀 사상이다.
  \item[\rm(iv)] $f:X\to Y$와 $f':X'\to Y'$이 두 아핀 $S$-사상이면,
    \[
      f\times_S f':X\times_S X'\to Y\times_S Y'
    \]
    도 아핀 사상이다.
  \item[\rm(v)] $f:X\to Y$와 $g:Y\to Z$가 두 사상이고 $g\circ f$가
    아핀 사상이며 $g$가 분리 사상이면, $f$는 아핀 사상이다.
  \item[\rm(vi)] $f$가 아핀 사상이면 $f_{\mathrm{red}}$도 아핀 사상이다.
\end{enumerate}
\end{proposition}

(\hyperref[I.5.5.12-ko]{I, 5.5.12})에 따라 (i), (ii), (iii)만 증명하면
충분하다. 그런데 (i)는 (\hyperref[II.1.2.2-ko]{1.2.2})의 예이고,
(ii)는 (\hyperref[II.1.3.5-ko]{1.3.5})이며, 마지막으로 (iii)은
$X_{(S')}$이 곱 $X\times_Y Y_{(S')}$과 동일시되므로
(\hyperref[II.1.5.1-ko]{1.5.1})에서 따른다
(\hyperref[I.3.3.11-ko]{I, 3.3.11}).

\begin{corollary}[1.6.3]
\label{II.1.6.3-ko}
$X$가 아핀 스킴이고 $Y$가 스킴이면, 모든 사상 $X\to Y$는 아핀
사상이다.
\end{corollary}

\begin{proposition}[1.6.4]
\label{II.1.6.4-ko}
$Y$를 국소 뇌터 준스킴, $f:X\to Y$를 유한형 사상이라 하자. $f$가
아핀 사상이기 위한 필요충분조건은 $f_{\mathrm{red}}$가 아핀 사상인
것이다.
\end{proposition}

(\hyperref[II.1.6.2-ko]{1.6.2, (vi)})에 따라 조건의 충분성만 증명하면
된다. $Y$가 아핀이고 뇌터일 때 $X$가 아핀임을 보이면 충분하다. 이때
$Y_{\mathrm{red}}$는 아핀이므로, 가정에 따라 $X_{\mathrm{red}}$도
아핀이다. 한편 $X$는 뇌터이므로, 결론은
(\hyperref[I.6.1.7-ko]{I, 6.1.7})에서 따른다.

\subsection{가군의 층에 대응하는 벡터 다발}
\label{II.1.7-ko}

\begin{env}[1.7.1]
\label{II.1.7.1-ko}
$A$를 환, $E$를 $A$-가군이라 하자. $E$의 \emph{대칭 대수}라 하고
$\mathbf{S}(E)$(또는 $\mathbf{S}_A(E)$)로 나타내는 것은 텐서 대수
$\mathbf{T}(E)$를 $x,y\in E$에 대한 원소
$x\otimes y-y\otimes x$들이 생성하는 양쪽 아이디얼로 나눈 몫대수임을
상기하자. 대수 $\mathbf{S}(E)$는 다음 보편 성질로 특징지어진다.
$\sigma$를 표준 사상 $E\to\mathbf{S}(E)$, 즉
$E\to\mathbf{T}(E)$와 표준 사상
$\mathbf{T}(E)\to\mathbf{S}(E)$의 합성이라 하자. $B$가 가환
$A$-대수이면, 모든 $A$-선형 사상 $E\to B$는 유일하게
$E\xrightarrow{\sigma}\mathbf{S}(E)\xrightarrow{g}B$로 인수분해된다. 여기서
$g$는 $A$-대수 준동형이다. 이 특징지음에서 두 $A$-가군 $E$, $F$에
대하여
\[
  \mathbf{S}(E\oplus F)=\mathbf{S}(E)\otimes\mathbf{S}(F)
\]

\oldpage[II]{15}
가 표준 동형을 제외하면 성립함이 곧바로 따른다. 또한
$E\mapsto\mathbf{S}(E)$는 $A$-가군의 범주에서 가환 $A$-대수의 범주로
가는 공변 함자이다. 마지막으로 앞의 특징지음에 따라
$E=\varinjlim E_\lambda$이면
$\mathbf{S}(E)=\varinjlim\mathbf{S}(E_\lambda)$가 표준 동형을 제외하면
성립한다. 말의 남용으로 $x_i\in E$일 때의 곱
$\sigma(x_1)\sigma(x_2)\cdots\sigma(x_n)$을 혼동의 우려가 없으면 흔히
$x_1x_2\cdots x_n$으로 나타낸다. 대수 $\mathbf{S}(E)$는
\emph{등급} 대수이며, $\mathbf{S}_n(E)$는 $E$의 $n$개 원소의 곱들의
$A$-선형 결합 전체이다($n\geq 0$). 대수 $\mathbf{S}(A)$는 한 부정원의
다항식 대수 $A[T]$와 표준적으로 동형이고, 대수 $\mathbf{S}(A^n)$은
$n$개 부정원의 다항식 대수 $A[T_1,\ldots,T_n]$과 표준적으로 동형이다.
\end{env}

\begin{env}[1.7.2]
\label{II.1.7.2-ko}
$\varphi:A\to B$를 환 준동형이라 하자. $F$가 $B$-가군이면, 표준 사상
$F\to\mathbf{S}(F)$는 표준 사상
$F_{[\varphi]}\to\mathbf{S}(F)_{[\varphi]}$를 주고, 따라서 이 사상은
$F_{[\varphi]}\to\mathbf{S}(F_{[\varphi]})\to
\mathbf{S}(F)_{[\varphi]}$로 인수분해된다. 표준 $A$-대수 준동형
$\mathbf{S}(F_{[\varphi]})\to\mathbf{S}(F)_{[\varphi]}$는 전사이지만
반드시 전단사인 것은 아니다. $E$가 $A$-가군이면, 모든 디-준동형
$E\to F$, 즉 모든 $A$-가군 준동형 $E\to F_{[\varphi]}$는 따라서
표준적으로 $A$-대수 준동형
$\mathbf{S}(E)\to\mathbf{S}(F_{[\varphi]})\to
\mathbf{S}(F)_{[\varphi]}$를 준다.
이는 곧 대수의 디-준동형 $\mathbf{S}(E)\to\mathbf{S}(F)$이다.

같은 표기 아래에서 모든 $A$-가군 $E$에 대하여
$\mathbf{S}(E\otimes_A B)$는 대수
$\mathbf{S}(E)\otimes_A B$와 표준적으로 동일시된다. 이는
$\mathbf{S}(E)$의 보편 성질에서 곧바로 따른다
(\hyperref[II.1.7.1-ko]{1.7.1}).
\end{env}

\begin{env}[1.7.3]
\label{II.1.7.3-ko}
$R$을 환 $A$의 곱셈적 부분집합이라 하자.
(\hyperref[II.1.7.2-ko]{1.7.2})를 환 $B=R^{-1}A$에 적용하고
$R^{-1}E=E\otimes_A R^{-1}A$임을 상기하면,
$\mathbf{S}(R^{-1}E)=R^{-1}\mathbf{S}(E)$가 표준 동형을 제외하면
성립함을 알 수 있다. 또한 $R'\supset R$이 $A$의 또 하나의 곱셈적
부분집합이면 도식
\[
  \xymatrix{
    R^{-1}E\ar[r]\ar[d] & {R'}^{-1}E\ar[d]\\
    \mathbf{S}(R^{-1}E)\ar[r] & \mathbf{S}({R'}^{-1}E)
  }
\]
은 가환한다.
\end{env}

\begin{env}[1.7.4]
\label{II.1.7.4-ko}
이제 $(S,\mathcal{A})$를 환 달린 공간, $\mathcal{E}$를 $S$ 위의
$\mathcal{A}$-가군이라 하자. 모든 열린집합 $U\subset S$에
$\Gamma(U,\mathcal{A})$-가군
$\mathbf{S}(\Gamma(U,\mathcal{E}))$을 대응시키면,
$\mathbf{S}(E)$의 함자성에 따라 대수의 준층이 정의된다
(\hyperref[II.1.7.2-ko]{1.7.2}). 이 준층의 층화를
$\mathbf{S}(\mathcal{E})$ 또는
$\mathbf{S}_{\mathcal{A}}(\mathcal{E})$로 나타내고,
$\mathcal{A}$-가군 $\mathcal{E}$의 \emph{대칭 대수}라 한다.
(\hyperref[II.1.7.1-ko]{1.7.1})에서
$\mathbf{S}(\mathcal{E})$가 보편 문제의 해임이 곧바로 따른다.
$\mathcal{B}$가 $\mathcal{A}$-대수일 때 모든 $\mathcal{A}$-가군
준동형 $\mathcal{E}\to\mathcal{B}$는 유일하게
$\mathcal{E}\to\mathbf{S}(\mathcal{E})\to\mathcal{B}$로 인수분해되며,
둘째 화살표는 $\mathcal{A}$-대수 준동형이다. 따라서
$\mathcal{A}$-가군 준동형 $\mathcal{E}\to\mathcal{B}$들과
$\mathcal{A}$-대수 준동형
$\mathbf{S}(\mathcal{E})\to\mathcal{B}$들 사이에는 전단사 대응이
있다. 특히 모든 $\mathcal{A}$-가군 준동형
$u:\mathcal{E}\to\mathcal{F}$는 $\mathcal{A}$-대수 준동형
$\mathbf{S}(u):\mathbf{S}(\mathcal{E})\to
\mathbf{S}(\mathcal{F})$를 정의하며, 따라서
$\mathcal{E}\mapsto\mathbf{S}(\mathcal{E})$는 공변 함자이다.

\oldpage[II]{16}
(\hyperref[II.1.7.2-ko]{1.7.2})와 $\mathbf{S}$가 귀납극한과 가환한다는
사실에 따라, 모든 점 $s\in S$에 대하여
$(\mathbf{S}(\mathcal{E}))_s=\mathbf{S}(\mathcal{E}_s)$이다.
$\mathcal{E}$, $\mathcal{F}$가 두 $\mathcal{A}$-가군이면
$\mathbf{S}(\mathcal{E}\oplus\mathcal{F})$는
$\mathbf{S}(\mathcal{E})\otimes_{\mathcal{A}}
\mathbf{S}(\mathcal{F})$와 표준적으로 동일시된다. 이는 대응하는
준층들에서 확인된다.

또한 $\mathbf{S}(\mathcal{E})$는 $\mathbf{S}_n(\mathcal{E})$들의 무한
직합인 등급 $\mathcal{A}$-대수이다. 여기서 $\mathcal{A}$-가군
$\mathbf{S}_n(\mathcal{E})$는 준층
$U\mapsto\mathbf{S}_n(\Gamma(U,\mathcal{E}))$의 층화이다. 특히
$\mathcal{E}=\mathcal{A}$로 두면
$\mathbf{S}_{\mathcal{A}}(\mathcal{A})$는
$\mathcal{A}[T]=\mathcal{A}\otimes_{\mathbf{Z}}\mathbf{Z}[T]$와
동일시된다($T$는 부정원이고, $\mathbf{Z}$는 상수층으로 본다).
\end{env}

\begin{env}[1.7.5]
\label{II.1.7.5-ko}
$(T,\mathcal{B})$를 또 하나의 환 달린 공간,
$f:(S,\mathcal{A})\to(T,\mathcal{B})$를 환 달린 공간의 사상이라 하자.
$\mathcal{F}$가 $\mathcal{B}$-가군이면
$\mathbf{S}(f^*(\mathcal{F}))$는
$f^*(\mathbf{S}(\mathcal{F}))$와 표준적으로 동일시된다. 실제로
$f=(\psi,\theta)$이면 정의상
(\hyperref[0.4.3.1-ko]{0, 4.3.1})
\[
  \mathbf{S}(f^*(\mathcal{F}))
  =\mathbf{S}(\psi^*(\mathcal{F})
    \otimes_{\psi^*(\mathcal{B})}\mathcal{A})
  =\mathbf{S}(\psi^*(\mathcal{F}))
    \otimes_{\psi^*(\mathcal{B})}\mathcal{A}
\]
이다(\hyperref[II.1.7.2-ko]{1.7.2}). $S$의 임의의 열린집합 $U$와
$U$ 위의 $\mathbf{S}(\psi^*(\mathcal{F}))$의 임의의 단면 $h$에
대하여, 각 점 $s\in U$의 어떤 근방 $V$에서 $h$는
$\mathbf{S}(\Gamma(V,\psi^*(\mathcal{F})))$의 한 원소와 일치한다.
$\psi^*(\mathcal{F})$의 정의
(\hyperref[0.3.7.1-ko]{0, 3.7.1})를 참조하고, 가군 $E$에 대한
$\mathbf{S}(E)$의 모든 원소가 $E$의 원소들의 곱을 유한 개만 사용한
선형 결합임을 고려하면, $T$에서 $\psi(s)$의 어떤 근방 $W$,
$W$ 위의 $\mathbf{S}(\mathcal{F})$의 어떤 단면 $h'$, 그리고 $s$의
어떤 근방 $V'\subset V\cap\psi^{-1}(W)$가 존재하여 $V'$ 위에서
$h$는 $t\mapsto h'(\psi(t))$와 일치함을 알 수 있다. 따라서 주장이
따른다.
\end{env}

\begin{proposition}[1.7.6]
\label{II.1.7.6-ko}
$A$를 환, $S=\operatorname{Spec}(A)$를 그 소 스펙트럼,
$\mathcal{E}=\widetilde{M}$을 $A$-가군 $M$에 대응하는
$\mathcal{O}_S$-가군이라 하자. 그러면 $\mathcal{O}_S$-대수
$\mathbf{S}(\mathcal{E})$는 $A$-대수 $\mathbf{S}(M)$에 대응하는
$\mathcal{O}_S$-대수이다.
\end{proposition}

\begin{proof}
실제로 모든 $f\in A$에 대하여
$\mathbf{S}(M_f)=(\mathbf{S}(M))_f$이다
(\hyperref[II.1.7.3-ko]{1.7.3}). 따라서 명제는 정의들
(\hyperref[I.1.3.4-ko]{I, 1.3.4})에서 따른다.
\end{proof}

\begin{corollary}[1.7.7]
\label{II.1.7.7-ko}
$S$가 준스킴이고 $\mathcal{E}$가 준연접 $\mathcal{O}_S$-가군이면,
$\mathcal{O}_S$-대수 $\mathbf{S}(\mathcal{E})$는 준연접이다. 또한
$\mathcal{E}$가 유한 생성 $\mathcal{O}_S$-가군이면, 각
$\mathcal{O}_S$-가군 $\mathbf{S}_n(\mathcal{E})$도 유한 생성이다.
\end{corollary}

\begin{proof}
첫째 주장은 (\hyperref[II.1.7.6-ko]{1.7.6})과
(\hyperref[I.1.4.1-ko]{I, 1.4.1})에서 곧바로 따른다. 둘째 주장은
$E$가 유한 생성 $A$-가군이면 $\mathbf{S}_n(E)$도 유한 생성
$A$-가군이라는 사실에서 따른다. 이제
(\hyperref[I.1.3.13-ko]{I, 1.3.13})을 적용하면 된다.
\end{proof}

\begin{definition}[1.7.8]
\label{II.1.7.8-ko}
$\mathcal{E}$를 준연접 $\mathcal{O}_S$-가군이라 하자. 준연접
$\mathcal{O}_S$-대수 $\mathbf{S}(\mathcal{E})$의 스펙트럼
(\hyperref[II.1.3.1-ko]{1.3.1})을 $\mathbf{V}(\mathcal{E})$로 나타내고,
이를 $\mathcal{E}$로 정의되는 $S$ 위의 \emph{벡터 다발}이라 한다.
\end{definition}

(\hyperref[II.1.2.7-ko]{1.2.7})에 따라 임의의 $S$-준스킴 $X$에 대하여,
$S$-사상 $X\to\mathbf{V}(\mathcal{E})$들과 $\mathcal{O}_S$-대수 준동형
$\mathbf{S}(\mathcal{E})\to\mathcal{A}(X)$들 사이에는 표준적인 전단사
대응이 있고, 따라서 이 $S$-사상들과 $\mathcal{O}_S$-가군 준동형
$\mathcal{E}\to\mathcal{A}(X)=f_*(\mathcal{O}_X)$들 사이에도 그러한
대응이 있다(여기서 $f$는 구조 사상 $X\to S$이다). 특히 다음을 얻는다.

\begin{env}[1.7.9]
\label{II.1.7.9-ko}
$X$로 $S$가 열린집합 $U\subset S$ 위에 유도하는 부분준스킴을 잡자.
그러면 $S$-사상 $U\to\mathbf{V}(\mathcal{E})$들은 열린집합
$p^{-1}(U)$ 위에 $\mathbf{V}(\mathcal{E})$가 유도하는 $U$-준스킴의
$U$-단면들(\hyperref[I.2.5.5-ko]{I, 2.5.5})에 지나지 않는다(여기서
$p$는 구조 사상 $\mathbf{V}(\mathcal{E})\to S$이다). 방금 본 바에
따르면 이 $U$-단면들은 $\mathcal{O}_S$-가군 준동형
$\mathcal{E}\to j_*(\mathcal{O}_S|U)$들과 전단사로 대응한다(여기서
$j$는 표준 포함 사상 $U\to S$이다). 또는, 이와 동치로

\oldpage[II]{17}
(\hyperref[0.4.4.3-ko]{0, 4.4.3}), 이들은 $(\mathcal{O}_S|U)$-가군
준동형
$j^*(\mathcal{E})=\mathcal{E}|U\to\mathcal{O}_S|U$들과 대응한다. 또한
$S$-사상 $U\to\mathbf{V}(\mathcal{E})$을 열린집합 $U'\subset U$로
제한한 것에는 대응하는 준동형
$\mathcal{E}|U\to\mathcal{O}_S|U$를 $U'$로 제한한 것이 대응함은
자명하다. 따라서 $\mathbf{V}(\mathcal{E})$의 \emph{$S$-단면의 싹으로
이루어진 층}은 $\mathcal{E}$의 \emph{쌍대 $\check{\mathcal{E}}$}와
표준적으로 동일시된다.

특히 $X=U=S$로 두면 영 준동형
$\mathcal{E}\to\mathcal{O}_S$에는 $\mathbf{V}(\mathcal{E})$의 표준
$S$-단면 하나가 대응하며, 이를 \emph{영 $S$-단면}이라 한다
(8.3.3 참조).
\end{env}

\begin{env}[1.7.10]
\label{II.1.7.10-ko}
이제 $X$로 체 $K$의 스펙트럼 $\{\xi\}$를 잡자. 그러면 구조 사상
$f:X\to S$는 $s=f(\xi)$일 때 체의 단사 준동형
$k(s)\to K$에 대응한다(\hyperref[I.2.4.6-ko]{I, 2.4.6}). 따라서
$S$-사상 $\{\xi\}\to\mathbf{V}(\mathcal{E})$들은 $k(s)$의 확대체
$K$를 값체로 하는 $\mathbf{V}(\mathcal{E})$의 \emph{기하적 점들}
(\hyperref[I.3.4.5-ko]{I, 3.4.5})에 지나지 않으며, 이 기하적 점들은
$p^{-1}(s)$의 점들에 놓인다. 이 점들의 집합은 점 $s$ 위의
$\mathbf{V}(\mathcal{E})$의 \emph{$K$-유리 기하적 섬유}라고도 할 수
있으며, (\hyperref[II.1.7.8-ko]{1.7.8})에 따라 $\mathcal{O}_S$-가군
준동형 $\mathcal{E}\to f_*(\mathcal{O}_X)$들의 집합과 동일시된다.
또는, 이와 동치로 (\hyperref[0.4.4.3-ko]{0, 4.4.3}),
$\mathcal{O}_X$-가군 준동형
$f^*(\mathcal{E})\to\mathcal{O}_X=K$들의 집합과 동일시된다. 한편 정의상
(\hyperref[0.4.3.1-ko]{0, 4.3.1})
$f^*(\mathcal{E})=\mathcal{E}_s\otimes_{\mathcal{O}_s}K
=\mathcal{E}^s\otimes_{k(s)}K$이며, 여기서
$\mathcal{E}^s=\mathcal{E}_s/\mathfrak{m}_s\mathcal{E}_s$로 둔다.
따라서 점 $s$ 위의 $\mathbf{V}(\mathcal{E})$의 $K$-유리 기하적
섬유는 \emph{$K$-벡터 공간
$\mathcal{E}^s\otimes_{k(s)}K$}의 \emph{쌍대}와 동일시된다.
$\mathcal{E}^s$ 또는 $K$가 $k(s)$ 위에서 유한 차원이면 이 쌍대는
$(\mathcal{E}^s)^\vee\otimes_{k(s)}K$와도 동일시된다. 여기서
$(\mathcal{E}^s)^\vee$는 $k(s)$-벡터 공간 $\mathcal{E}^s$의 쌍대를
나타낸다.
\end{env}

\begin{proposition}[1.7.11]
\label{II.1.7.11-ko}
\begin{enumerate}
  \item[\rm(i)] $\mathbf{V}(\mathcal{E})$는 준연접
    $\mathcal{O}_S$-가군의 범주에서 아핀 $S$-스킴의 범주로 가는,
    $\mathcal{E}$에 관한 반변 함자이다.
  \item[\rm(ii)] $\mathcal{E}$가 유한 생성 $\mathcal{O}_S$-가군이면,
    $\mathbf{V}(\mathcal{E})$는 $S$ 위에서 유한형이다.
  \item[\rm(iii)] $\mathcal{E}$와 $\mathcal{F}$가 두 준연접
    $\mathcal{O}_S$-가군이면,
    $\mathbf{V}(\mathcal{E}\oplus\mathcal{F})$는
    $\mathbf{V}(\mathcal{E})\times_S\mathbf{V}(\mathcal{F})$와
    표준적으로 동일시된다.
  \item[\rm(iv)] $g:S'\to S$를 사상이라 하자. 임의의 준연접
    $\mathcal{O}_S$-가군 $\mathcal{E}$에 대하여,
    $\mathbf{V}(g^*(\mathcal{E}))$는
    $\mathbf{V}(\mathcal{E})_{(S')}
    =\mathbf{V}(\mathcal{E})\times_S S'$와 표준적으로 동일시된다.
  \item[\rm(v)] 준연접 $\mathcal{O}_S$-가군의 전사 준동형
    $\mathcal{E}\to\mathcal{F}$에는 닫힌 몰입 사상
    $\mathbf{V}(\mathcal{F})\to\mathbf{V}(\mathcal{E})$가 대응한다.
\end{enumerate}
\end{proposition}

\begin{proof}
임의의 $\mathcal{O}_S$-가군 준동형 $\mathcal{E}\to\mathcal{F}$가
$\mathcal{O}_S$-대수 준동형
$\mathbf{S}(\mathcal{E})\to\mathbf{S}(\mathcal{F})$를 표준적으로
정한다는 점을 고려하면, (i)는 (\hyperref[II.1.2.7-ko]{1.2.7})의
직접적인 귀결이다. (ii)는 정의들
(\hyperref[I.6.3.1-ko]{I, 6.3.1})과, $E$가 유한 생성 $A$-가군이면
$\mathbf{S}(E)$가 유한 생성 $A$-대수라는 사실에서 곧바로 따른다.
(iii)를 증명하려면 표준 동형
$\mathbf{S}(\mathcal{E}\oplus\mathcal{F})
\simeq\mathbf{S}(\mathcal{E})
\otimes_{\mathcal{O}_S}\mathbf{S}(\mathcal{F})$
(\hyperref[II.1.7.4-ko]{1.7.4})에서 출발하여
(\hyperref[II.1.4.6-ko]{1.4.6})을 적용하면 충분하다. 마찬가지로
(iv)를 증명하려면 표준 동형
$\mathbf{S}(g^*(\mathcal{E}))
\simeq g^*(\mathbf{S}(\mathcal{E}))$
(\hyperref[II.1.7.5-ko]{1.7.5})에서 출발하여
(\hyperref[II.1.5.2-ko]{1.5.2})를 적용하면 충분하다. 끝으로 (v)를
보이려면 준동형 $\mathcal{E}\to\mathcal{F}$가 전사일 때 이에 대응하는
$\mathcal{O}_S$-대수 준동형
$\mathbf{S}(\mathcal{E})\to\mathbf{S}(\mathcal{F})$도 전사임을
관찰하면 충분하며, 결론은 (\hyperref[II.1.4.10-ko]{1.4.10})에서
따른다.
\end{proof}

\begin{env}[1.7.12]
\label{II.1.7.12-ko}
특히 $\mathcal{E}=\mathcal{O}_S$로 잡자. 준스킴
$\mathbf{V}(\mathcal{O}_S)$는 $\mathcal{O}_S$-대수
$\mathbf{S}(\mathcal{O}_S)$의 스펙트럼인 아핀 $S$-스킴이며, 이
대수는 $\mathcal{O}_S$-대수
$\mathcal{O}_S[T]=\mathcal{O}_S\otimes_{\mathbf{Z}}\mathbf{Z}[T]$와
\oldpage[II]{18}
동일시된다($T$는 부정원이다). $S=\operatorname{Spec}(\mathbf{Z})$일
때에는 (\hyperref[II.1.7.6-ko]{1.7.6})에 따라 자명하며, 구조 사상
$S\to\operatorname{Spec}(\mathbf{Z})$를 고려하고
(\hyperref[II.1.7.11-ko]{1.7.11, (iv)})를 사용하면 여기서 일반적인
경우로 넘어갈 수 있다. 이 결과에 따라
$\mathbf{V}(\mathcal{O}_S)=S[T]$로도 쓰며, 따라서 다음 공식을 얻는다.
\[
  S[T]=S\otimes_{\mathbf{Z}}\mathbf{Z}[T].
  \tag{1.7.12.1}\label{II.1.7.12.1-ko}
\]

(\hyperref[I.3.3.15-ko]{I, 3.3.15})에서 이미 본, $S[T]$의
$S$-단면의 싹으로 이루어진 층과 $\mathcal{O}_S$의 동일시는 여기에서
(\hyperref[II.1.7.9-ko]{1.7.9})의 특수한 경우로서 더 일반적인 맥락
안에서 다시 얻어진다.
\end{env}

\begin{env}[1.7.13]
\label{II.1.7.13-ko}
임의의 $S$-준스킴 $X$에 대하여,
(\hyperref[II.1.7.8-ko]{1.7.8})에서
$\operatorname{Hom}_S(X,S[T])$가
$\operatorname{Hom}_{\mathcal{O}_S}(\mathcal{O}_S,\mathcal{A}(X))$와
표준적으로 동일시됨을 보았다. 이 집합은
$\Gamma(S,\mathcal{A}(X))$와 표준적으로 동형이므로 환 구조가 주어진다.
또한 임의의 $S$-사상 $h:X\to Y$에는 이 환 구조들에 관한 사상
$\Gamma(\mathcal{A}(h)):\Gamma(S,\mathcal{A}(Y))
\to\Gamma(S,\mathcal{A}(X))$가 대응한다
(\hyperref[II.1.1.2-ko]{1.1.2}).
$\operatorname{Hom}_S(X,S[T])$에 이와 같이 정의된 환 구조를 주면,
따라서 $\operatorname{Hom}_S(X,S[T])$를 $S$-준스킴의 범주에서 환의
범주로 가는 \emph{$X$에 관한 반변 함자}로 볼 수 있다.
한편 $\operatorname{Hom}_S(X,\mathbf{V}(\mathcal{E}))$도 마찬가지로
$\operatorname{Hom}_{\mathcal{O}_S}(\mathcal{E},\mathcal{A}(X))$와
동일시된다(여기서는 $\mathcal{A}(X)$를
\emph{$\mathcal{O}_S$-가군}으로 본다). 따라서 여기에 환
$\operatorname{Hom}_S(X,S[T])$ 위의 \emph{가군} 구조를 표준적으로
줄 수 있으며, 위와 마찬가지로 쌍
\[
  \bigl(\operatorname{Hom}_S(X,S[T]),
  \operatorname{Hom}_S(X,\mathbf{V}(\mathcal{E}))\bigr)
\]
은 $X$에 관한 반변 함자임을 알 수 있다. 이 함자는 환 $A$와
$A$-가군 $M$으로 이루어진 쌍 $(A,M)$을 대상으로 하고 디-준동형들을
사상으로 하는 범주에 값을 갖는다.

이 사실들을 $S[T]$가 \emph{$S$ 위의 환 스킴}이고
$\mathbf{V}(\mathcal{E})$가 $S$ 위의 환 스킴 $S[T]$ 위의
\emph{$S$ 위의 가군 스킴}이라는 말로 해석하겠다
(0장, \textsection8 참조).
\end{env}

\begin{env}[1.7.14]
\label{II.1.7.14-ko}
$S$-스킴 $\mathbf{V}(\mathcal{E})$ 위에 정의된 $S$ 위의 가군 스킴
구조를 사용하여 $\mathcal{O}_S$-가군 $\mathcal{E}$를 유일한 동형을
제외하고 복원할 수 있음을 보이겠다. 이를 위해 $\mathcal{E}$가 이
구조를 사용하여 정의되는
$\mathbf{S}(\mathcal{E})=\mathcal{A}(\mathbf{V}(\mathcal{E}))$의 한
부분-$\mathcal{O}_S$-가군과 표준적으로 동형임을 보이겠다. 실제로
(\hyperref[II.1.7.4-ko]{1.7.4})에 따라
\emph{$\mathcal{O}_S$-대수} 준동형들의 집합
$\operatorname{Hom}_{\mathcal{O}_S}(\mathbf{S}(\mathcal{E}),
\mathcal{A}(X))$는 \emph{$\mathcal{O}_S$-가군} 준동형들의 집합
$\operatorname{Hom}_{\mathcal{O}_S}(\mathcal{E},\mathcal{A}(X))$와
표준적으로 동일시된다. $h$, $h'$을 이 뒤의 집합의 두 원소,
$s_i$ ($1\leq i\leq n$)를 열린집합 $U\subset S$ 위의
$\mathcal{E}$의 단면들, $t$를 $U$ 위의 $\mathcal{A}(X)$의 단면이라
하면, 정의에 따라
\[
  (h+h')(s_1s_2\cdots s_n)
  =\prod_{i=1}^n\bigl(h(s_i)+h'(s_i)\bigr)
\]
이고
\[
  (t.h)(s_1s_2\cdots s_n)=t^n\prod_{i=1}^n h(s_i).
\]

이제 $z$가 $U$ 위의 $\mathbf{S}(\mathcal{E})$의 단면이면,
$h\mapsto h(z)$는
$\operatorname{Hom}_S(X,\mathbf{V}(\mathcal{E}))
=\operatorname{Hom}_{\mathcal{O}_S}(\mathbf{S}(\mathcal{E}),
\mathcal{A}(X))$에서 $\Gamma(U,\mathcal{A}(X))$로 가는 사상이다.
이제 다음을
\oldpage[II]{19}
보이겠다. $\mathcal{E}$는 다음 성질을 갖는
$\mathbf{S}(\mathcal{E})$의 부분가군과 동일시된다.
\emph{모든 열린집합 $U\subset S$, 이 부분-$\mathcal{O}_S$-가군의
$U$ 위의 모든 단면 $z$, 모든 $S$-준스킴 $X$에 대하여,
$\operatorname{Hom}_{\mathcal{O}_S|U}(\mathbf{S}(\mathcal{E})|U,
\mathcal{A}(X)|U)$에서 $\Gamma(U,\mathcal{A}(X))$로 가는 사상
$h\mapsto h(z)$가 $\Gamma(U,\mathcal{A}(X))$-가군 준동형이다.}

실제로 $\mathcal{E}$가 이 성질을 갖는다는 것은 곧바로 알 수 있다.
그 역을 증명하려면, $S=\operatorname{Spec}(A)$,
$\mathcal{E}=\widetilde{M}$일 때 $S$ 위의
$\mathbf{S}(\mathcal{E})$의 단면 $z$가 ($U=S$에 대하여) 위의
성질을 가지면 반드시 $\mathcal{E}$의 단면임을 증명하는 것으로
충분하다. 이때 $z=\sum_{n=0}^{\infty}z_n$이고
$z_n\in\mathbf{S}_n(M)$이며, $n\ne1$이면 $z_n=0$임을 보이면 된다.
$B=\mathbf{S}(M)$으로 놓고 $X$로 준스킴
$\operatorname{Spec}(B[T])$를 택하자. 여기서 $T$는 부정원이다.
집합
$\operatorname{Hom}_{\mathcal{O}_S}(\mathbf{S}(\mathcal{E}),
\mathcal{A}(X))$는 환 준동형 $h:B\to B[T]$들의 집합과 동일시된다
(\hyperref[I.1.3.13-ko]{I, 1.3.13}). 위에서 본 바에 따라
$(T.h)(z)=\sum_{n=0}^{\infty}T^nh(z_n)$이며, $z$에 관한 가정은 모든
준동형 $h$에 대하여
$\sum_{n=0}^{\infty}T^nh(z_n)=T.\sum_{n=0}^{\infty}h(z_n)$임을
뜻한다. 특히 $h$로 표준 포함 사상을 택하면
$\sum_{n=0}^{\infty}T^nz_n=T.\sum_{n=0}^{\infty}z_n$을 얻으며,
여기서 실제로 $n\ne1$이면 $z_n=0$이라는 결론이 따른다.
\end{env}

\begin{proposition}[1.7.15]
\label{II.1.7.15-ko}
$Y$를 바탕공간이 뇌터인 준스킴 또는 준콤팩트 스킴이라 하자.
$Y$ 위에서 유한형인 임의의 아핀 $Y$-스킴 $X$는
$\mathbf{V}(\mathcal{E})$ 꼴의 $Y$-스킴의 닫힌 부분-$Y$-스킴과
$Y$-동형이다. 여기서 $\mathcal{E}$는 유한 생성 준연접
$\mathcal{O}_Y$-가군이다.
\end{proposition}

\begin{proof}
실제로 준연접 $\mathcal{O}_Y$-대수 $\mathcal{A}(X)$는 유한 생성이다
(\hyperref[II.1.3.7-ko]{1.3.7}). 가정에 따라 $\mathcal{A}(X)$는
유한 생성 준연접 부분-$\mathcal{O}_Y$-가군 $\mathcal{E}$에 의해
대수로서 생성된다(\hyperref[I.9.6.5-ko]{I, 9.6.5}). 정의에 따라
이는 포함 사상 $\mathcal{E}\to\mathcal{A}(X)$를 표준적으로 연장하는
표준 준동형 $\mathbf{S}(\mathcal{E})\to\mathcal{A}(X)$가
\emph{전사}임을 뜻한다. 그러면 결론은
(\hyperref[II.1.4.10-ko]{1.4.10})에서 따른다.
\end{proof}

\section{동차 소 스펙트럼}
\label{section:II.2-ko}

\subsection{등급환과 등급가군에 관한 일반 사항.}
\label{subsection:II.2.1-ko}

\begin{notation}[2.1.1]
\label{II.2.1.1-ko}
음이 아닌 차수로 \emph{등급}이 주어진 환 $S$에 대하여,
차수가 $n$인 동차 원소들로 이루어진 $S$의 부분을 $S_n$으로
($n\geq0$), $n>0$인 $S_n$들의 합(직합)을 $S_+$로 나타내겠다.
$1\in S_0$이고, $S_0$는 $S$의 부분환이며, $S_+$는 $S$의 동차
아이디얼이고, $S$는 $S_0$와 $S_+$의 직합이다. $M$이 $S$ 위의
\emph{등급가군}이면(음의 차수도 허용한다), 마찬가지로
$M$의 차수가 $n$인 동차 원소들로 이루어진 $S_0$-가군을
$M_n$으로 나타내겠다(여기서 $n\in\mathbf{Z}$).

임의의 정수 $d>0$에 대하여, $S_{nd}$들의 직합을 $S^{(d)}$로
나타낸다. $S_{nd}$의 원소들을 차수가 $n$인 동차 원소로 보면,
$S_{nd}$들은 $S^{(d)}$ 위에 등급환 구조를 정의한다.

$0\leq k\leq d-1$인 임의의 정수 $k$에 대하여, $M^{(d,k)}$로
다음 직합을 나타내겠다.
\oldpage[II]{20}
$M_{nd+k}$ ($n\in\mathbf{Z}$)들의 직합이다.
$M_{nd+k}$의 원소들을 차수가 $n$인 동차 원소로 보면, 이는 등급
$S^{(d)}$-가군이다. $M^{(d,0)}$ 대신 $M^{(d)}$로 쓰겠다.

앞의 표기법 아래에서, 임의의 정수 $n$(음수도 허용한다)에 대하여,
모든 $k\in\mathbf{Z}$에 대해 $(M(n))_k=M_{n+k}$로 정의되는 등급
$S$-가군을 $M(n)$으로 나타내겠다. 특히 $S(n)$은
$(S(n))_k=S_{n+k}$를 만족하는 등급 $S$-가군이며,
$n<0$이면 $S_n=0$으로 놓기로 한다. 등급 $S$-가군 $M$이
\emph{등급가군으로서} $S(n)$ 꼴의 가군들의 직합과 동형이면,
\emph{자유}라고 한다. $S(n)$은 $S$의 원소 $1$을 차수가
$-n$인 원소로 보아 이를 생성원으로 갖는
$S$-가군이므로, 이는 $M$이 \emph{동차} 원소들로 이루어진
$S$ 위의 \emph{기저}를 갖는다고 말하는 것과 같다.

등급 $S$-가군 $M$이 \emph{유한 표시를 갖는다}는 것은, $P$와 $Q$가
$S(n)$ 꼴의 가군들의 유한 직합이고 준동형들의 차수가 $0$인 완전열
$P\to Q\to M\to0$이 존재한다는 뜻이다
(\agref{II.2.1.2-ko}{2.1.2} 참조).
\end{notation}

\begin{env}[2.1.2]
\phantomsection
\label{II.2.1.2-ko}
$M$, $N$을 두 등급 $S$-가군이라 하자. $M\otimes_S N$ 위에 다음과
같이 \emph{등급} $S$-가군 구조를 정의한다. 텐서곱
$M\otimes_{\mathbf{Z}}N$ 위에는
$\mathbf{Z}$-가군의 등급 구조($\mathbf{Z}$에는
$\mathbf{Z}_0=\mathbf{Z}$, $n\ne0$일 때 $\mathbf{Z}_n=0$인 등급을
준다)를 다음과 같이 정의할 수 있다.
\[
  (M\otimes_{\mathbf{Z}}N)_q
  =\bigoplus_{m+n=q}M_m\otimes_{\mathbf{Z}}N_n
\]
($M$과 $N$은 각각 $M_m$들과 $N_n$들의 직합이므로,
$M\otimes_{\mathbf{Z}}N$을 모든 $M_m\otimes_{\mathbf{Z}}N_n$의
직합과 표준적으로 동일시할 수 있음을 알고 있다.) 이제
$M\otimes_S N=(M\otimes_{\mathbf{Z}}N)/P$이다. 여기서 $P$는
$x\in M$, $y\in N$, $s\in S$에 대하여
$(xs)\otimes y-x\otimes(sy)$ 꼴의 원소들로 생성되는
$M\otimes_{\mathbf{Z}}N$의 부분-$\mathbf{Z}$-가군이다. $P$가
$M\otimes_{\mathbf{Z}}N$의 \emph{등급} 부분-$\mathbf{Z}$-가군임은
명백하며, 몫을 취하면 실제로 $M\otimes_S N$ 위에 등급 $S$-가군
구조가 얻어짐을 곧바로 알 수 있다.

두 등급 $S$-가군 $M$, $N$에 대하여, $S$-가군 준동형
$u:M\to N$이 모든 $j\in\mathbf{Z}$에 대하여
$u(M_j)\subset N_{j+k}$를 만족하면 \emph{차수 $k$}라고 함을
상기하자. $M$에서 $N$으로 가는 차수 $n$인 모든 준동형의 집합을
$H_n$이라 하고, $M$에서 $N$으로 가는 모든 ($S$-가군) 준동형들의
$S$-가군 $H$ 안에서 $H_n$들의 \emph{합}(직합)($n\in\mathbf{Z}$)을
$\operatorname{Hom}_S(M,N)$으로 나타낸다. 일반적으로
$\operatorname{Hom}_S(M,N)$은 앞서 말한 가군 전체와 같지 않다. 다만 $M$이
\emph{유한 생성}이면 $H=\operatorname{Hom}_S(M,N)$이다. 실제로 이때
$M$이 유한 개의 동차 원소 $x_i$ ($1\leq i\leq n$)로 생성된다고
할 수 있고, 모든 준동형 $u\in H$는 유일하게
$\sum_{k\in\mathbf{Z}}u_k$로 쓰인다. 여기서 각 $k$에 대하여
$u_k(x_i)$는 $u(x_i)$의 차수가 $k+\deg(x_i)$인 동차 성분과 같고
($1\leq i\leq n$), 따라서 유한 개의 지표를 제외하면 $u_k=0$이다.
정의에 따라 $u_k\in H_k$이므로 결론을 얻는다.

$\operatorname{Hom}_S(M,N)$의 차수 $0$인 원소들을 \emph{등급
$S$-가군의 준동형}이라 한다. $S_mH_n\subset H_{m+n}$임은 명백하므로,
$H_n$들은 $\operatorname{Hom}_S(M,N)$ 위에 등급 $S$-가군 구조를
정의한다.

이 정의들로부터 두 등급 $S$-가군 $M$, $N$에 대하여 다음을 곧바로
얻는다.
\[
  \label{II.2.1.2.1-ko}
  M(m)\otimes_S N(n)=(M\otimes_S N)(m+n),
  \tag{2.1.2.1}
\]
\[
  \label{II.2.1.2.2-ko}
  \operatorname{Hom}_S(M(m),N(n))
  =(\operatorname{Hom}_S(M,N))(n-m)
  \tag{2.1.2.2}
\]
\oldpage[II]{21}
$S$, $S'$을 두 등급환이라 하자. \emph{등급환}의 준동형
$\varphi:S\to S'$은 모든 $n\in\mathbf{Z}$에 대하여
$\varphi(S_n)\subset S'_n$을 만족하는 환 준동형이다(다시 말해,
$\varphi$는 등급 $\mathbf{Z}$-가군의 \emph{차수 $0$인} 준동형이어야
한다). 이러한 준동형 하나를 주면 $S'$ 위에 \emph{등급} $S$-가군 구조가
정의된다. 이 구조와 원래의 등급환 구조를 갖춘 $S'$을
\emph{등급 $S$-대수}라고 한다.

이제 $M$이 등급 $S$-가군이면, 등급 $S$-가군의 텐서곱
$M\otimes_S S'$ 위에 \emph{등급} $S'$-가군 구조가 자연스럽게 주어지며,
그 등급은 위에서 정의한 것이다.
\end{env}

\begin{lemma}[2.1.3]
\phantomsection
\label{II.2.1.3-ko}
$S$를 음이 아닌 차수로 등급이 주어진 환이라 하자. 동차 원소들로
이루어진 $S_+$의 부분집합 $E$가 $S_+$를 $S$-가군으로서 생성할
필요충분조건은 $E$가 $S$를 $S_0$-대수로서 생성하는 것이다.
\end{lemma}

\begin{proof}
조건이 충분함은 명백하다. 필요함을 보이자. $E$의 차수가 $n$인
원소들의 집합(각각 차수가 $\leq n$인 원소들의 집합)을 $E_n$
(각각 $E^n$)이라 하자. $n>0$에 대한 귀납법으로, $S_n$이
$E^n$의 원소들의 곱인 차수 $n$의 원소들로 생성되는 $S_0$-가군임을
보이면 충분하다. 그런데 $n=1$일 때 이는 가정에 의해 명백하다.
또한 가정으로부터
$S_n=\sum_{p=0}^{n-1}S_pE_{n-p}$이고, 이제 귀납 논법은 즉시
따른다.
\end{proof}

\begin{corollary}[2.1.4]
\phantomsection
\label{II.2.1.4-ko}
$S_+$가 유한 생성 아이디얼이기 위한 필요충분조건은 $S$가
$S_0$ 위의 유한 생성 대수인 것이다.
\end{corollary}

\begin{proof}
실제로 $S_0$-대수 $S$ (각각 $S$-아이디얼 $S_+$)의 유한 생성계가
동차 원소들로 이루어져 있다고 언제나 가정할 수 있다. 주어진 각 생성원을
그 동차 성분들로 바꾸면 되기 때문이다.
\end{proof}

\begin{corollary}[2.1.5]
\phantomsection
\label{II.2.1.5-ko}
$S$가 뇌터이기 위한 필요충분조건은 $S_0$가 뇌터이고 $S$가
$S_0$ 위의 유한 생성 대수인 것이다.
\end{corollary}

\begin{proof}
조건이 충분함은 명백하다. 필요함은 $S_0$가 $S/S_+$와 동형이고
$S_+$가 유한 생성 아이디얼이어야 한다는 데서 따른다
(\hyperref[II.2.1.4-ko]{2.1.4}).
\end{proof}

\begin{lemma}[2.1.6]
\phantomsection
\label{II.2.1.6-ko}
$S$를 음이 아닌 차수로 등급이 주어진 환이며 $S_0$ 위의 유한 생성
대수라고 하자. $M$을 유한 생성 등급 $S$-가군이라 하자. 그러면 다음이
성립한다.
\begin{enumerate}
  \item[\rm(i)] $M_n$들은 유한 생성 $S_0$-가군이고, 모든
    $n\leq n_0$에 대하여 $M_n=0$이 되는 정수 $n_0$가 존재한다.
  \item[\rm(ii)] 모든 정수 $n\geq n_1$에 대하여
    $M_{n+h}=S_hM_n$이 되는 정수 $n_1$과 정수 $h>0$이 존재한다.
  \item[\rm(iii)] $d>0$, $0\leq k\leq d-1$을 만족하는 모든 정수쌍
    $(d,k)$에 대하여 $M^{(d,k)}$는 유한 생성 $S^{(d)}$-가군이다.
  \item[\rm(iv)] 모든 정수 $d>0$에 대하여 $S^{(d)}$는 $S_0$ 위의
    유한 생성 대수이다.
  \item[\rm(v)] 모든 $m>0$에 대하여 $S_{mh}=(S_h)^m$이 되는 정수
    $h>0$이 존재한다.
  \item[\rm(vi)] 모든 정수 $n>0$에 대하여, 모든 $m\geq m_0$에 대해
    $S_m\subset S_+^n$이 되는 정수 $m_0$가 존재한다.
\end{enumerate}
\end{lemma}

\begin{proof}
$S$가 $S_0$-대수로서 차수가 $h_i$인 동차 원소 $f_i$
($1\leq i\leq r$)들로 생성되고, $M$이 $S$-가군으로서 차수가 $k_j$인
동차 원소 $x_j$ ($1\leq j\leq s$)들로 생성된다고 가정할 수 있다.
$M_n$은 $\alpha_i$들이 $\geq0$인 정수이고
$k_j+\sum_i\alpha_i h_i=n$을 만족할 때의 원소
\oldpage[II]{22}
$f_1^{\alpha_1}\cdots f_r^{\alpha_r}x_j$들의, 계수가 $S_0$에 속하는
선형결합들로
이루어짐이 명백하다. 각 $j$에 대하여 이 방정식을 만족하는
$(\alpha_i)$는 유한 개뿐이다. 이는 $h_i$들이 $>0$이기 때문이며, 따라서
(i)의 첫 번째 주장이 따른다. 두 번째 주장은 명백하다. 한편 $h$를
$h_i$들의 최소공배수라 하고, 모든 $g_i$의 차수가 $h$가 되도록
$g_i=f_i^{h/h_i}$ ($1\leq i\leq r$)로 놓자. 또 $z_\mu$들을
$0\leq\alpha_i<h/h_i$ ($1\leq i\leq r$)인
$f_1^{\alpha_1}\cdots f_r^{\alpha_r}x_j$ 꼴의 $M$의 원소들이라 하자.
이 원소들은 유한 개이고, 그 차수들 가운데 가장 큰 것을 $n_1$이라 하자.
$n\geq n_1$이면 $M_{n+h}$의 모든 원소가 $z_\mu$들의 선형결합이며 그
계수들은 $g_i$들에 관한 차수가 $>0$인 단항식임이 명백하다. 따라서
$M_{n+h}=S_hM_n$이고, 이로써 (ii)가 증명된다. 같은 방식으로 (모든
$d>0$에 대하여) $M^{(d,k)}$의 원소는 $g$가 $S$의 동차 원소이고
$0\leq\alpha_i<d$일 때의
$g^df_1^{\alpha_1}\cdots f_r^{\alpha_r}x_j$ 꼴의 원소들의, 계수가
$S_0$에 속하는 선형결합임을 알 수 있다. 따라서 (iii)가 성립한다. 이제
$M=S_+$로 놓고 $(S_+)^{(d)}=(S^{(d)})_+$임을 이용하면, (iv)는
(iii)와 보조정리 (\hyperref[II.2.1.3-ko]{2.1.3})에서 따른다. (v)는
(ii)에서 $M=S$로 놓으면 얻어진다. 마지막으로 주어진 $n$에 대하여
$\alpha_i\geq0$이고 $\sum_i\alpha_i<n$인 계 $(\alpha_i)$는 유한 개뿐이다.
따라서 이 계들에 대한 합 $\sum_i\alpha_i h_i$의 최댓값을 $m_0$라 하면
$m>m_0$에 대하여 $S_m\subset S_+^n$이고, 이로써 (vi)가 증명된다.
\footnote{\textbf{번역 주.} 인쇄 원문의 (vi)는 $m\geq m_0$라고 서술하지만,
여기서 $m_0$로 정의한 바로 그 최댓값을 쓰면 증명이 직접 주는 범위는
$m>m_0$이다. 존재명제 자체는 이 최댓값에 $1$을 더한 수를 임계값으로
취하면 원문의 $m\geq m_0$ 형태로 성립한다.}
\end{proof}

\begin{corollary}[2.1.7]
\phantomsection
\label{II.2.1.7-ko}
$S$가 뇌터이면 모든 정수 $d>0$에 대하여 $S^{(d)}$도 뇌터이다.
\end{corollary}

\begin{proof}
이는 (\hyperref[II.2.1.5-ko]{2.1.5})와
(\hyperref[II.2.1.6-ko]{2.1.6, (iv)})에서 따른다.
\end{proof}

\begin{env}[2.1.8]
\phantomsection
\label{II.2.1.8-ko}
$\mathfrak{p}$를 등급환 $S$의 동차 소아이디얼이라 하자. 그러면
$\mathfrak{p}$는 부분군 $\mathfrak{p}_n=\mathfrak{p}\cap S_n$들의
직합이다. $\mathfrak{p}$가 $S_+$를 포함하지 않는다고 가정하자. 그러면
$\mathfrak{p}$에 속하지 않는 $f\in S_+$에 대하여, 관계
$f^nx\in\mathfrak{p}$는 $x\in\mathfrak{p}$와 동치이다. 특히
$f\in S_d$ ($d>0$)이면, 모든 $x\in S_{m-nd}$에 대하여 관계
$f^nx\in\mathfrak{p}_m$은 $x\in\mathfrak{p}_{m-nd}$와 동치이다.
\end{env}

\begin{proposition}[2.1.9]
\phantomsection
\label{II.2.1.9-ko}
$n_0$를 $>0$인 정수라 하고, 각 $n\geq n_0$에 대하여
$\mathfrak{p}_n$을 $S_n$의 부분군이라 하자. $S_+$를 포함하지 않으며
모든 $n\geq n_0$에 대하여
$\mathfrak{p}\cap S_n=\mathfrak{p}_n$을 만족하는 $S$의 동차
소아이디얼 $\mathfrak{p}$가 존재하기 위한 필요충분조건은 다음 조건들이
성립하는 것이다.
\begin{enumerate}
  \item[$1^\circ$] 모든 $m\geq 0$과 모든 $n\geq n_0$에 대하여
    $S_m\mathfrak{p}_n\subset\mathfrak{p}_{m+n}$이다.
  \item[$2^\circ$] $m\geq n_0$, $n\geq n_0$, $f\in S_m$, $g\in S_n$일 때,
    관계 $fg\in\mathfrak{p}_{m+n}$은 $f\in\mathfrak{p}_m$ 또는
    $g\in\mathfrak{p}_n$을 함의한다.
  \item[$3^\circ$] 적어도 하나의 $n\geq n_0$에 대하여
    $\mathfrak{p}_n\neq S_n$이다.
\end{enumerate}
또한 이때 동차 소아이디얼 $\mathfrak{p}$는 유일하다.
\end{proposition}

\begin{proof}
조건 $1^\circ$와 $2^\circ$가 필요함은 명백하다. 또한
$\mathfrak{p}\not\supset S_+$이면
$\mathfrak{p}\cap S_k\neq S_k$인 $k>0$이 적어도 하나 존재한다.
$f\in S_k$가 $\mathfrak{p}$에 속하지 않으면, 관계
$\mathfrak{p}\cap S_n=S_n$은
(\hyperref[II.2.1.8-ko]{2.1.8})에 따라
$\mathfrak{p}\cap S_{n-mk}=S_{n-mk}$를 함의한다. 따라서 어떤 값 이상의
모든 $n$에 대하여 $\mathfrak{p}\cap S_n=S_n$이라면
$\mathfrak{p}\supset S_+$가 되어 가정에 모순이다. 이로써 $3^\circ$가
필요함이 증명된다. 역으로 조건 $1^\circ$, $2^\circ$, $3^\circ$가
성립한다고 가정하자. 어떤 정수 $d\geq n_0$에 대하여
$f\in S_d$가 $\mathfrak{p}_d$에 속하지 않으면, $\mathfrak{p}$가
존재할 때 $m<n_0$에 대한 $\mathfrak{p}_m$은 유한 개의 값을 제외한 모든
$r$에 대하여 $f^rx\in\mathfrak{p}_{m+rd}$를 만족하는
$x\in S_m$들의 집합과 반드시 같음에 유의하자. 이는 $\mathfrak{p}$가 존재한다면
유일함을 이미 증명한다. 이제 앞의 조건으로 $m<n_0$에 대한
$\mathfrak{p}_m$들을 정의하면
$\mathfrak{p}=\sum_{n=0}^{\infty}\mathfrak{p}_n$이 소아이디얼임을 보이는
일만 남는다. 먼저 $2^\circ$에 따라, $m\geq n_0$에 대해서도
$\mathfrak{p}_m$은 유한 개의 값을 제외한 모든 $r$에 대하여
$f^rx\in\mathfrak{p}_{m+rd}$를 만족하는 $x\in S_m$들의 집합으로
정의됨에 유의하자. 이제
\oldpage[II]{23}
$g\in S_m$, $x\in\mathfrak{p}_n$이면, 유한 개의 값을 제외한 모든
$r$에 대하여 $f^rgx\in\mathfrak{p}_{m+n+rd}$이다. 따라서
$gx\in\mathfrak{p}_{m+n}$이고, 이는 $\mathfrak{p}$가 $S$의
아이디얼임을 증명한다. 이 아이디얼이 소아이디얼임을, 다시 말해 부분군
$S_n/\mathfrak{p}_n$들로 등급이 주어진 환 $S/\mathfrak{p}$가 정역임을
보이기 위해서는 ($S/\mathfrak{p}$의 두 원소의 최고차 동차성분들을
고려하여) $x\in S_m$, $y\in S_n$이
$x\notin\mathfrak{p}_m$, $y\notin\mathfrak{p}_n$을 만족하면
$xy\notin\mathfrak{p}_{m+n}$임을 보이면 충분하다. 그렇지 않다면 충분히
큰 $r$에 대하여 $f^{2r}xy\in\mathfrak{p}_{m+n+2rd}$일 것이다. 그러나
모든 $r>0$에 대하여 $f^ry\notin\mathfrak{p}_{n+rd}$이다. 그러면
$2^\circ$에 따라 유한 개의 값을 제외한 모든 $r$에 대하여
$f^rx\in\mathfrak{p}_{m+rd}$이고, 이로부터
$x\in\mathfrak{p}_m$을 얻어 가정에 모순이다.
\end{proof}

\begin{env}[2.1.10]
\phantomsection
\label{II.2.1.10-ko}
$S_+$의 부분집합 $\mathfrak{J}$가 $S$의 아이디얼이면 이를
\emph{$S_+$의 아이디얼}이라 한다. 또한 $\mathfrak{J}$가 $S_+$와,
\emph{$S_+$를 포함하지 않는} $S$의 동차 소아이디얼의 교집합이면 이를
\emph{$S_+$의 동차 소아이디얼}이라 한다(더욱이 이 소아이디얼은
(\hyperref[II.2.1.9-ko]{2.1.9})에 따라 유일하다). $\mathfrak{J}$가
$S_+$의 아이디얼이면, \emph{$S_+$에서의 $\mathfrak{J}$의 근기}는 어떤
거듭제곱이 $\mathfrak{J}$에 속하는 $S_+$의 원소들로 이루어진 집합이다.
다시 말해 이는 $r_+(\mathfrak{J})=r(\mathfrak{J})\cap S_+$이다. 특히
$S_+$에서 $0$의 근기는 다시 $S_+$의 \emph{멱영근기}라 하고
$\mathfrak{N}_+$로 나타낸다. 따라서 이는 $S_+$의 멱영원 전체의
집합이다. $\mathfrak{J}$가 $S_+$의 \emph{동차} 아이디얼이면 그 근기
$r_+(\mathfrak{J})$도 동차 아이디얼이다. 실제로 몫환
$S/\mathfrak{J}$로 넘어가면 $\mathfrak{J}=0$인 경우로 환원할 수 있고,
$x=x_h+x_{h+1}+\cdots+x_k$가 멱영원이면
$x_i\in S_i$ ($1\leq h\leq i\leq k$)들도 모두 멱영원임을 보이면
충분하다. $x_k\neq 0$이라고 가정할 수 있고, 이때 $x^n$의 최고차
성분은 $x_k^n$이므로 $x_k$는 멱영원이다. 이제 $k$에 대한 귀납법을
적용한다. 등급환 $S$가 \emph{본질적으로 축소}라는 것은
$\mathfrak{N}_+=0$이라는 뜻이다. 다시 말해 $S_+$가 $\neq 0$인 멱영원을
포함하지 않는다는 뜻이다.
\end{env}

\begin{env}[2.1.11]
\phantomsection
\label{II.2.1.11-ko}
등급환 $S$에서 원소 $x$가 어떤 영이 아닌 원소와 곱하여 $0$이 되는
영인자이면 그 최고차 동차성분도 영인자임에 유의하자. 환 $S$가
\emph{본질적으로 정역}이라는 것은 $S_+$가
(\emph{단위원을 갖지 않는 환으로서}) 영인자를 포함하지 않아 영이 아닌
두 원소의 곱이 $0$일 수 없고, 또한 $\neq 0$인 것을 뜻한다. 이를 위해서는
$S_+$의 동차 원소 가운데 $\neq 0$인 것은 모두 이 환에서 영인자가 아니어서
어떤 영이 아닌 원소와 곱해도 $0$이 되지 않으면 충분하다.
$\mathfrak{p}$가 $S_+$의
동차 소아이디얼이면 $S/\mathfrak{p}$가 본질적으로 정역임은 명백하다.

$S$를 본질적으로 정역인 등급환이라 하고 $x_0\in S_0$라 하자. 이때
$S_+$의 동차 원소 $f\neq 0$가 \emph{하나라도} 존재하여
$x_0f=0$이면 $x_0S_+=0$이다. 실제로 모든 $g\in S_+$에
대하여 $(x_0g)f=(x_0f)g=0$이고, 가정에 따라 $x_0g=0$이다. 따라서
$S$가 정역이기 위한 필요충분조건은 $S_0$가 정역이고 $S_0$에서
$S_+$의 소멸자가 $0$뿐인 것이다.
\end{env}

\subsection{등급환의 분수환.}
\label{subsection:II.2.2-ko}

\begin{env}[2.2.1]
\phantomsection
\label{II.2.2.1-ko}
$S$를 음이 아닌 차수로 등급이 주어진 환이라 하고, $f$를 차수가
$d>0$인 $S$의 \emph{동차} 원소라 하자. 그러면 분수환 $S'=S_f$는
$S'_n$을 $x\in S_{n+kd}$이고 $k\geq 0$일 때의 $x/f^k$ 전체의
집합으로 놓음으로써 등급환이 된다(여기서 $n$은 임의의 음의 값도 취할
수 있음에 유의하자). $S'$의 \emph{차수가 $0$인} 원소들로 이루어진
부분환 $S'_0=(S_f)_0$을 $S_{(f)}$로 나타내겠다.

$f\in S_d$이면 $S_f$의 단항식 $(f/1)^h$들($h$는 영을 포함한 임의의
정수)은 환 $S_{(f)}$ 위의 \emph{자유계}를 이루며, 그 선형결합 전체의
집합은 바로
\oldpage[II]{24}
환 $(S^{(d)})_f$이다. 따라서 이 환은
\emph{$S_{(f)}[T,T^{-1}]=S_{(f)}\otimes_{\mathbf{Z}}\mathbf{Z}[T,T^{-1}]$와
동형이다}($T$는 부정원이다). 실제로
$\sum_{h=-a}^{b}z_h(f/1)^h=0$인 관계가 있고 $z_h=x_h/f^m$이며 모든
$x_h$가 $S_{md}$에 속한다고 하자. 정의에 따라 이 관계는
어떤 $k\geq a$가 존재하여
$\sum_{h=-a}^{b}f^{h+k}x_h=0$이 되는 것과 동치이다. 이 합의 항들은
차수가 서로 다르므로 모든 $h$에 대하여 $f^{h+k}x_h=0$이고, 따라서
모든 $h$에 대하여 $z_h=0$이다.

$M$이 등급 $S$-가군이면 $M'=M_f$는 등급 $S_f$-가군이다. 여기서
$M'_n$은 $z/f^k$ 꼴의 원소들로 이루어진 집합이며, 여기서
$z\in M_{n+kd}$이고 $k\geq 0$이다.
$M'$의 차수가 $0$인 동차 원소들의 집합을 $M_{(f)}$로 나타낸다.
$M_{(f)}$가 $S_{(f)}$-가군이고
$(M^{(d)})_f=M_{(f)}\otimes_{S_{(f)}}(S^{(d)})_f$임은 바로 알 수 있다.
\end{env}

\begin{lemma}[2.2.2]
\phantomsection
\label{II.2.2.2-ko}
$d$, $e$를 $>0$인 두 정수, $f\in S_d$, $g\in S_e$라 하자. 다음과 같은
환의 표준 동형이 존재한다.
\[
  S_{(fg)}\xrightarrow{\sim}(S_{(f)})_{g^d/f^e}~;
\]
이 두 환을 표준적으로 동일시하면, 다음과 같은 가군의 표준 동형이
존재한다.
\[
  M_{(fg)}\xrightarrow{\sim}(M_{(f)})_{g^d/f^e}.
\]
\end{lemma}

\begin{proof}
실제로 $fg$는 $f^e g^d$의 약수이고, 후자는 $(fg)^{de}$의 약수이므로
등급환 $S_{fg}$와 $S_{f^e g^d}$는 표준적으로 동일시된다. 다른 한편
$S_{f^e g^d}$는 $(S_{f^e})_{g^d/1}$와도 동일시되고
(\hyperref[0.1.4.6-ko]{0, I.4.6}), $f^e/1$은 $S_{f^e}$에서 가역이므로
$S_{f^e g^d}$는 $(S_{f^e})_{g^d/f^e}$와도 동일시된다. 그런데
$g^d/f^e$는 차수가 $0$인 $S_{f^e}$의 원소이다. 따라서
$(S_{f^e})_{g^d/f^e}$에서 차수가 $0$인 원소들로 이루어진 부분환은
$(S_{(f^e)})_{g^d/f^e}$임을 곧바로 얻는다. 또 명백히
$S_{(f^e)}=S_{(f)}$이므로, 이는 명제의 첫 번째 부분을 증명한다.
두 번째 부분도 같은 방식으로 확립된다.
\end{proof}

\begin{env}[2.2.3]
\phantomsection
\label{II.2.2.3-ko}
(\hyperref[II.2.2.2-ko]{2.2.2})의 가정 아래에서, 표준 준동형
$S_f\to S_{fg}$ (\hyperref[0.1.4.1-ko]{0, I.4.1})은 $x/f^k$를
$g^k x/(fg)^k$로 보내며 차수가 $0$임은 명백하다. 따라서
제한함으로써
\emph{표준 준동형 $S_{(f)}\to S_{(fg)}$}을 얻으며, 다음 도식은
가환한다.
\[
  \xymatrix{
    & S_{(f)}\ar[dl]\ar[dr]\\
    S_{(fg)}\ar[rr]^-{\sim} && (S_{(f)})_{g^d/f^e}
  }
\]
마찬가지로 표준 준동형 $M_{(f)}\to M_{(fg)}$을 정의한다.
\end{env}

\begin{lemma}[2.2.4]
\phantomsection
\label{II.2.2.4-ko}
$f$, $g$가 $S_+$의 두 동차 원소이면, 환 $S_{(fg)}$는 $S_{(f)}$와
$S_{(g)}$의 표준 상들의 합집합에 의해 생성된다.
\end{lemma}

\begin{proof}
(\hyperref[II.2.2.2-ko]{2.2.2})에 따라
$1/(g^d/f^e)=f^{d+e}/(fg)^d$가 $S_{(g)}$의 $S_{(fg)}$ 안에서의 표준
상에 속함을 보이면 충분하며, 이는 정의상 명백하다.
\end{proof}

\begin{proposition}[2.2.5]
\phantomsection
\label{II.2.2.5-ko}
$d$를 $>0$인 정수, $f\in S_d$라 하자. 그러면 다음과 같은 환의
표준 동형이 존재한다.
$S_{(f)}\xrightarrow{\sim}S^{(d)}/(f-1)S^{(d)}$~; 이 동형으로 두 환을
표준적으로 동일시하면, 다음과 같은 가군의 표준 동형이 존재한다.
$M_{(f)}\xrightarrow{\sim}M^{(d)}/(f-1)M^{(d)}$.
\end{proposition}

\begin{proof}
이 동형들 중 첫 번째 것은 $x/f^n$($x\in S_{nd}$)에 원소
$\overline{x}$, 곧 $x$의 $(f-1)S^{(d)}$에 대한 잉여류를 대응시켜 정의한다.
이 사상은 잘 정의된다. 실제로 모든 양의 정수 지수에 대하여 합동식
$f^h x\equiv x\pmod{(f-1)S^{(d)}}$이 모든 $x\in S^{(d)}$에 대해
성립한다. 따라서 $f^h x=0$인 $h>0$가 있으면
\oldpage[II]{25}
$\overline{x}=0$이다. 한편 $x\in S_{nd}$이고 $x=(f-1)y$라 하자. 여기서
$y=y_{hd}+y_{(h+1)d}+\cdots+y_{kd}$이고 $y_{jd}\in S_{jd}$이며
$y_{hd}\neq 0$이다. 그러면 반드시 $h=n$이고 $x=-y_{hd}$이며, 또한
관계 $y_{(j+1)d}=fy_{jd}$가 $h\leq j\leq k-1$에 대해 성립하고
$fy_{kd}=0$이다.
이로부터 마침내 $f^{k-n+1}x=0$을 얻는다. 따라서 $x\in S_{nd}$인 원소가
나타내는 $(f-1)S^{(d)}$에 대한 각 잉여류 $\overline{x}$에 $S_{(f)}$의 원소
$x/f^n$을 대응시킬 수 있다.
앞의 고찰에 따라 이 사상은 잘 정의된다. 이렇게 정의된 두 사상이
환 준동형이며 서로 역임을 곧바로 알 수 있다. $M$에 대해서도 완전히
같은 방식으로 한다.
\end{proof}

\begin{corollary}[2.2.6]
\phantomsection
\label{II.2.2.6-ko}
$S$가 뇌터이면, 차수가 $>0$인 임의의 동차 원소 $f$에 대하여
$S_{(f)}$도 뇌터이다.
\end{corollary}

\begin{proof}
이는 (\hyperref[II.2.1.7-ko]{2.1.7})과
(\hyperref[II.2.2.5-ko]{2.2.5})에서 곧바로 따른다.
\end{proof}

\begin{env}[2.2.7]
\phantomsection
\label{II.2.2.7-ko}
$T$를 \emph{동차} 원소들로 이루어진 $S_+$의 곱셈적 부분집합이라 하자.
그러면 $T_0=T\cup\{1\}$은 $S$의 곱셈적 부분집합이다. $T_0$의
원소들은 동차이므로, 환 $T_0^{-1}S$에는 명백한 방식으로 여전히 등급이
주어진다. $S_{(T)}$로 $T_0^{-1}S$에서 차수가 $0$인 원소들, 즉
$x/h$ 꼴의 원소들로 이루어진 부분환을 나타내겠다. 여기서 $h\in T$이고
$x$는 $h$와 차수가 같은 동차 원소이다.
(\hyperref[0.1.4.5-ko]{0, 1.4.5})에 따라 $T_0^{-1}S$는 환
$S_f$들($f$가 $T$를 달리며 전이사상은 표준 준동형
$S_f\to S_{fg}$이다)의 귀납극한과 표준적으로 동일시됨이 알려져 있다.
이 동일시는 차수를 보존하므로 $S_{(T)}$를
$S_{(f)}$들의($f\in T$) \emph{귀납극한}과 동일시한다. 모든 등급
$S$-가군 $M$에 대하여도 마찬가지로 가군 $M_{(T)}$(환 $S_{(T)}$ 위의
가군)를 차수가 $0$인 $T_0^{-1}M$의 원소들로 이루어진 것으로 정의한다.
이 가군은 $M_{(f)}$들의($f\in T$) 귀납극한임을 알 수 있다.

$\mathfrak{p}$가 $S_+$의 동차 소아이디얼이면, 각각
$S_{(\mathfrak{p})}$와 $M_{(\mathfrak{p})}$로 환 $S_{(T)}$와 가군
$M_{(T)}$를 나타내겠다. 여기서 $T$는 $S_+$에서 $\mathfrak{p}$에
속하지 않는 \emph{동차} 원소들의 집합이다.
\end{env}

\subsection{등급환의 동차 소 스펙트럼.}
\label{subsection:II.2.3-ko}

\begin{env}[2.3.1]
\phantomsection
\label{II.2.3.1-ko}
음이 아닌 차수로 등급이 주어진 환 $S$에 대하여, $S$의
\emph{동차 소 스펙트럼}이라 하고 $\operatorname{Proj}(S)$로 나타내는 것은
$S_+$의 동차 소아이디얼들의 집합
(\hyperref[II.2.1.10-ko]{2.1.10}), 또는 같은 말로 $S$의 동차
소아이디얼 가운데 $S_+$를 \emph{포함하지 않는} 것들의 집합이다. 이제
$\operatorname{Proj}(S)$를 바탕집합으로 갖는 \emph{스킴} 구조를 정의하겠다.
\end{env}

\begin{env}[2.3.2]
\phantomsection
\label{II.2.3.2-ko}
$E$를 $S$의 임의의 부분집합이라 하고, $V_+(E)$를 $S$의 동차
소아이디얼 가운데 $E$를 포함하고 $S_+$를 포함하지 않는 것들의 집합이라
하자. 그러므로 이는 $V(E)\cap\operatorname{Proj}(S)$라는
$\operatorname{Spec}(S)$의 부분집합이다. 따라서
(\hyperref[I.1.1.2-ko]{I, 1.1.2})에서 다음을 얻는다.
\[
  \label{II.2.3.2.1-ko}
  V_+(0)=\operatorname{Proj}(S),\quad
  V_+(S)=V_+(S_+)=\varnothing
  \tag{2.3.2.1}
\]
\[
  \label{II.2.3.2.2-ko}
  V_+\left(\bigcup_\lambda E_\lambda\right)
  =\bigcap_\lambda V_+(E_\lambda)
  \tag{2.3.2.2}
\]
\[
  \label{II.2.3.2.3-ko}
  V_+(EE')=V_+(E)\cup V_+(E')
  \tag{2.3.2.3}
\]
$V_+(E)$는 집합 $E$를 $E$가 생성하는 동차 아이디얼로 바꾸어도 변하지
않는다. 또한 $\mathfrak{J}$가 $S$의 동차 아이디얼이면 다음이 성립한다.
\[
  \label{II.2.3.2.4-ko}
  V_+(\mathfrak{J})
  =V_+\left(\bigcup_{q\geq n}(\mathfrak{J}\cap S_q)\right)
  \tag{2.3.2.4}
\]
\oldpage[II]{26}
이는 모든 $n>0$에 대하여 성립한다. 실제로
$\mathfrak{p}\in\operatorname{Proj}(S)$가 $\mathfrak{J}$의 차수가
$\geq n$인 동차 원소들을 포함한다고 하자. 가정에 따라 동차 원소
$f\in S_d$ 가운데 $\mathfrak{p}$에 속하지 않는 것이 존재하므로, 모든
$m\geq 0$과 모든 $x\in S_m\cap\mathfrak{J}$에 대하여 관계
$f^r x\in\mathfrak{J}\cap S_{m+rd}$는 모든 음이 아닌 정수 $r$에 대해
성립한다. 한편 유한 개의 값을 제외한 모든 지수에 대해서는
$m+rd\geq n$이므로 위 교집합은 해당 소아이디얼에 포함되고, 따라서
$f^r x\in\mathfrak{p}\cap S_{m+rd}$이다. 앞서 고른 동차 원소가 해당
소아이디얼에 속하지 않으므로 소아이디얼의 성질로 그 거듭제곱을 소거하면
$x\in\mathfrak{p}\cap S_m$을 얻는다
(\hyperref[II.2.1.8-ko]{2.1.8}).

마지막으로 $\mathfrak{J}$가 $S$의 임의의 동차 아이디얼이면
\[
  \label{II.2.3.2.5-ko}
  V_+(\mathfrak{J})=V_+(r_+(\mathfrak{J})).
  \tag{2.3.2.5}
\]
\end{env}

\begin{env}[2.3.3]
\phantomsection
\label{II.2.3.3-ko}
정의에 따라 $V_+(E)$ 꼴의 집합들은 $X=\operatorname{Proj}(S)$ 위에서
$\operatorname{Spec}(S)$의 스펙트럼 위상으로부터 유도된 위상의
닫힌집합들이다. 이 위상도 $X$ 위의 \emph{스펙트럼 위상}이라 하겠다.
모든 $f\in S$에 대하여 다음과 같이 놓는다.
\[
  \label{II.2.3.3.1-ko}
  D_+(f)=D(f)\cap\operatorname{Proj}(S)
  =\operatorname{Proj}(S)\setminus V_+(f)
  \tag{2.3.3.1}
\]
따라서 $f$, $g$가 $S$의 두 원소이면
(\hyperref[I.1.1.9.1-ko]{I, 1.1.9.1})
\[
  \label{II.2.3.3.2-ko}
  D_+(fg)=D_+(f)\cap D_+(g).
  \tag{2.3.3.2}
\]
\end{env}

\begin{proposition}[2.3.4]
\phantomsection
\label{II.2.3.4-ko}
$f$가 $S_+$의 동차 원소 전체를 달릴 때, $D_+(f)$들은
$X=\operatorname{Proj}(S)$의 위상의 기저를 이룬다.
\end{proposition}

\begin{proof}
실제로 (\hyperref[II.2.3.2.2-ko]{2.3.2.2})와
(\hyperref[II.2.3.2.4-ko]{2.3.2.4})에 따라 $X$의 모든 닫힌 부분집합은
$V_+(f)$ 꼴 집합들의 교집합이다. 여기서 $f$는 차수가 $>0$인 동차
원소이다.
\end{proof}

\begin{env}[2.3.5]
\phantomsection
\label{II.2.3.5-ko}
$f$를 $S_+$의 차수가 $d>0$인 \emph{동차} 원소라 하자.
$\mathfrak{p}$가 $S$의 동차 소아이디얼이고 $f$를 포함하지 않을 때,
$x/f^n$ 꼴 원소
가운데 $x\in\mathfrak{p}$이고 $n\geq 0$인 것들의 집합은 분수환
$S_f$의 소아이디얼임이 알려져 있다
(\hyperref[0.1.2.6-ko]{0, 1.2.6}). 그 아이디얼과 $S_{(f)}$의 교집합도
이 환의 소아이디얼이므로 이를 $\psi_f(\mathfrak{p})$로 나타내겠다.
이는 $x/f^n$ 가운데 $n\geq 0$이고
$x\in\mathfrak{p}\cap S_{nd}$인 것들의 집합이다. 이로써 사상
\[
  \psi_f:D_+(f)\longrightarrow\operatorname{Spec}(S_{(f)})~;
\]
을 정의하였다. 또한 $g\in S_e$가 $S_+$의 또 하나의 동차 원소이면 다음
도식은 가환한다.
\[
  \label{II.2.3.5.1-ko}
  \xymatrix{
    D_+(f)\ar[r]^{\psi_f}
    & \operatorname{Spec}(S_{(f)})\\
    D_+(fg)\ar[u]\ar[r]_{\psi_{fg}}
    & \operatorname{Spec}(S_{(fg)})\ar[u]
  }
  \tag{2.3.5.1}
\]
여기서 왼쪽 수직 화살표는 포함 사상이고, 오른쪽 수직 화살표는 사상
${}^a\omega_{fg,f}$이다. 이는 표준 준동형
$\omega=\omega_{fg,f}:S_{(f)}\to S_{(fg)}$에서 유도된다
(\hyperref[I.1.2.1-ko]{I, 1.2.1}). 실제로
$x/f^n\in\omega^{-1}(\psi_{fg}(\mathfrak{p}))$이고
$fg\notin\mathfrak{p}$라 하자. 정의에 따라
$g^n x/(fg)^n\in\psi_{fg}(\mathfrak{p})$이므로
$g^n x\in\mathfrak{p}$이고, 따라서 $x\in\mathfrak{p}$이다. 역은
자명하다.
\end{env}

\begin{proposition}[2.3.6]
\phantomsection
\label{II.2.3.6-ko}
사상 $\psi_f$는 $D_+(f)$에서 $\operatorname{Spec}(S_{(f)})$로의
위상동형이다.
\end{proposition}

\begin{proof}
먼저 $\psi_f$는 연속이다. 실제로 $h\in S_{nd}$이고
$h/f^n\in\psi_f(\mathfrak{p})$이면 정의에 따라
$h\in\mathfrak{p}$이며, 그 역도 성립한다. 따라서
\[
  \psi_f^{-1}(D(h/f^n))=D_+(hf)
\]
이고, 주장은 공식
(\hyperref[II.2.3.3.2-ko]{2.3.3.2})에서 따른다. 또한 $D_+(hf)$들, 즉
$h$가 여러 집합 $S_{nd}$의 원소 전체를 달릴 때 얻는 이 집합들은,
(\hyperref[II.2.3.4-ko]{2.3.4})와 공식
(\hyperref[II.2.3.3.2-ko]{2.3.3.2})에 따라 $D_+(f)$의 위상의 기저를
이룬다.
\oldpage[II]{27}
그러므로 공리 $(T_0)$가 $D_+(f)$와 $\operatorname{Spec}(S_{(f)})$에서
성립함을 위의 논의와 함께 고려하면, $\psi_f$는 단사이고 그 역인 사상
$\psi_f(D_+(f))\longrightarrow D_+(f)$는 연속이다. 마지막으로
$\psi_f$가 전사임을 보이자. $\mathfrak{q}_0$를 $S_{(f)}$의
소아이디얼이라 하고, 모든 $n>0$에 대하여 $\mathfrak{p}_n$을
$x\in S_n$ 가운데 $x^d/f^n\in\mathfrak{q}_0$인 것들의 집합이라 하자.
그러면 $\mathfrak{p}_n$들은 (\hyperref[II.2.1.9-ko]{2.1.9})의 조건들을
만족한다. 실제로 $x\in S_n$, $y\in S_n$이고
$x^d/f^n\in\mathfrak{q}_0$, $y^d/f^n\in\mathfrak{q}_0$라 하자. 그러면
$(x+y)^{2d}/f^{2n}\in\mathfrak{q}_0$이고, 따라서
$(x+y)^d/f^n\in\mathfrak{q}_0$이다. 이는 $\mathfrak{q}_0$가
소아이디얼이기 때문이다. 이로써
$\mathfrak{p}_n$들이 $S_n$의 부분군임이 증명된다.
(\hyperref[II.2.1.9-ko]{2.1.9})의 다른 조건들도 $\mathfrak{q}_0$가
소아이디얼임을 고려하면 자명하게 확인된다. 이렇게 정의한 $\mathfrak{p}$는 $S$의
동차 소아이디얼이고 실제로
$\psi_f(\mathfrak{p})=\mathfrak{q}_0$이다. 왜냐하면 $x\in S_{nd}$일 때
두 관계 $x/f^n\in\mathfrak{q}_0$와
$x^d/f^{nd}\in\mathfrak{q}_0$는 $\mathfrak{q}_0$가 소아이디얼이므로
서로 동치이기 때문이다.
\end{proof}

\begin{corollary}[2.3.7]
\phantomsection
\label{II.2.3.7-ko}
$D_+(f)=\varnothing$이기 위한 필요충분조건은 $f$가 멱영원인 것이다.
\end{corollary}

\begin{proof}
실제로 $\operatorname{Spec}(S_{(f)})=\varnothing$이기 위한 필요충분조건은
$S_{(f)}=0$인 것이며, 이는 다시 $1=0$이 $S_f$에서 성립하는 것과 동치이다.
정의에 따라 이는 $f$가 멱영원이라는 뜻이다.
\end{proof}

\begin{corollary}[2.3.8]
\phantomsection
\label{II.2.3.8-ko}
$E$를 $S_+$의 부분집합이라 하자. 다음 조건들은 서로 동치이다.
\begin{enumerate}
  \item[\rm a)] $V_+(E)=X=\operatorname{Proj}(S)$.
  \item[\rm b)] $E$의 모든 원소는 멱영원이다.
  \item[\rm c)] $E$의 각 원소의 모든 동차 성분은 멱영원이다.
\end{enumerate}
\end{corollary}

\begin{proof}
c)가 b)를 함의하고 b)가 a)를 함의함은 명백하다. $\mathfrak{J}$가
$S$의 동차 아이디얼이고 $E$로 생성된다면, 조건 a)는
$V_+(\mathfrak{J})=X$와 동치이다. \emph{특히}, a)는 모든 동차 원소
$f\in\mathfrak{J}$가 $V_+(f)=X$를 만족함을
함의하므로, (\hyperref[II.2.3.7-ko]{2.3.7})에 따라 $f$는 멱영원이다.
\end{proof}

\begin{corollary}[2.3.9]
\phantomsection
\label{II.2.3.9-ko}
$\mathfrak{J}$가 $S_+$의 동차 아이디얼이면,
$r_+(\mathfrak{J})$는 $S_+$의 동차 소아이디얼들 가운데
$\mathfrak{J}$를 포함하는 것들의 교집합이다.
\end{corollary}

\begin{proof}
등급환 $S/\mathfrak{J}$를 고려하여 $\mathfrak{J}=0$인 경우로 환원할 수
있다. $f\in S_+$가 멱영원이 아니면 $S$의 동차 소아이디얼 가운데
$f$를 포함하지 않는 것이 존재함을 증명해야 한다. 그런데 $f$의 동차
성분 중 적어도 하나는 멱영원이 아니므로 $f$가 동차라고 가정할 수 있고,
이때 주장은 (\hyperref[II.2.3.7-ko]{2.3.7})에서 따른다.
\end{proof}

\begin{env}[2.3.10]
\phantomsection
\label{II.2.3.10-ko}
$Y$를 $X=\operatorname{Proj}(S)$의 임의의 부분집합이라 하자.
$\mathfrak{j}_+(Y)$를 모든 $f\in S_+$ 가운데 $Y\subset V_+(f)$를
만족하는 것들의 집합으로 둔다. 다시 말해,
$\mathfrak{j}_+(Y)=\mathfrak{j}(Y)\cap S_+$이다. 따라서
$\mathfrak{j}_+(Y)$는 $S_+$의 아이디얼이며, $S_+$에서 취한 그 근기와
같다.
\end{env}

\begin{proposition}[2.3.11]
\phantomsection
\label{II.2.3.11-ko}
\begin{enumerate}
  \item[\rm(i)] $E$를 $S_+$의 임의의 부분집합이라 하면,
    $\mathfrak{j}_+(V_+(E))$는 $S_+$에서 취한, $S_+$에서 $E$가 생성하는
    동차 아이디얼의 근기이다.
  \item[\rm(ii)] $Y$를 $X$의 임의의 부분집합이라 하면,
    $V_+(\mathfrak{j}_+(Y))=\overline{Y}$이다. 즉 우변은 $Y$의 폐포이며,
    이 폐포는 $X$ 안에서 취한다.
\end{enumerate}
\end{proposition}

\begin{proof}
\begin{enumerate}
  \item[\rm(i)] $\mathfrak{J}$가 $S_+$의 동차 아이디얼이고 $E$로
    생성된다면 $V_+(E)=V_+(\mathfrak{J})$이고, 이때 주장은
    (\hyperref[II.2.3.9-ko]{2.3.9})에서 따른다.
  \item[\rm(ii)] $V_+(\mathfrak{J})=\bigcap_{f\in\mathfrak{J}}V_+(f)$이므로,
    관계 $Y\subset V_+(\mathfrak{J})$는 $Y\subset V_+(f)$가 모든
    $f\in\mathfrak{J}$에 대하여 성립함을 함의한다. 따라서
    $\mathfrak{j}_+(Y)\supset\mathfrak{J}$이고, 이로부터
    $V_+(\mathfrak{j}_+(Y))\subset V_+(\mathfrak{J})$를 얻는다. 따라서
    닫힌 부분집합의 정의에 의해 (ii)가 증명된다.
\end{enumerate}
\end{proof}

\begin{corollary}[2.3.12]
\phantomsection
\label{II.2.3.12-ko}
$Y$가 $X=\operatorname{Proj}(S)$의 닫힌 부분집합일 때, 그러한 부분집합
전체와 $S_+$의 동차 아이디얼 가운데 $S_+$에서 취한 자신의 근기와 같은
것 전체는 포함관계를 역전하는 두 사상
$Y\longmapsto\mathfrak{j}_+(Y)$,
$\mathfrak{J}\longmapsto V_+(\mathfrak{J})$에 의해 서로 전단사로 대응한다.
또 $X$의 두 닫힌 부분집합의 합집합 $Y_1\cup Y_2$에는
\oldpage[II]{28}
$\mathfrak{j}_+(Y_1)\cap\mathfrak{j}_+(Y_2)$가 대응하고, 임의의 닫힌
부분집합족 $(Y_\lambda)$의 교집합에는 $S_+$에서
$\mathfrak{j}_+(Y_\lambda)$들의 합의 근기가 대응한다.
\end{corollary}

\begin{corollary}[2.3.13]
\phantomsection
\label{II.2.3.13-ko}
$\mathfrak{J}$를 $S_+$의 동차 아이디얼이라 하자.
$V_+(\mathfrak{J})=\varnothing$이기 위한 필요충분조건은 $S_+$의 각
원소마다 그 원소의 양의 정수 거듭제곱 중 하나가 $\mathfrak{J}$에 속하는
것이다.
\end{corollary}

마지막 따름정리는 다음과 같은 서로 동치인 형태들로도 표현된다.

\begin{corollary}[2.3.14]
\phantomsection
\label{II.2.3.14-ko}
$(f_\alpha)$를 $S_+$의 동차 원소족이라 하자. $D_+(f_\alpha)$들이
$X=\operatorname{Proj}(S)$의 덮개를 이루기 위한 필요충분조건은 $S_+$의
각 원소마다 그 원소의 양의 정수 거듭제곱 중 하나가 $f_\alpha$들이
생성하는 아이디얼에 속하는 것이다.
\end{corollary}

\begin{corollary}[2.3.15]
\phantomsection
\label{II.2.3.15-ko}
$(f_\alpha)$를 $S_+$의 동차 원소족, $f$를 $S_+$의 한 원소라 하자.
다음 관계들은 서로 동치이다.
\begin{enumerate}
  \item[\rm a)] $D_+(f)\subset\bigcup_\alpha D_+(f_\alpha)$.
  \item[\rm b)] $V_+(f)\supset\bigcap_\alpha V_+(f_\alpha)$.
  \item[\rm c)] $f$의 어떤 양의 정수 거듭제곱이 $f_\alpha$들이 생성하는
    아이디얼에 속한다.
\end{enumerate}
\end{corollary}

\begin{corollary}[2.3.16]
\phantomsection
\label{II.2.3.16-ko}
$X=\operatorname{Proj}(S)$가 공집합이기 위한 필요충분조건은 $S_+$의
모든 원소가 멱영원인 것이다.
\end{corollary}

\begin{corollary}[2.3.17]
\phantomsection
\label{II.2.3.17-ko}
(\hyperref[II.2.3.12-ko]{2.3.12})에서 서술한 전단사 대응에서 $X$의
기약 닫힌 부분집합들에는 $S_+$의 동차 소아이디얼들이 대응한다.
\end{corollary}

\begin{proof}
실제로 $Y=Y_1\cup Y_2$이고 $Y_1$, $Y_2$가 $Y$의 진부분집합인
닫힌집합들이면
\[
  \mathfrak{j}_+(Y)
  =\mathfrak{j}_+(Y_1)\cap\mathfrak{j}_+(Y_2)
\]
이며, 아이디얼 $\mathfrak{j}_+(Y_1)$과 $\mathfrak{j}_+(Y_2)$는 둘 다
$\mathfrak{j}_+(Y)$와 다르므로 $\mathfrak{j}_+(Y)$는 소아이디얼이
아니다. 역으로 $\mathfrak{J}$가 $S_+$의 소아이디얼이 아닌 동차
아이디얼이면 두 원소 $f$, $g$가 $S_+$에 존재하여
$f\notin\mathfrak{J}$, $g\notin\mathfrak{J}$, $fg\in\mathfrak{J}$를
만족한다. 그러면
$V_+(f)\not\supset V_+(\mathfrak{J})$,
$V_+(g)\not\supset V_+(\mathfrak{J})$이지만
(\hyperref[II.2.3.2.3-ko]{2.3.2.3})에 따라
$V_+(\mathfrak{J})\subset V_+(f)\cup V_+(g)$이다. 따라서
$V_+(\mathfrak{J})$는 닫힌집합
$V_+(f)\cap V_+(\mathfrak{J})$와
$V_+(g)\cap V_+(\mathfrak{J})$의 합집합이고, 이 두 닫힌집합은 모두
$V_+(\mathfrak{J})$와 다르다.
\end{proof}

\subsection{$\operatorname{Proj}(S)$의 스킴 구조}
\label{subsection:II.2.4-ko}

\begin{env}[2.4.1]
\phantomsection
\label{II.2.4.1-ko}
$f$, $g$를 $S_+$의 두 동차 원소라 하자. 아핀 스킴
$Y_f=\operatorname{Spec}(S_{(f)})$,
$Y_g=\operatorname{Spec}(S_{(g)})$,
$Y_{fg}=\operatorname{Spec}(S_{(fg)})$를 생각하자.
(\hyperref[II.2.2.2-ko]{2.2.2})에 따라, 사상
$w_{fg,f}=({}^a\omega_{fg,f},\widetilde{\omega}_{fg,f})$는 $Y_{fg}$에서
$Y_f$로 가며, 표준 준동형
$\omega_{fg,f}:S_{(f)}\to S_{(fg)}$에 대응하고,
\emph{열린 몰입 사상}이다
(\hyperref[I.1.3.6-ko]{I, 1.3.6}).
위상동형 $\psi_f:D_+(f)\to Y_f$의 역
(\hyperref[II.2.3.6-ko]{2.3.6})을 이용하여 $D_+(f)$에 $Y_f$의 아핀
스킴 구조를 옮길 수 있다. 도식
(\hyperref[II.2.3.5.1-ko]{2.3.5.1})의 가환성에 따라, 이렇게 얻은 아핀
스킴 $D_+(fg)$는 아핀 스킴 $D_+(f)$의 바탕공간의 열린 부분집합
$D_+(fg)$에 유도된 스킴과 동일시된다. 그러면
(\hyperref[II.2.3.4-ko]{2.3.4})를 고려할 때
$X=\operatorname{Proj}(S)$에는 유일한 \emph{준스킴} 구조가 주어지며,
그 구조를 각 $D_+(f)$에 제한하면 방금 정의한 아핀 스킴이 된다. 또한
다음이 성립한다.
\end{env}

\begin{proposition}[2.4.2]
\phantomsection
\label{II.2.4.2-ko}
준스킴 $\operatorname{Proj}(S)$는 스킴이다.
\end{proposition}

\begin{proof}
(\hyperref[I.5.5.6-ko]{I, 5.5.6})에 따라, $f$, $g$가 $S_+$의 임의의
동차 원소일 때 $D_+(f)\cap D_+(g)=D_+(fg)$가 아핀이고 그 환이
$D_+(f)$와 $D_+(g)$의 환들의 표준 상들에 의해 생성됨을 보이면
충분하다. 첫째 주장은 정의상 명백하고, 둘째 주장은
(\hyperref[II.2.2.4-ko]{2.2.4})에서 따른다.
\end{proof}
\oldpage[II]{29}
$\operatorname{Proj}(S)$라는 동차 소 스펙트럼을 \emph{스킴}이라고 할
때에는 언제나 방금 정의한 구조를 뜻한다.

\begin{example}[2.4.3]
\phantomsection
\label{II.2.4.3-ko}
$S=K[T_1,T_2]$라 하자. 여기서 $K$는 체이고, $T_1$, $T_2$는 두
부정원이며, $S$에는 전체 차수에 따른 등급이 주어져 있다.
(\hyperref[II.2.3.14-ko]{2.3.14})에 따라
$\operatorname{Proj}(S)$는 $D_+(T_1)$과 $D_+(T_2)$의 합집합이다.
이 두 아핀 스킴은 각각 $\operatorname{Spec}(K[T])$와 표준적으로
동형이며, $\operatorname{Proj}(S)$는
(\hyperref[I.2.3.2-ko]{I, 2.3.2})에서 서술한 대로 이 두 아핀 스킴을
붙여서 얻어진다((7.4.14) 참조).
\end{example}

\begin{proposition}[2.4.4]
\phantomsection
\label{II.2.4.4-ko}
$S$를 음이 아닌 차수로 등급이 주어진 환이라 하고, $X$를 스킴
$\operatorname{Proj}(S)$라 하자.
\begin{enumerate}
  \item[\rm(i)] $\mathfrak{N}_+$가 $S_+$의 멱영근기이면
    (\hyperref[II.2.1.10-ko]{2.1.10}), 스킴 $X_{\mathrm{red}}$는
    $\operatorname{Proj}(S/\mathfrak{N}_+)$와 표준적으로 동형이다. 특히
    $S$가 본질적으로 축소이면 $\operatorname{Proj}(S)$는 축소이다.
  \item[\rm(ii)] $S$가 본질적으로 축소라고 가정하자. 이때 $X$가 정역
    스킴이기 위한 필요충분조건은 $S$가 본질적으로 정역인 것이다.
\end{enumerate}
\end{proposition}

\begin{proof}
\begin{enumerate}
  \item[\rm(i)] $\overline{S}$를 등급환 $S/\mathfrak{N}_+$라 하고, 차수가
    $0$인 표준 준동형 $S\to\overline{S}$를
    $x\mapsto\overline{x}$로 나타내자. 모든 $f\in S_d$ ($d>0$)에 대하여
    표준 준동형
    $S_f\to\overline{S}_{\overline{f}}$
    (\hyperref[0.1.5.1-ko]{0, 1.5.1})는 전사이고 차수가 $0$이므로, 이를
    제한하면 전사 준동형
    $S_{(f)}\to\overline{S}_{(\overline{f})}$를 얻는다.
    $f\notin\mathfrak{N}_+$라고 가정하면
    $\overline{S}_{(\overline{f})}$가 축소환이고 앞의 준동형의 핵이
    $S_{(f)}$의 멱영근기임을 바로 확인할 수 있다. 다시 말해
    $\overline{S}_{(\overline{f})}=(S_{(f)})_{\mathrm{red}}$이다. 따라서 이
    준동형에는 $D_+(\overline{f})$를 $(D_+(f))_{\mathrm{red}}$와 동일시하는
    닫힌 몰입 사상
    $D_+(\overline{f})\to D_+(f)$가 대응한다
    (\hyperref[I.5.1.2-ko]{I, 5.1.2}). 특히 이 사상은 두 아핀 스킴의
    바탕공간 사이의 위상동형이다. 더욱이
    $g\notin\mathfrak{N}_+$가 $S_+$의 또 다른 동차 원소이면 도식
    \[
      \xymatrix{
        S_{(f)}\ar[r]\ar[d]
        & \overline{S}_{(\overline{f})}\ar[d]\\
        S_{(fg)}\ar[r]
        & \overline{S}_{(\overline{fg})}
      }
    \]
    은 가환한다. 한편 $D_+(f)$들은 $f$가 차수가 $>0$인 동차 원소이고
    $f\notin\mathfrak{N}_+$일 때
    $X=\operatorname{Proj}(S)$의 덮개를 이룬다
    (\hyperref[II.2.3.7-ko]{2.3.7}). 따라서 사상
    $D_+(\overline{f})\to D_+(f)$들은 닫힌 몰입 사상
    $\operatorname{Proj}(\overline{S})\to\operatorname{Proj}(S)$의
    제한들이고, 이 닫힌 몰입 사상은 바탕공간 사이의 위상동형이다. 이제
    (\hyperref[I.5.1.2-ko]{I, 5.1.2})에 의해 결론이 따른다.
  \item[\rm(ii)] $S$가 본질적으로 정역이라고 가정하자. 다시 말해
    $(0)$은 $S_+$와 다른 $S_+$의 등급 소아이디얼이다. 그러면 (i)에 따라
    $X$는 축소이고, (\hyperref[II.2.3.17-ko]{2.3.17})에 따라 기약이다.
    반대로 $S$가 본질적으로 축소이고 $X$가 정역 스킴이라고 가정하자.
    그러면 $S_+$의 동차 원소 $f\neq0$에 대하여
    $D_+(f)\neq\varnothing$이다
    (\hyperref[II.2.3.7-ko]{2.3.7}). $X$가 기약이라는 가정에 따라
    $S_+$의 동차 원소 $f$, $g$가 모두 $\neq0$이면
    $D_+(f)\cap D_+(g)\neq\varnothing$이다. 따라서
    (\hyperref[II.2.3.3.2-ko]{2.3.3.2})에 의해 $fg\neq0$이고, 이로부터
    $S_+$에는 영인자가 없음을 얻으므로 처음의 주장이 따른다.
\end{enumerate}
\end{proof}

\begin{env}[2.4.5]
\phantomsection
\label{II.2.4.5-ko}
$A$를 가환환이라 하자. 등급환 $S$에 각 부분군 $S_n$이 $A$-가군이
되도록 하는 $A$-대수 구조가 주어져 있으면 이를
\emph{등급 $A$-대수}라고 한다는 것을 상기하자. 이를 위해서는 $S_0$가
\oldpage[II]{30}
$A$-대수이면 충분하다. 다시 말해 $S_0$에 $A$-대수 구조를 정의하고
$\alpha\in A$, $x\in S_n$에 대하여
$\alpha\cdot x=(\alpha\cdot1)x$로 놓으면 $S$에 등급 $A$-대수 구조가
정의된다.
\end{env}

\begin{proposition}[2.4.6]
\phantomsection
\label{II.2.4.6-ko}
$S$가 등급 $A$-대수라고 가정하자. 그러면
$X=\operatorname{Proj}(S)$ 위에서 구조층 $\mathcal{O}_X$는
$A$-대수들의 층이다($A$는 $X$ 위의 상수층으로 본다). 다시 말해 $X$는
$\operatorname{Spec}(A)$ 위의 스킴이다.
\end{proposition}

\begin{proof}
$S_+$의 모든 동차 원소 $f$에 대하여 $S_{(f)}$가 $A$ 위의 대수이고,
$S_+$의 동차 원소 $f$, $g$에 대하여 도식
\[
  \xymatrix{
    S_{(f)}\ar[rr] && S_{(fg)}\\
    & A\ar[ul]\ar[ur] &
  }
\]
이 가환함에 유의하면 충분하다.
\end{proof}

\begin{proposition}[2.4.7]
\phantomsection
\label{II.2.4.7-ko}
$S$를 음이 아닌 차수로 등급이 주어진 환이라 하자.
\begin{enumerate}
  \item[\rm(i)] 모든 정수 $d>0$에 대하여 스킴
    $\operatorname{Proj}(S)$에서 스킴
    $\operatorname{Proj}(S^{(d)})$로 가는 표준 동형이 존재한다.
  \item[\rm(ii)] $S'$를 다음 조건을 만족하는 등급환이라 하자:
    $S'_0=\mathbf{Z}$이고, $S'_n=S_n$은 ($\mathbf{Z}$-가군으로서)
    $n>0$일 때 성립한다. 그러면 스킴 $\operatorname{Proj}(S)$에서 스킴
    $\operatorname{Proj}(S')$로 가는 표준 동형이 존재한다.
\end{enumerate}
\end{proposition}

\begin{proof}
\begin{enumerate}
  \item[\rm(i)] 먼저 사상
    $\mathfrak{p}\mapsto\mathfrak{p}\cap S^{(d)}$가 집합
    $\operatorname{Proj}(S)$에서
    $\operatorname{Proj}(S^{(d)})$로 가는 전단사임을 보이자. 실제로 동차
    소아이디얼 $\mathfrak{p}'\in\operatorname{Proj}(S^{(d)})$가 주어졌다고
    하고 $\mathfrak{p}_{nd}=\mathfrak{p}'\cap S_{nd}$ ($n\geq0$)로 놓자.
    $n>0$이고 $d$의 배수가 아닌 모든 경우에 $\mathfrak{p}_n$을
    $x\in S_n$이고 $x^d\in\mathfrak{p}_{nd}$인 원소들의 집합으로 정의하자.
    $x\in\mathfrak{p}_n$, $y\in\mathfrak{p}_n$이면
    $(x+y)^{2d}\in\mathfrak{p}_{2nd}$이고, 따라서
    $(x+y)^d\in\mathfrak{p}_{nd}$이다. 이는 $\mathfrak{p}'$이
    소아이디얼이기 때문이다. 이와 같이 정의된 $\mathfrak{p}_n$들이 모든
    $n\geq0$에 대하여
    (\hyperref[II.2.1.9-ko]{2.1.9})의 조건을 만족함은 자명하다. 따라서
    유일한 소아이디얼 $\mathfrak{p}\in\operatorname{Proj}(S)$가 존재하며,
    이 소아이디얼은 $\mathfrak{p}\cap S^{(d)}=\mathfrak{p}'$을 만족한다. 한편 $f$가
    $S_+$의 동차 원소일 때마다
    $V_+(f)=V_+(f^d)$이므로
    (\hyperref[II.2.3.2.3-ko]{2.3.2.3}), 앞의 전단사는 위상공간들의
    위상동형이다. 마지막으로 같은 표기 아래에서 $S_{(f)}$와
    $S^{(d)}_{(f^d)}$는 표준적으로 동일시되므로
    (\hyperref[II.2.2.2-ko]{2.2.2}), $\operatorname{Proj}(S)$와
    $\operatorname{Proj}(S^{(d)})$도 스킴으로서 표준적으로 동일시된다.
  \item[\rm(ii)] 각 $\mathfrak{p}\in\operatorname{Proj}(S)$에 유일한
    소아이디얼 $\mathfrak{p}'\in\operatorname{Proj}(S')$을 대응시키되,
    $\mathfrak{p}'\cap S_n=\mathfrak{p}\cap S_n$이 모든 $n>0$에 대하여
    성립하게 하면
    (\hyperref[II.2.1.9-ko]{2.1.9}), 바탕공간들의 표준 위상동형
    $\operatorname{Proj}(S)\xrightarrow{\sim}\operatorname{Proj}(S')$을
    정의함이 분명하다. 실제로 $V_+(f)$는 $S$에 대해서나 $S'$에 대해서나
    같은 집합이다($f$가 $S_+$의 동차 원소일 때). 더욱이
    $S_{(f)}=S'_{(f)}$이므로 $\operatorname{Proj}(S)$와
    $\operatorname{Proj}(S')$은 스킴으로서 동일시된다.
\end{enumerate}
\end{proof}

\begin{corollary}[2.4.8]
\phantomsection
\label{II.2.4.8-ko}
$S$가 등급 $A$-대수이고, $S'_A$가 다음 조건을 만족하는 등급
$A$-대수라고 하자: $(S'_A)_0=A$, $(S'_A)_n=S_n$ ($n>0$).
$\operatorname{Proj}(S)$에서 $\operatorname{Proj}(S'_A)$로 가는 표준
동형이 존재한다.
\end{corollary}

\begin{proof}
실제로 (\hyperref[II.2.4.7-ko]{2.4.7, (ii)})의 표기 아래에서 이 두
스킴은 모두 $\operatorname{Proj}(S')$와 표준적으로 동형이다.
\end{proof}

\subsection{등급가군에 대응하는 층.}
\label{subsection:II.2.5-ko}

\begin{env}[2.5.1]
\phantomsection
\label{II.2.5.1-ko}
$M$을 등급 $S$-가군이라 하자. $f$가 $S_+$의 동차 원소이면
$M_{(f)}$는 $S_{(f)}$-가군이고, 따라서 이 가군에 대응하는 준연접층
$\widetilde{M_{(f)}}$가 아핀 스킴
$\operatorname{Spec}(S_{(f)})$ 위에 있다. 이 스킴은 $D_+(f)$와
동일시된다(\hyperref[I.1.3.4-ko]{I, 1.3.4}).
\end{env}
\oldpage[II]{31}

\begin{proposition}[2.5.2]
\phantomsection
\label{II.2.5.2-ko}
$X=\operatorname{Proj}(S)$ 위에 다음 조건을 만족하는 준연접
$\mathcal{O}_X$-가군 $\widetilde{M}$이 유일하게 존재한다. 즉 $f$가
$S_+$의 동차 원소일 때마다
$\Gamma(D_+(f),\widetilde{M})=M_{(f)}$이고, 제한 준동형
$\Gamma(D_+(f),\widetilde{M})\to
\Gamma(D_+(fg),\widetilde{M})$은 $f$, $g$가 $S_+$의 동차 원소일 때
표준 준동형 $M_{(f)}\to M_{(fg)}$이다
(\hyperref[II.2.2.3-ko]{2.2.3}).
\end{proposition}

\begin{proof}
$f\in S_d$, $g\in S_e$라 하자.
(\hyperref[II.2.2.2-ko]{2.2.2})에 따라 $D_+(fg)$는
$(S_{(f)})_{g^d/f^e}$의 소 스펙트럼과 동일시된다. 따라서
$D_+(fg)$로 제한한 층 $\widetilde{M_{(f)}}$---여기서 제한 전의 층은
$D_+(f)$ 위에 있다---은 가군 $(M_{(f)})_{g^d/f^e}$에 대응하는 층과 표준적으로
동일시되고(\hyperref[I.1.3.6-ko]{I, 1.3.6}), 따라서
$\widetilde{M_{(fg)}}$와도 동일시된다
(\hyperref[II.2.2.2-ko]{2.2.2}). 그러므로 표준 동형
\[
  \theta_{g,f}:\widetilde{M_{(f)}}|D_+(fg)
  \xrightarrow{\sim}\widetilde{M_{(g)}}|D_+(fg)
\]
이 존재하며, $h$가 $S_+$의 세 번째 동차 원소이면
$\theta_{f,h}=\theta_{f,g}\circ\theta_{g,h}$가 $D_+(fgh)$에서 성립한다.
따라서
(\hyperref[0.3.3.1-ko]{0, 3.3.1}) $X$ 위에 준연접
$\mathcal{O}_X$-가군 $\mathcal{F}$가 존재하고, 각 $f$가 $S_+$의 동차
원소일 때 동형 $\eta_f$가 $\mathcal{F}|D_+(f)$에서
$\widetilde{M_{(f)}}$로 가며 $\theta_{g,f}=\eta_g\circ\eta_f^{-1}$을
만족한다.
이제 층 $\mathcal{G}$를 생각하자. 이는 $X$의 위상의 기저인
$D_+(f)$들 위에서 $D_+(f)\to M_{(f)}$로 정의되고, 표준 준동형
$M_{(f)}\to M_{(fg)}$를 제한 준동형으로 갖는 준층을 층화하여 얻는
층이다. 앞의 논의로부터 $\mathcal{F}$와 $\mathcal{G}$가 동형임을 안다
((\hyperref[I.1.3.7-ko]{I, 1.3.7})도 고려한다). 층 $\mathcal{G}$를
$\widetilde{M}$으로 나타내면 이 층은 명제의 조건들을 만족한다. 특히
$\widetilde{S}=\mathcal{O}_X$이다.
\end{proof}

\begin{definition}[2.5.3]
\phantomsection
\label{II.2.5.3-ko}
(\hyperref[II.2.5.2-ko]{2.5.2})에서 정의한 준연접
$\mathcal{O}_X$-가군 $\widetilde{M}$을 등급 $S$-가군 $M$에
\emph{대응하는} 층이라 한다.
\end{definition}

등급 $S$-가군들은 등급가군의 준동형을 \emph{차수 $0$의} 준동형으로
제한하면 하나의 범주를 이룸을 상기하자. 이 약속 아래에서 다음이
성립한다.

\begin{proposition}[2.5.4]
\phantomsection
\label{II.2.5.4-ko}
함자 $M\mapsto\widetilde{M}$은 등급 $S$-가군의 범주에서 준연접
$\mathcal{O}_X$-가군의 범주로 가는 공변 가법 완전 함자이고,
귀납극한 및 직합과 가환한다.
\end{proposition}

\begin{proof}
실제로 이 성질들은 국소적이므로 층
$\widetilde{M}|D_+(f)=\widetilde{M_{(f)}}$에 대하여 확인하면 충분하다.
그런데 함자 $M\mapsto M_f$, 함자 $N\mapsto N_0$ (등급
$S_f$-가군의 범주에서) 및 함자 $P\mapsto\widetilde{P}$
($S_{(f)}$-가군의 범주에서)는 모두 완전성과 귀납극한 및 직합과의
가환성을 갖는다(\hyperref[I.1.3.5-ko]{I, 1.3.5} 및
\hyperref[I.1.3.9-ko]{I, 1.3.9}). 이로부터 명제가 따른다.
\end{proof}

$\widetilde{u}$로 준동형 $\widetilde{M}\to\widetilde{N}$을 나타내기로
한다. 이는 차수 $0$의 준동형 $u:M\to N$에 대응하는 준동형이다.
(\hyperref[II.2.5.4-ko]{2.5.4})에서 곧바로, 등급 $S$-가군과 차수
$0$의 준동형에 대해서도 (여기서 정한 $\widetilde{M}$의 뜻으로)
(\hyperref[I.1.3.9-ko]{I, 1.3.9} 및
\hyperref[I.1.3.10-ko]{I, 1.3.10})의 결과들이 여전히 성립함을 알 수
있다. 실제로 그 증명들은 순전히 형식적이다.

\begin{proposition}[2.5.5]
\phantomsection
\label{II.2.5.5-ko}
모든 $\mathfrak{p}\in X=\operatorname{Proj}(S)$에 대하여
$\widetilde{M}_{\mathfrak{p}}=M_{(\mathfrak{p})}$이다.
\end{proposition}

\begin{proof}
실제로 정의에 따라
$\widetilde{M}_{\mathfrak{p}}=
\varinjlim\Gamma(D_+(f),\widetilde{M})$이다. 여기서 $f$의 범위는
$S_+$의 동차 원소 가운데 $f\notin\mathfrak{p}$인 원소 전체이다. 한편
$\Gamma(D_+(f),\widetilde{M})=M_{(f)}$이므로, 명제는
$M_{(\mathfrak{p})}$의 정의
(\hyperref[II.2.2.7-ko]{2.2.7})에서 따른다.
\end{proof}
\oldpage[II]{32}

특히 \emph{국소환 $(\mathcal{O}_X)_\mathfrak{p}$}은 바로 환
$S_{(\mathfrak{p})}$이다. 이는 분수 $x/f$들의 집합으로, 여기서 $f$는
$S_+$의 동차 원소이고 $\mathfrak{p}$에 속하지 않으며, $x$는 $f$와
같은 차수의 동차 원소이다.

더욱 구체적으로, $S$가 \emph{본질적으로 정역}이어서
$\operatorname{Proj}(S)=X$가 \emph{정역 스킴}이고
(\hyperref[II.2.4.4-ko]{2.4.4}), $\xi=(0)$가 $X$의 일반점이면,
\emph{유리함수체} $R(X)=\mathcal{O}_\xi$는 $fg^{-1}$ 꼴의 원소들로 이루어진
체이다. 여기서 $f$와 $g$는 $S_+$에 속하는 같은 차수의 동차 원소이고
$g\neq0$이다.

\begin{proposition}[2.5.6]
\phantomsection
\label{II.2.5.6-ko}
모든 $z\in M$에 대하여, $f$가 $S_+$의 동차 원소일 때마다 $f$의 어떤
거듭제곱이 $z$를 소멸시키면, $\widetilde{M}=0$이다. 이 충분조건은
$S_0$-대수 $S$가 집합 $S_1$에 의해 생성되고, 이 집합이 차수 $1$의
동차 원소 전체로 이루어질 때에는 필요조건이기도 하다.
\end{proposition}

\begin{proof}
실제로 조건 $\widetilde{M}=0$은 $M_{(f)}=0$이 $f$가 $S_+$의 동차
원소일 때마다 성립하는 것과 동치이다. 한편 $f\in S_d$일 때
$M_{(f)}=0$이라는 것은 모든 동차 $z\in M$ 가운데 차수가 $d$의 배수인
각 원소에 대하여 어떤 거듭제곱 $f^n$이 존재하여 $f^nz=0$임을 뜻한다.
$M_{(f)}=0$이 모든 $f\in S_1$에 대해 성립하면 명제의 조건은 모든
$f\in S_1$에 대하여 성립한다. 또한 앞의 생성 가정 아래에서는 $f$가
$S_+$의 동차 원소일 때마다 이 조건이 \emph{더 나아가}
성립한다. 이 생성 가정은 $S_1$이 $S$를 생성한다는 것이며, 실제로
$S_+$의 모든 동차 원소는 $S_1$의 원소들의 곱들의 선형결합이기 때문이다.
\end{proof}

\begin{proposition}[2.5.7]
\phantomsection
\label{II.2.5.7-ko}
$d$를 $>0$인 정수라 하고 $f\in S_d$라 하자. 그러면 모든
$n\in\mathbf{Z}$에 대하여 $(\mathcal{O}_X|D_+(f))$-가군
$\widetilde{S(nd)}|D_+(f)$는 $\mathcal{O}_X|D_+(f)$와 표준적으로
동형이다.
\end{proposition}

\begin{proof}
실제로 $(f/1)^n$은 $S_f$의 가역원이며, 이를 곱하는 사상은
$S_{(f)}=(S_f)_0$에서
$(S_f)_{nd}=(S_f(nd))_0=(S(nd)_f)_0=S(nd)_{(f)}$로 가는 전단사이다.
다시 말해 $S_{(f)}$-가군 $S_{(f)}$와 $S(nd)_{(f)}$는 표준적으로
동형이고, 따라서 명제가 성립한다.
\end{proof}

\begin{corollary}[2.5.8]
\phantomsection
\label{II.2.5.8-ko}
열린부분집합
\[
  U=\bigcup_{f\in S_d}D_+(f),
\]
위에서 $\mathcal{O}_X$-가군 $\widetilde{S(nd)}$의 제한은 가역인
$(\mathcal{O}_X|U)$-가군이다
(\hyperref[0.5.4.1-ko]{0, 5.4.1}).
\end{corollary}

\begin{corollary}[2.5.9]
\phantomsection
\label{II.2.5.9-ko}
아이디얼 $S_+$가 $S$에서 집합 $S_1$에 의해 생성되고, 이 집합이 차수
$1$의 동차 원소 전체로 이루어지면, $\mathcal{O}_X$-가군
$\widetilde{S(n)}$은 모든 $n\in\mathbf{Z}$에 대하여 가역이다.
\end{corollary}

\begin{proof}
$X=\bigcup_{f\in S_1}D_+(f)$임은 가정과
(\hyperref[II.2.3.14-ko]{2.3.14})에서 따른다. 이제 $U=X$로 두고
(\hyperref[II.2.5.8-ko]{2.5.8})을 적용하면 충분하다.
\end{proof}

\begin{env}[2.5.10]
\phantomsection
\label{II.2.5.10-ko}
이 절의 나머지 부분 전체에서
\[
\label{II.2.5.10.1-ko}
  \mathcal{O}_X(n)=\widetilde{S(n)}
\tag{2.5.10.1}
\]
로 두기로 한다. 이는 모든 $n\in\mathbf{Z}$에 대한 약속이다. 또한
$U$를 $X$의 임의의 열린부분집합이라 하고, 임의의
$(\mathcal{O}_X|U)$-가군 $\mathcal{F}$에 대하여
\[
\label{II.2.5.10.2-ko}
  \mathcal{F}(n)=
  \mathcal{F}\otimes_{\mathcal{O}_X|U}(\mathcal{O}_X(n)|U)
\tag{2.5.10.2}
\]
로 두기로 한다. 이는 모든 $n\in\mathbf{Z}$에 대한 약속이다. 아이디얼
$S_+$가 $S_1$에 의해 생성되면, $\mathcal{F}(n)$을 대응시키는 함자는
$\mathcal{F}$에 관해 모든 $n\in\mathbf{Z}$에 대하여 \emph{완전}하다. 이 경우
$\mathcal{O}_X(n)$이 가역인 $\mathcal{O}_X$-가군이기 때문이다.
\end{env}

\begin{env}[2.5.11]
\phantomsection
\label{II.2.5.11-ko}
$M$과 $N$을 두 등급 $S$-가군이라 하자. 모든 $f\in S_d$ ($d>0$)에
대하여 다음과 같은 $S_{(f)}$-가군의 함자적인 표준 준동형을 정의한다.
\[
\label{II.2.5.11.1-ko}
  \lambda_f:M_{(f)}\otimes_{S_{(f)}}N_{(f)}
  \longrightarrow(M\otimes_S N)_{(f)}
\tag{2.5.11.1}
\]
\oldpage[II]{33}

정의는 준동형
$M_{(f)}\otimes_{S_{(f)}}N_{(f)}\to M_f\otimes_{S_f}N_f$
---이는 단사 준동형들 $M_{(f)}\to M_f$, $N_{(f)}\to N_f$ 및
$S_{(f)}\to S_f$에서 유도된다---과 표준 동형
$M_f\otimes_{S_f}N_f\xrightarrow{\sim}(M\otimes_S N)_f$
(\hyperref[0.1.3.4-ko]{0, 1.3.4})을 합성하여 얻는다. 또한 두 등급가군의
텐서곱의 정의에 따라 뒤의 동형이 차수를 보존한다는 점을 사용한다. 따라서
$x\in M_{md}$, $y\in N_{nd}$ ($m\geq0$, $n\geq0$)에 대하여
\[
  \lambda_f((x/f^m)\otimes(y/f^n))=(x\otimes y)/f^{m+n}.
\]

이 정의에서 곧바로, $g\in S_e$ ($e>0$)이면 다음 도식이 가환함을 알 수
있다.
\[
  \xymatrix{
    M_{(f)}\otimes_{S_{(f)}}N_{(f)}\ar[r]^{\lambda_f}\ar[d]
    & (M\otimes_S N)_{(f)}\ar[d]\\
    M_{(fg)}\otimes_{S_{(fg)}}N_{(fg)}\ar[r]_{\lambda_{fg}}
    & (M\otimes_S N)_{(fg)}
  }
\]
여기서 오른쪽 세로 화살표는 표준 준동형이고 왼쪽 세로 화살표는 표준
준동형들에서 유도된다. 따라서 이들 $\lambda$로부터
$\mathcal{O}_X$-가군의 함자적인 표준 준동형
\[
\label{II.2.5.11.2-ko}
  \lambda:\widetilde{M}\otimes_{\mathcal{O}_X}\widetilde{N}
  \longrightarrow\widetilde{M\otimes_S N}.
\tag{2.5.11.2}
\]
을 얻는다.

특히 두 동차 아이디얼 $\mathfrak{J}$, $\mathfrak{K}$를 생각하자. 이들은
$S$의 아이디얼이다. $\widetilde{\mathfrak{J}}$와
$\widetilde{\mathfrak{K}}$는 $\mathcal{O}_X$의 아이디얼층이므로 표준 준동형
$\widetilde{\mathfrak{J}}\otimes_{\mathcal{O}_X}
\widetilde{\mathfrak{K}}\to\mathcal{O}_X$가 존재하며, 다음 도식은 가환한다.
\[
\label{II.2.5.11.3-ko}
  \xymatrix{
    \widetilde{\mathfrak{J}}\otimes_{\mathcal{O}_X}
    \widetilde{\mathfrak{K}}\ar[rr]^\lambda\ar[dr]
    && \widetilde{\mathfrak{J}\otimes_S\mathfrak{K}}\ar[dl]\\
    & \mathcal{O}_X &
  }
\tag{2.5.11.3}
\]
실제로 각 열린집합 $D_+(f)$ ($f$는 $S_+$의 동차 원소)에서 이를 확인하는
경우로 귀착할 수 있고, 이는 $\lambda_f$의 정의
(\hyperref[II.2.5.11.1-ko]{2.5.11.1})와
(\hyperref[I.1.3.13-ko]{I, 1.3.13})에서 곧바로 따른다.

끝으로 $M$, $N$, $P$를 세 등급 $S$-가군이라 하면 다음 도식을 얻는다.
\[
\label{II.2.5.11.4-ko}
  \xymatrix{
    \widetilde{M}\otimes_{\mathcal{O}_X}\widetilde{N}
    \otimes_{\mathcal{O}_X}\widetilde{P}
    \ar[r]^{\lambda\otimes1}\ar[d]_{1\otimes\lambda}
    & \widetilde{M\otimes_S N}\otimes_{\mathcal{O}_X}\widetilde{P}
    \ar[d]^\lambda\\
    \widetilde{M}\otimes_{\mathcal{O}_X}\widetilde{N\otimes_S P}
    \ar[r]_\lambda
    & \widetilde{M\otimes_S N\otimes_S P}
  }
\tag{2.5.11.4}
\]
\oldpage[II]{34}

이 도식은 가환한다. 다시 각 $D_+(f)$에서 확인하면 충분하며, 이는 정의와
(\hyperref[I.1.3.13-ko]{I, 1.3.13})에서 곧바로 따른다.
\end{env}

\begin{env}[2.5.12]
\phantomsection
\label{II.2.5.12-ko}
(\hyperref[II.2.5.11-ko]{2.5.11})의 가정 아래에서 다음과 같은
$S_{(f)}$-가군의 함자적인 표준 준동형을 정의한다.
\[
\label{II.2.5.12.1-ko}
  \mu_f:\bigl(\operatorname{Hom}_S(M,N)\bigr)_{(f)}
  \longrightarrow
  \operatorname{Hom}_{S_{(f)}}(M_{(f)},N_{(f)})
\tag{2.5.12.1}
\]
이는 $u/f^n$---여기서 $u$는 차수 $nd$인 준동형이다---에 준동형
$M_{(f)}\to N_{(f)}$을 대응시켜 정의한다. 이 준동형은 $x/f^m$
($x\in M_{md}$)을 $u(x)/f^{m+n}$으로 보낸다. $g\in S_e$ ($e>0$)에
대해서도 다음 도식은 가환한다.
\[
  \xymatrix{
    \bigl(\operatorname{Hom}_S(M,N)\bigr)_{(f)}
    \ar[r]^{\mu_f}\ar[d]
    & \operatorname{Hom}_{S_{(f)}}(M_{(f)},N_{(f)})\ar[d]\\
    \bigl(\operatorname{Hom}_S(M,N)\bigr)_{(fg)}
    \ar[r]_{\mu_{fg}}
    & \operatorname{Hom}_{S_{(fg)}}(M_{(fg)},N_{(fg)})
  }
\]
여기서 왼쪽 화살표는 표준 준동형이고 오른쪽 화살표는 표준 준동형들에서
유도된다. 따라서 (\hyperref[I.1.3.8-ko]{I, 1.3.8})을 고려하면
$\mu_f$들이 다음과 같은 $\mathcal{O}_X$-가군의 함자적인 표준 준동형을
정의함을 다시 얻는다.
\[
\label{II.2.5.12.2-ko}
  \mu:\widetilde{\operatorname{Hom}_S(M,N)}
  \longrightarrow
  \mathcal{H}om_{\mathcal{O}_X}(\widetilde{M},\widetilde{N}).
\tag{2.5.12.2}
\]
\end{env}

\begin{proposition}[2.5.13]
\phantomsection
\label{II.2.5.13-ko}
아이디얼 $S_+$가 $S_1$에 의해 생성된다고 가정하자. 그러면 준동형
$\lambda$(\hyperref[II.2.5.11.2-ko]{2.5.11.2})는 동형사상이고, 준동형
$\mu$(\hyperref[II.2.5.12.2-ko]{2.5.12.2})도 등급 $S$-가군 $M$이
유한 표시를 가지는 경우(\hyperref[II.2.1.1-ko]{2.1.1})에는 동형사상이다.
\end{proposition}

\begin{proof}
$X$가 $D_+(f)$ ($f\in S_1$)들의 합집합이므로
(\hyperref[II.2.3.14-ko]{2.3.14}), 고려하는 각각의 가정 아래에서
$\lambda_f$와 $\mu_f$가 $f$가 동차이고 \emph{차수가 $1$}일 때
동형사상임을 보이는 것으로 귀착된다. 이때 $\mathbf{Z}$-쌍선형 사상
\[
  M_m\times N_n\longrightarrow
  M_{(f)}\otimes_{S_{(f)}}N_{(f)}
\]
을 $(x,y)$에 원소 $(x/f^m)\otimes(y/f^n)$을 대응시켜 정의한다
($m<0$이면 $x/f^m$은 $f^{-m}x/1$을 뜻한다). 이 사상들로
$\mathbf{Z}$-선형 사상
\[
  M\otimes_{\mathbf{Z}}N\longrightarrow
  M_{(f)}\otimes_{S_{(f)}}N_{(f)},
\]
을 정의할 수 있다. $s\in S_q$이면 이 사상은 $(sx)\otimes y$를
\[
  (s/f^q)\bigl((x/f^m)\otimes(y/f^n)\bigr)
\]
으로 보낸다($x\in M_m$, $y\in N_n$). 따라서 가군의 쌍준동형
\[
  \gamma_f:M\otimes_S N\longrightarrow
  M_{(f)}\otimes_{S_{(f)}}N_{(f)},
\]
을 얻는다. 이는 표준 준동형 $S\to S_{(f)}$---$s\in S_q$를
$s/f^q$로 보내는 준동형---에 관한 것이다. 더 나아가 원소
$\sum_i(x_i\otimes y_i)$가 $M\otimes_S N$에 속하고, 여기서
$x_i$, $y_i$는 각각 차수가
$m_i$, $n_i$인 동차 원소이다---에 대하여
\[
  f^r\sum_i(x_i\otimes y_i)=0,
\]
즉 $\sum_i(f^rx_i\otimes y_i)=0$이라고 가정하자.
(\hyperref[0.1.3.4-ko]{0, 1.3.4})에 따라
\[
  \sum_i(f^rx_i/f^{m_i+r})\otimes(y_i/f^{n_i})=0,
\]
을 얻는다. 다시 말해 $\gamma_f(\sum_i(x_i\otimes y_i))=0$이다. 따라서
$\gamma_f$는 다음과 같이 인수분해된다.
\[
  M\otimes_S N\longrightarrow(M\otimes_S N)_f
  \xrightarrow{\gamma'_f}
  M_{(f)}\otimes_{S_{(f)}}N_{(f)}~;
\]
$\lambda'_f$를 $\gamma'_f$의
\oldpage[II]{35}
$(M\otimes_S N)_{(f)}$로의 제한이라 하면,
$\lambda_f$와 $\lambda'_f$가 서로 역인 $S_{(f)}$-준동형임을 곧바로
확인할 수 있다. 이로써 명제의 첫 번째 부분이 증명된다.

두 번째 부분을 증명하기 위해, $M$이 준동형 $P\to Q$의 여핵이라고
가정하자. 이는 등급 $S$-가군의 준동형이며, $P$와 $Q$는 $S(n)$ 꼴인
가군들의 유한 직합이다. $\operatorname{Hom}_S(L,N)$의 $L$에 관한
왼쪽 완전성과 $M_{(f)}$의 $M$에 관한 완전성을 사용하면, $\mu_f$가
$M=S(n)$일 때 동형사상임을 보이는 것으로 곧바로 귀착된다. 이제 임의의
동차 원소 $z$를 $N$에서 택하고, $u_z$를 $S(n)$에서 $N$으로 가며
$u_z(1)=z$를 만족하는 준동형이라 하자. 그러면
$\eta:z\mapsto u_z$는 차수 $0$이고 $N(-n)$에서
$\operatorname{Hom}_S(S(n),N)$으로 가는 동형사상임을 곧바로
알 수 있다. 이에 대응하는 동형사상
\[
  \eta_f:\bigl(N(-n)\bigr)_{(f)}
  \longrightarrow
  \bigl(\operatorname{Hom}_S(S(n),N)\bigr)_{(f)}.
\]
이 있다. 한편 $\eta'_f$를 다음 동형사상이라 하자.
\[
  N_{(f)}\longrightarrow
  \operatorname{Hom}_{S_{(f)}}\bigl(S(n)_{(f)},N_{(f)}\bigr)
\]
이는 각 $z'\in N_{(f)}$에 준동형 $v_{z'}$를 대응시키며, 이 준동형은
$v_{z'}(s/f^k)=sz'/f^{n+k}$를 만족한다
($s\in S_{n+k}=(S(n))_k$). 합성 사상
\[
  \bigl(N(-n)\bigr)_{(f)}
  \xrightarrow{\eta_f}
  \bigl(\operatorname{Hom}_S(S(n),N)\bigr)_{(f)}
  \xrightarrow{\mu_f}
  \operatorname{Hom}_{S_{(f)}}\bigl(S(n)_{(f)},N_{(f)}\bigr)
  \xrightarrow{{\eta'_f}^{-1}}N_{(f)}
\]
은 $z/f^h\mapsto z/f^{h-n}$으로 주어지는
$\bigl(N(-n)\bigr)_{(f)}$에서 $N_{(f)}$로 가는 동형사상임을 쉽게 확인할 수
있다. 따라서 $\mu_f$는 동형사상이다.
\end{proof}

아이디얼 $S_+$가 $S_1$에 의해 생성되면,
(\hyperref[II.2.5.13-ko]{2.5.13})에서 $S$의 모든 등급 아이디얼
$\mathfrak{J}$와 모든 등급 $S$-가군 $M$에 대하여
\[
\label{II.2.5.13.1-ko}
  \widetilde{\mathfrak{J}}\cdot\widetilde{M}
  =\widetilde{\mathfrak{J}\cdot M}
\tag{2.5.13.1}
\]
가 표준 동형까지 성립함을 얻는다. 실제로 이는 다음 도식의 가환성에서
따른다.
\[
  \xymatrix{
    \widetilde{\mathfrak{J}}\otimes_{\mathcal{O}_X}\widetilde{M}
    \ar[rr]^\lambda\ar[dr]
    && \widetilde{\mathfrak{J}\otimes_S M}\ar[dl]\\
    & \widetilde{M} &
  }
\]
이 가환성은 (\hyperref[II.2.5.11.3-ko]{2.5.11.3})의 경우와 마찬가지로
확인한다.

\begin{corollary}[2.5.14]
\phantomsection
\label{II.2.5.14-ko}
$S$가 $S_1$에 의해 생성된다고 하자. 어떠한 $m,n\in\mathbf{Z}$에
대해서도 다음 공식들이 성립한다.
\[
\label{II.2.5.14.1-ko}
  \mathcal{O}_X(m)\otimes_{\mathcal{O}_X}\mathcal{O}_X(n)
  =\mathcal{O}_X(m+n)
\tag{2.5.14.1}
\]
\[
\label{II.2.5.14.2-ko}
  \mathcal{O}_X(n)=\bigl(\mathcal{O}_X(1)\bigr)^{\otimes n}
\tag{2.5.14.2}
\]
여기서 등식들은 표준 동형까지 성립한다.
\end{corollary}

\begin{proof}
첫째 공식은 (\hyperref[II.2.5.13-ko]{2.5.13})과 차수 $0$인 표준 동형
$S(m)\otimes_S S(n)\xrightarrow{\sim}S(m+n)$의 존재에서 따른다. 이 동형은
$1\otimes1$---여기서 첫째 인자 $1$은 $(S(m))_{-m}$에, 둘째 인자 $1$은
$(S(n))_{-n}$에 속한다---을 원소 $1\in(S(m+n))_{-m-n}$에 대응시킨다.
이제 둘째 공식은 $n=-1$인 경우만 증명하면 충분하다. 그리고
(\hyperref[II.2.5.13-ko]{2.5.13})에 의하면 이는
$\operatorname{Hom}_S(S(1),S)$가 $S(-1)$과 표준적으로 동형임을 보이는
것으로 귀착한다. 정의들 (\hyperref[II.2.1.2-ko]{2.1.2})을 되짚어 보고
$S(1)$이 한 원소로 생성되는 $S$-가군임을 상기하면 곧바로 확인된다.
\end{proof}

\oldpage[II]{36}
\begin{corollary}[2.5.15]
\phantomsection
\label{II.2.5.15-ko}
$S$가 $S_1$에 의해 생성된다고 하자. 모든 등급 $S$-가군 $M$과 모든
$n\in\mathbf{Z}$에 대하여
\[
\label{II.2.5.15.1-ko}
  \widetilde{M(n)}=\widetilde{M}(n)
\tag{2.5.15.1}
\]
가 표준 동형까지 성립한다.
\end{corollary}

\begin{proof}
이는 정의들 (\hyperref[II.2.5.10.2-ko]{2.5.10.2})와
(\hyperref[II.2.5.10.1-ko]{2.5.10.1}), 명제
(\hyperref[II.2.5.13-ko]{2.5.13}), 그리고 차수 $0$인 표준 동형
$M(n)\xrightarrow{\sim}M\otimes_S S(n)$의 존재에서 따른다. 이 동형은 모든
$z\in(M(n))_h=M_{n+h}$를
$z\otimes1\in M_{n+h}\otimes(S(n))_{-n}\subset(M\otimes_S S(n))_h$에
대응시킨다.
\end{proof}

\begin{env}[2.5.16]
\phantomsection
\label{II.2.5.16-ko}
$S'_0=\mathbf{Z}$이고 $n>0$에 대하여 $S'_n=S_n$인 등급환을 $S'$라 하자.
그러면 $f\in S_d$ ($d>0$)일 때 모든 $n\in\mathbf{Z}$에 대하여
$(S(n))_{(f)}=(S'(n))_{(f)}$이다. 실제로 $(S'(n))_{(f)}$의 원소는
$x\in S'_{n+kd}$ ($k>0$)인 $x/f^k$ 꼴이며, 항상 $n+kd\neq0$이 되도록
$k$를 택할 수 있다. $X=\operatorname{Proj}(S)$와
$X'=\operatorname{Proj}(S')$가 표준적으로 동일시되므로
(\hyperref[II.2.4.7-ko]{2.4.7, (ii)}), 모든 $n\in\mathbf{Z}$에 대하여
$\mathcal{O}_X(n)$과 $\mathcal{O}_{X'}(n)$은 앞의 동일시에서 서로
대응한다.

한편 모든 $d>0$과 모든 $n\in\mathbf{Z}$에 대하여
\[
  (S^{(d)}(n))_h=S_{(n+h)d}=(S(nd))_{hd}
\]
이며, $f\in S_d$이면 따라서
$(S^{(d)}(n))_{(f)}=(S(nd))_{(f)}$이다. 스킴
$X=\operatorname{Proj}(S)$와 $X^{(d)}=\operatorname{Proj}(S^{(d)})$가
표준적으로 동일시됨을 알고 있다
(\hyperref[II.2.4.7-ko]{2.4.7, (i)}). 앞의 논의에 따르면 $S_0$-대수
$S^{(d)}$가 $S_d$에 의해 생성될 때, 모든 $n\in\mathbf{Z}$에 대하여
$\mathcal{O}_X(nd)$와 $\mathcal{O}_{X^{(d)}}(n)$은 이 동일시에서 서로
대응한다.
\end{env}

\begin{proposition}[2.5.17]
\phantomsection
\label{II.2.5.17-ko}
$d$를 $>0$인 정수라 하고
$U=\bigcup_{f\in S_d}D_+(f)$라 하자. 표준 준동형
$\mathcal{O}_X(nd)\otimes_{\mathcal{O}_X}\mathcal{O}_X(-nd)
\to\mathcal{O}_X$의 $U$에 대한 제한은 모든 정수 $n$에 대하여
동형이다.
\end{proposition}

\begin{proof}
(\hyperref[II.2.5.16-ko]{2.5.16})에 의해 $d=1$인 경우로 환원할 수 있고,
그 경우 결론은 (\hyperref[II.2.5.13-ko]{2.5.13})의 증명에서 따른다.
\end{proof}

\subsection{$\operatorname{Proj}(S)$ 위의 층에 대응하는 등급
$S$-가군.}
\label{subsection:II.2.6-ko}

\emph{이 절 전체에서 $S$의 아이디얼 $S_+$가 차수 $1$의 동차 원소들의
집합 $S_1$에 의해 생성된다고 가정한다.}

\begin{env}[2.6.1]
\phantomsection
\label{II.2.6.1-ko}
그러면 $\mathcal{O}_X$-가군 $\mathcal{O}_X(1)$은
\emph{가역}이다(\hyperref[II.2.5.9-ko]{2.5.9}). 이때 모든
$\mathcal{O}_X$-가군 $\mathcal{F}$
(\hyperref[0.5.4.6-ko]{0, 5.4.6})에 대하여
\[
\label{II.2.6.1.1-ko}
  \Gamma_*(\mathcal{F})
  =\Gamma_*(\mathcal{O}_X(1),\mathcal{F})
  =\bigoplus_{n\in\mathbf{Z}}\Gamma(X,\mathcal{F}(n))
\tag{2.6.1.1}
\]
로 둔다. 여기서는 (\hyperref[II.2.5.14.2-ko]{2.5.14.2})를 고려하였다.
(\hyperref[0.5.4.6-ko]{0, 5.4.6})에서 상기한 바와 같이
$\Gamma_*(\mathcal{O}_X)$는 \emph{등급환}의 구조를 가지며,
$\Gamma_*(\mathcal{F})$는 $\Gamma_*(\mathcal{O}_X)$ 위의
\emph{등급가군} 구조를 가진다.

$\mathcal{O}_X(n)$이 국소 자유이므로
$\mathcal{F}\mapsto\Gamma_*(\mathcal{F})$는 $\mathcal{F}$에 관한
공변 가법 \emph{왼쪽 완전} 함자이다. 특히 $\mathcal{J}$가
$\mathcal{O}_X$의 아이디얼층이면, $\Gamma_*(\mathcal{J})$는
$\Gamma_*(\mathcal{O}_X)$의 \emph{등급 아이디얼}과 표준적으로
동일시된다.
\end{env}

\begin{env}[2.6.2]
\phantomsection
\label{II.2.6.2-ko}
$M$을 등급 $S$-가군이라 하자. 모든 $f\in S_d$ ($d>0$)에 대하여
$x\mapsto x/1$은 아벨군의 준동형 $M_0\to M_{(f)}$이고,
$M_{(f)}$는 표준적으로

\oldpage[II]{37}

$\Gamma(D_+(f),\widetilde{M})$와 동일시되므로 아벨군의 준동형
\[
  \alpha_0^f:M_0\longrightarrow\Gamma(D_+(f),\widetilde{M}).
\]
을 얻는다. 모든 $g\in S_e$ ($e>0$)에 대하여 다음 도식이 가환함은
명백하다.
\[
  \xymatrix{
    & \Gamma(D_+(f),\widetilde{M})\ar[dd]\\
    M_0\ar[ur]^{\alpha_0^f}\ar[dr]_{\alpha_0^{fg}}\\
    & \Gamma(D_+(fg),\widetilde{M})
  }
\]
이는 모든 $x\in M_0$에 대하여 $\widetilde{M}$의 단면
$\alpha_0^f(x)$와 $\alpha_0^g(x)$가 $D_+(f)\cap D_+(g)$에서
일치함을 뜻한다. 따라서 각 $D_+(f)$로의 제한이 $\alpha_0^f(x)$인
유일한 단면 $\alpha_0(x)\in\Gamma(X,\widetilde{M})$이 존재한다.
이로써 ($S$가 $S_1$에 의해 생성된다는 가정을 사용하지 않고) 아벨군의
준동형
\[
\label{II.2.6.2.1-ko}
  \alpha_0:M_0\longrightarrow\Gamma(X,\widetilde{M}).
\tag{2.6.2.1}
\]
을 정의하였다.

이 결과를 모든 $n\in\mathbf{Z}$에 대하여 등급 $S$-가군 $M(n)$에
적용하면, 각 $n\in\mathbf{Z}$에 대하여 아벨군의 준동형
\[
\label{II.2.6.2.2-ko}
  \alpha_n:M_n=(M(n))_0\longrightarrow\Gamma(X,\widetilde{M}(n))
\tag{2.6.2.2}
\]
을 얻는다((\hyperref[II.2.5.15-ko]{2.5.15})를 고려한다). 따라서 등급
아벨군의 함자적인 준동형(차수 $0$)
\[
\label{II.2.6.2.3-ko}
  \alpha:M\longrightarrow\Gamma_*(\widetilde{M})
\tag{2.6.2.3}
\]
을 얻으며, 이를 $\alpha_M$으로도 나타낸다. 이 준동형은 각 $M_n$에서
$\alpha_n$과 일치한다.

특히 $M=S$로 두면, $\Gamma_*(\mathcal{O}_X)$의 곱셈의 정의
(\hyperref[0.5.4.6-ko]{0, 5.4.6})를 고려하여
$\alpha:S\to\Gamma_*(\mathcal{O}_X)$가 등급환의 준동형임을 곧바로
확인할 수 있다. 또한 모든 등급 $S$-가군 $M$에 대하여
(\hyperref[II.2.6.2.3-ko]{2.6.2.3})은 등급 가군의 쌍준동형이다.
\end{env}

\begin{proposition}[2.6.3]
\phantomsection
\label{II.2.6.3-ko}
모든 $f\in S_d$ ($d>0$)에 대하여, $D_+(f)$는 $\mathcal{O}_X(d)$의
단면 $\alpha_d(f)$가 소멸하지 않는 $\mathfrak{p}\in X$들의 집합과
일치한다(\hyperref[0.5.5.2-ko]{0, 5.5.2}).
\end{proposition}

\begin{proof}
가정에 따라 $X=\bigcup_{g\in S_1}D_+(g)$이므로, 모든 $g\in S_1$에
대하여 $\alpha_d(f)$가 소멸하지 않는 $\mathfrak{p}\in D_+(g)$들의
집합이 $D_+(fg)$와 일치함을 보이면 충분하다. 그런데 $\alpha_d(f)$의
$D_+(g)$에 대한 제한은 정의상 $(S(d))_{(g)}$의 원소 $f/1$에 대응하는
단면이다. 표준 동형
$(S(d))_{(g)}\xrightarrow{\sim}S_{(g)}$
(\hyperref[II.2.5.7-ko]{2.5.7}) 아래에서, $D_+(g)$ 위의
$\mathcal{O}_X(d)$의 이 단면은 $S_{(g)}$의 원소 $f/g^d$에 대응하는
$D_+(g)$ 위의 $\mathcal{O}_X$의 단면과 동일시된다. 이 단면이
$\mathfrak{p}\in D_+(g)$에서 소멸한다는 것은 $f/g^d\in\mathfrak{q}$임을
뜻한다. 여기서 $\mathfrak{q}$는 $\mathfrak{p}$에 대응하는 $S_{(g)}$의
소아이디얼이다(\hyperref[II.2.3.6-ko]{2.3.6}). 정의상 이는
$f\in\mathfrak{p}$라는 뜻이므로 명제가 성립한다.
\end{proof}

\begin{env}[2.6.4]
\phantomsection
\label{II.2.6.4-ko}
이제 $\mathcal{F}$를 $\mathcal{O}_X$-가군이라 하고
$M=\Gamma_*(\mathcal{F})$로 두자. 등급환 준동형
$\alpha:S\to\Gamma_*(\mathcal{O}_X)$가 존재하므로, $M$을 등급
$S$-가군으로 볼 수 있다. 모든 $f\in S_d$ ($d>0$)에 대하여
(\hyperref[II.2.6.3-ko]{2.6.3})에서 $\mathcal{O}_X(d)$의 단면
$\alpha_d(f)$의 $D_+(f)$에 대한 제한이 가역임을 알 수 있다. 따라서
모든 $n>0$에 대하여 $\mathcal{O}_X(nd)$의 단면
$\alpha_{nd}(f^n)$\footnote{\textbf{번역 주.} 인쇄 원문은 이 문단의 세
곳에서 각각 $\alpha_d(f^n)$, $\alpha_d(f^k)$, $\alpha_d(f^n)$을
쓰지만, $f^n\in S_{nd}$와 $f^k\in S_{kd}$이므로
EGA2-CANON-DISP-0002에 따라 각각 $\alpha_{nd}(f^n)$,
$\alpha_{kd}(f^k)$, $\alpha_{nd}(f^n)$으로 교정했다.}의 $D_+(f)$에
대한 제한도 가역이다. 이제
$z\in M_{nd}=\Gamma(X,\mathcal{F}(nd))$ ($n>0$)라 하자. 어떤 정수
$k\geq0$에 대하여 $f^kz$의 $D_+(f)$에 대한 제한, 즉
\oldpage[II]{38}
단면
$(z|D_+(f))(\alpha_{kd}(f^k)|D_+(f))$가
$\mathcal{F}((n+k)d)$의 단면으로서 영이면, 앞의 설명에 따라
$z|D_+(f)=0$이기도 하다. 이로써 원소 $z/f^n$에 $D_+(f)$ 위의
$\mathcal{F}$의 단면
$(z|D_+(f))(\alpha_{nd}(f^n)|D_+(f))^{-1}$을 대응시켜
$S_{(f)}$-가군 준동형
$\beta_f:M_{(f)}\to\Gamma(D_+(f),\mathcal{F})$를 정의할 수 있다.
또한 $g\in S_e$ ($e>0$)일 때 다음 도식이 가환함을 곧바로 확인할 수
있다.
\[
\label{II.2.6.4.1-ko}
  \xymatrix{
    M_{(f)}\ar[r]^{\beta_f}\ar[d]
    & \Gamma(D_+(f),\mathcal{F})\ar[d]\\
    M_{(fg)}\ar[r]_{\beta_{fg}}
    & \Gamma(D_+(fg),\mathcal{F})
  }
\tag{2.6.4.1}
\]
$M_{(f)}$가 $\Gamma(D_+(f),\widetilde{M})$과 표준적으로 동일시되고
$D_+(f)$들이 $X$의 위상의 기저를 이룬다는 사실
(\hyperref[II.2.3.4-ko]{2.3.4})을 상기하면, $\beta_f$들이 유일한 표준
$\mathcal{O}_X$-가군 준동형
\[
\label{II.2.6.4.2-ko}
  \beta:\widetilde{\Gamma_*(\mathcal{F})}
  \longrightarrow\mathcal{F}
\tag{2.6.4.2}
\]
에서 유도됨을 알 수 있다. 이 준동형은 $\beta_{\mathcal{F}}$로도
표기하며 명백히 함자적이다.
\end{env}

\begin{proposition}[2.6.5]
\phantomsection
\label{II.2.6.5-ko}
$M$을 등급 $S$-가군, $\mathcal{F}$를 $\mathcal{O}_X$-가군이라 하자.
그러면 다음 합성 준동형들은 각각 항등 준동형이다.
\[
\label{II.2.6.5.1-ko}
  \widetilde{M}\xrightarrow{\widetilde{\alpha}}
  \widetilde{\Gamma_*(\widetilde{M})}
  \xrightarrow{\beta}\widetilde{M}
\tag{2.6.5.1}
\]
\[
\label{II.2.6.5.2-ko}
  \Gamma_*(\mathcal{F})\xrightarrow{\alpha}
  \Gamma_*(\widetilde{\Gamma_*(\mathcal{F})})
  \xrightarrow{\Gamma_*(\beta)}\Gamma_*(\mathcal{F})
\tag{2.6.5.2}
\]
\end{proposition}

\begin{proof}
(\hyperref[II.2.6.5.1-ko]{2.6.5.1})의 확인은 국소적이다. 열린집합
$D_+(f)$에서는 정의들과, 준연접 층에 적용한 $\beta$가 $D_+(f)$ 위의
단면들에 대한 작용으로 정해진다는 사실
(\hyperref[I.1.3.8-ko]{I, 1.3.8})에서 곧바로 따른다.
(\hyperref[II.2.6.5.2-ko]{2.6.5.2})의 확인은 각 차수에서 따로 이루어진다.
$M=\Gamma_*(\mathcal{F})$로 두면
$M_n=\Gamma(X,\mathcal{F}(n))$이고
$(\Gamma_*(\widetilde{M}))_n=\Gamma(X,\widetilde{M}(n))
=\Gamma(X,\widetilde{M(n)})$이다. 그런데 $f\in S_1$이고 $z\in M_n$이면,
$\alpha_n^f(z)$는 $(M(n))_{(f)}$의 원소 $z/1$이며
$(f/1)^n(z/f^n)$과 같다. 이에 $\beta_f$로 대응하는 것은 $D_+(f)$ 위의
단면
\[
  ((\alpha_1(f))^n|D_+(f))
  ((z|D_+(f))((\alpha_1(f))^n|D_+(f))^{-1})
\]
이다. 이는 곧 $z$의 $D_+(f)$에 대한 제한이므로
(\hyperref[II.2.6.5.2-ko]{2.6.5.2})가 확인된다.
\end{proof}

\subsection{유한성 조건.}
\label{subsection:II.2.7-ko}

\begin{proposition}[2.7.1]
\phantomsection
\label{II.2.7.1-ko}
\begin{enumerate}
  \item[\rm(i)] $S$가 뇌터인 등급환이면
    $X=\operatorname{Proj}(S)$는 뇌터 스킴이다.
  \item[\rm(ii)] $S$가 유한형 등급 $A$-대수이면
    $X=\operatorname{Proj}(S)$는 $Y=\operatorname{Spec}(A)$ 위의
    유한형 스킴이다.
\end{enumerate}
\end{proposition}

\oldpage[II]{39}
\begin{proof}
\begin{enumerate}
  \item[\rm(i)] $S$가 뇌터이면 아이디얼 $S_+$는 동차 원소들
    $f_i$ ($1\leq i\leq p$)로 이루어진 유한 생성계를 갖는다. 따라서
    (\hyperref[II.2.3.14-ko]{2.3.14}) 바탕공간 $X$는
    $D_+(f_i)=\operatorname{Spec}(S_{(f_i)})$들의 합집합이고, 각
    $S_{(f_i)}$가 뇌터임을 보이면 충분하다. 이는
    (\hyperref[II.2.2.6-ko]{2.2.6})에서 따른다.
  \item[\rm(ii)] 가정으로부터 $S_0$가 유한형 $A$-대수이고 $S$가
    유한형 $S_0$-대수임이 따른다. 따라서 $S_+$는 유한 생성
    아이디얼이다(\hyperref[II.2.1.4-ko]{2.1.4}). 그러면 (i)에서와 같이
    $f\in S_d$일 때 $S_{(f)}$가 유한형 $A$-대수임을 증명하면
    충분하다. (\hyperref[II.2.2.5-ko]{2.2.5})에 따라 $S^{(d)}$가 유한형
    $A$-대수임을 보이면 충분하고, 이는
    (\hyperref[II.2.1.6-ko]{2.1.6})에서 따른다.
\end{enumerate}
\end{proof}

\begin{env}[2.7.2]
\phantomsection
\label{II.2.7.2-ko}
이하에서는 등급 $S$-가군 $M$에 관한 다음 유한성 조건들을 고려하겠다.
\begin{enumerate}
  \item[\rm(TF)] 어떤 정수 $n$이 존재하여 부분가군
    $\bigoplus_{k\geq n}M_k$가 유한 생성 $S$-가군이다.
  \item[\rm(TN)] 어떤 정수 $n$이 존재하여 $k\geq n$이면 $M_k=0$이다.
\end{enumerate}

$M$이 (TN)을 만족하면 $S_+$의 모든 동차 원소 $f$에 대하여
$M_{(f)}=0$이고, 따라서 $\widetilde{M}=0$이다.

$M$, $N$을 두 등급 $S$-가군이라 하자. 차수 $0$의 준동형
$u:M\to N$에 대하여, 어떤 정수 $n$이 존재하여 $k\geq n$일 때
$u_k:M_k\to N_k$가 단사(각각 전사, 전단사)이면 $u$를
(TN)-\emph{단사}(각각 (TN)-\emph{전사}, (TN)-\emph{전단사})라고 한다.
따라서 $u$가 (TN)-단사(각각 (TN)-전사)라는 것은
$\operatorname{Ker}u$(각각 $\operatorname{Coker}u$)가 (TN)을
만족한다는 것과 같다. (\hyperref[II.2.5.4-ko]{2.5.4})에 따라 $u$가
(TN)-단사(각각 (TN)-전사, (TN)-전단사)이면 $\widetilde{u}$는
단사(각각 전사, 전단사)이다. $u$가 (TN)-전단사일 때 $u$를
(TN)-\emph{동형사상}이라고도 한다.
\end{env}

\begin{proposition}[2.7.3]
\phantomsection
\label{II.2.7.3-ko}
$S$를 아이디얼 $S_+$가 유한 생성인 등급환이라 하고, $M$을 등급
$S$-가군이라 하자.
\begin{enumerate}
  \item[\rm(i)] $M$이 조건 (TF)를 만족하면 $\mathcal{O}_X$-가군
    $\widetilde{M}$은 유한 생성이다.
  \item[\rm(ii)] $M$이 (TF)를 만족한다고 가정하자. $\widetilde{M}=0$일
    필요충분조건은 $M$이 (TN)을 만족하는 것이다.
\end{enumerate}
\end{proposition}

\begin{proof}
방금 본 바와 같이 조건 (TN)은 $\widetilde{M}=0$을 함의한다. $M$이
(TF)를 만족하면, 가정에 따라 유한 생성인 등급 부분가군
$M'=\bigoplus_{k\geq n}M_k$에 대하여 $M/M'$은 (TN)을 만족한다. 따라서
$\widetilde{M/M'}=0$이고, 함자 $M\mapsto\widetilde{M}$의 완전성
(\hyperref[II.2.5.4-ko]{2.5.4})으로부터
$\widetilde{M}=\widetilde{M'}$을 얻는다. 그러므로 $\widetilde{M}$이 유한
생성임을 증명하기 위해 $M$이 유한 생성인 경우로 귀착할 수 있다. 문제는
국소적이므로 $M_{(f)}$가 유한 생성 $S_{(f)}$-가군임을 보이면 충분하다
(\hyperref[I.1.3.9-ko]{I, 1.3.9}). 그런데 $M^{(d)}$는 유한 생성
$S^{(d)}$-가군이고(\hyperref[II.2.1.6-ko]{2.1.6, (iii)}), 이제 주장은
(\hyperref[II.2.2.5-ko]{2.2.5})에서 따른다.

이제 $M$이 (TF)를 만족하고 $\widetilde{M}=0$이라고 가정하자. 위와 같은
표기에서 $\widetilde{M'}=0$이고, $M'$에 대한 조건 (TN)은 $M$에 대한
조건 (TN)과 동치이다. 따라서 $\widetilde{M}=0$이 $M$의 조건 (TN)을
함의함을 증명하기 위해, 다시 $M$이 유한 개의 동차 원소 $x_i$
($1\leq i\leq p$)에 의해 생성되는 경우만 고려하면 된다. 또
$(f_j)_{1\leq j\leq q}$를 아이디얼 $S_+$의 동차 생성계라 하자. 가정에
따라 모든 $j$에 대하여 $M_{(f_j)}=0$이므로, 어떤 정수 $n$이 존재하여
모든 $i$, $j$에 대하여 $f_j^nx_i=0$이다. $n_j$를 $f_j$의 차수라 하고,
$\sum_jr_j\leq nq$를 만족하는 정수족 $(r_j)$에 대하여
$\sum_jr_jn_j$가 취하는 최댓값을 $m$이라 하자. 그러면

\oldpage[II]{40}
$k>m$이면 모든 $i$에 대하여 $S_kx_i=0$임이 명백하다. $x_i$들의 차수 중
최댓값을 $h$라 하면 $k>h+m$일 때 $M_k=0$이고, 이것으로 증명이 끝난다.
\end{proof}

\begin{corollary}[2.7.4]
\phantomsection
\label{II.2.7.4-ko}
$S$를 아이디얼 $S_+$가 유한 생성인 등급환이라 하자. 그러면
$X=\operatorname{Proj}(S)=\varnothing$이기 위한 필요충분조건은 어떤 $n$이
존재하여 $k\geq n$일 때 $S_k=0$인 것이다.
\end{corollary}

\begin{proof}
실제로 조건 $X=\varnothing$은
$\mathcal{O}_X=\widetilde{S}=0$과 동치이고, $S$는 한 원소로 생성되는
$S$-가군이다.
\end{proof}

\begin{theorem}[2.7.5]
\phantomsection
\label{II.2.7.5-ko}
아이디얼 $S_+$가 차수 $1$인 유한 개의 동차 원소로 생성된다고 가정하고,
$X=\operatorname{Proj}(S)$라 하자. 그러면 모든 준연접
$\mathcal{O}_X$-가군 $\mathcal{F}$에 대하여 표준 준동형
$\beta:\widetilde{\Gamma_*(\mathcal{F})}
\to\mathcal{F}$ (\hyperref[II.2.6.4-ko]{2.6.4})는 동형이다.
\end{theorem}

\begin{proof}
실제로 $S_+$가 유한 개의 원소 $f_i\in S_1$로 생성되면, $X$는 준콤팩트인
부분공간 $\operatorname{Spec}(S_{(f_i)})$
(\hyperref[II.2.3.6-ko]{2.3.6})들의 합집합이므로 준콤팩트이다. 또한 $X$는
스킴이다(\hyperref[II.2.4.2-ko]{2.4.2}).
(\hyperref[I.9.3.2-ko]{I, 9.3.2}),
(\hyperref[II.2.5.14.2-ko]{2.5.14.2}) 및
(\hyperref[II.2.6.3-ko]{2.6.3})에 따라 모든 $f\in S_d$ ($d>0$)에 대하여
표준 동형
$(\Gamma_*(\mathcal{F}))_{(\alpha_d(f))}\xrightarrow{\sim}
\Gamma(D_+(f),\mathcal{F})$을 얻는다. 더욱이
$(\Gamma_*(\mathcal{F}))_{(\alpha_d(f))}$는
$\Gamma_*(\mathcal{F})$를 $\Gamma_*(\mathcal{O}_X)$-가군으로 볼 때의
가군인데, 이는 $\Gamma_*(\mathcal{F})$를 $S$-가군으로 볼 때의
$(\Gamma_*(\mathcal{F}))_{(f)}$와 다름없다. 앞의 표준 동형의 정의
(\hyperref[I.9.3.1-ko]{I, 9.3.1})를 참조하면 이 동형이 준동형
$\beta_f$와 일치함을 알 수 있고, 정리가 따른다.
\end{proof}

\begin{remark}[2.7.6]
\phantomsection
\label{II.2.7.6-ko}
등급환 $S$가 \emph{뇌터}라고 가정하면, 아이디얼 $S_+$가 차수 $1$의 동차
원소 전체의 집합 $S_1$에 의해 생성된다고 가정하는 것만으로
(\hyperref[II.2.7.5-ko]{2.7.5})의 조건이 \emph{자동으로} 성립한다.
\end{remark}

\begin{corollary}[2.7.7]
\phantomsection
\label{II.2.7.7-ko}
(\hyperref[II.2.7.5-ko]{2.7.5})의 가정 아래에서 모든 준연접
$\mathcal{O}_X$-가군 $\mathcal{F}$는 $\widetilde{M}$ 꼴의
$\mathcal{O}_X$-가군과 동형이다. 여기서 $M$은 등급 $S$-가군이다.
\end{corollary}

\begin{corollary}[2.7.8]
\phantomsection
\label{II.2.7.8-ko}
(\hyperref[II.2.7.5-ko]{2.7.5})의 가정 아래에서 모든 유한 생성 준연접
$\mathcal{O}_X$-가군 $\mathcal{F}$는 $\widetilde{N}$ 꼴의
$\mathcal{O}_X$-가군과 동형이다. 여기서 $N$은 유한 생성 등급
$S$-가군이다.
\end{corollary}

\begin{proof}
(\hyperref[II.2.7.7-ko]{2.7.7})에 따라 $\mathcal{F}=\widetilde{M}$이라고
가정할 수 있다. 여기서 $M$은 등급 $S$-가군이다.
$(f_\lambda)_{\lambda\in L}$를 $M$의 동차 생성계라 하자. $L$의 모든 유한
부분집합 $H$에 대하여, $\lambda\in H$인 $f_\lambda$들로 생성되는 $M$의
등급 부분가군을 $M_H$라 하자. $M$이 부분가군 $M_H$들의 귀납극한임은
명백하므로, $\mathcal{F}$는 그 $\mathcal{O}_X$-부분가군
$\widetilde{M_H}$들의 귀납극한이다
(\hyperref[II.2.5.4-ko]{2.5.4}). 그런데 $\mathcal{F}$가 유한 생성이고
바탕공간 $X$가 준콤팩트이므로, (\hyperref[0.5.2.3-ko]{0, 5.2.3})에서
$L$의 어떤 유한 부분집합 $H$에 대하여
$\mathcal{F}=\widetilde{M_H}$임이 따른다.
\end{proof}

\begin{corollary}[2.7.9]
\phantomsection
\label{II.2.7.9-ko}
(\hyperref[II.2.7.5-ko]{2.7.5})의 가정 아래에서 $\mathcal{F}$를 유한 생성
준연접 $\mathcal{O}_X$-가군이라 하자. 그러면 어떤 정수 $n_0$가 존재하여,
모든 $n\geq n_0$에 대하여 $\mathcal{F}(n)$은 $\mathcal{O}_X^k$의 어떤
몫가군과 동형이다. 여기서 $k>0$은 $n$에 의존한다. 따라서
$\mathcal{F}(n)$은 $X$ 위의 유한 개 단면으로 생성된다
(\hyperref[0.5.1.1-ko]{0, 5.1.1}).
\end{corollary}

\begin{proof}
(\hyperref[II.2.7.8-ko]{2.7.8})에 따라 $\mathcal{F}=\widetilde{M}$이라고
가정할 수 있다. 여기서 $M$은 $S(m_i)$ 꼴의 $S$-가군들의 유한 직합의
몫가군이다. 따라서 (\hyperref[II.2.5.4-ko]{2.5.4})에 따라 $M=S(m)$인
경우로 귀착할 수 있고, 이때
$\mathcal{F}(n)=\widetilde{S(m+n)}
=\mathcal{O}_X(m+n)$이다
(\hyperref[II.2.5.15-ko]{2.5.15}). 그러므로 다음 보조정리를 증명하면
충분하다.

\oldpage[II]{41}
\begin{lemma}[2.7.9.1]
\phantomsection
\label{II.2.7.9.1-ko}
(\hyperref[II.2.7.5-ko]{2.7.5})의 가정 아래에서 모든 $n\geq0$에 대하여,
어떤 정수 $k$가 존재하여(이는 $n$에 의존한다) 전사 준동형
$\mathcal{O}_X^k\to\mathcal{O}_X(n)$이 존재한다.
\end{lemma}

실제로 (\hyperref[II.2.7.2-ko]{2.7.2})에 따라, 적당한 $k$에 대하여 등급
$S$-가군 곱 $S^k$에서 $S(n)$으로 가는 \emph{차수 $0$}의
(TN)-전사 준동형 $u$가 존재함을 보이면 충분하다. 그런데
$(S(n))_0=S_n$이고 가정에 따라 모든 $h>0$에 대하여 $S_h=S_1^h$이므로
$SS_n=S_n+S_{n+1}+\cdots$이다. $S_n$은 유한 생성 $S_0$-가군이므로
(\hyperref[II.2.1.5-ko]{2.1.5} 및
\hyperref[II.2.1.6-ko]{2.1.6, (i)}), 이 가군의 생성계
$(a_i)_{1\leq i\leq k}$를 택하자. 준동형 $u$는 $S^k$의 $i$번째 표준 기저
원소를 $a_i$에 대응시킨다($1\leq i\leq k$). 이때
$\operatorname{Coker}u$는
$(S(n))_{-n}+\cdots+(S(n))_{-1}$과 동일시되므로, $u$는 요구한 조건을
충족한다.
\end{proof}

\begin{corollary}[2.7.10]
\phantomsection
\label{II.2.7.10-ko}
(\hyperref[II.2.7.5-ko]{2.7.5})의 가정 아래에서, $\mathcal{F}$를 유한 생성
준연접 $\mathcal{O}_X$-가군이라 하자. 그러면 어떤 정수 $n_0$가 존재하여,
모든 $n\geq n_0$에 대하여 $\mathcal{F}$는
$(\mathcal{O}_X(-n))^k$ 꼴의 $\mathcal{O}_X$-가군의 몫과 동형이다
($k$는 $n$에 의존한다).
\end{corollary}

\begin{proposition}[2.7.11]
\phantomsection
\label{II.2.7.11-ko}
(\hyperref[II.2.7.5-ko]{2.7.5})의 가정이 성립한다고 하고, $M$을 등급
$S$-가군이라 하자. 그러면 다음이 성립한다.
\begin{enumerate}
  \item[\rm(i)] 표준 준동형
    $\widetilde{\alpha}:\widetilde{M}\to
    \widetilde{\Gamma_*(\widetilde{M})}$은 동형이다.
  \item[\rm(ii)] $\mathcal{G}$를 $\widetilde{M}$의 준연접
    $\mathcal{O}_X$-부분가군이라 하고, $N$을 $\alpha$에 의한
    $\Gamma_*(\mathcal{G})$의 역상인 $M$의 등급 $S$-부분가군이라 하자.
    그러면 $\widetilde{N}=\mathcal{G}$이다. 여기서
    (\hyperref[II.2.5.4-ko]{2.5.4})에 따라 $\widetilde{N}$을
    $\widetilde{M}$의 $\mathcal{O}_X$-부분가군으로 동일시한다.
\end{enumerate}
\end{proposition}

\begin{proof}
(\hyperref[II.2.7.5-ko]{2.7.5})에 따라
$\beta:\widetilde{\Gamma_*(\widetilde{M})}\to\widetilde{M}$가 동형이므로,
(\hyperref[II.2.6.5.1-ko]{2.6.5.1})에 따라 $\widetilde{\alpha}$는 그
역동형이다. 이로써 (i)가 따른다. $P$를 $\Gamma_*(\widetilde{M})$의 등급
부분가군 $\alpha(M)$이라 하자. $M\mapsto\widetilde{M}$은 완전함자이므로
(\hyperref[II.2.5.4-ko]{2.5.4}), $\widetilde{\alpha}$에 의한
$\widetilde{M}$의 상은 $\widetilde{P}$와 같다. 따라서 (i)에 따라
$\widetilde{P}=\widetilde{\Gamma_*(\widetilde{M})}$이다. 이제
$Q=\Gamma_*(\mathcal{G})\cap P$로 놓자. 앞의 결과와
(\hyperref[II.2.5.4-ko]{2.5.4})에 따라
$\widetilde{Q}=\widetilde{\Gamma_*(\mathcal{G})}$이므로,
(\hyperref[II.2.7.5-ko]{2.7.5})에 따라 $\beta$를 $\widetilde{Q}$에 제한한
사상은 이 $\mathcal{O}_X$-가군에서 $\mathcal{G}$ 위로 가는
\emph{동형}이다. 한편 $N$의 정의와
(\hyperref[II.2.5.4-ko]{2.5.4})에 따라, 동형 $\widetilde{\alpha}$를
$\widetilde{N}$에 제한한 사상은 $\widetilde{N}$에서
$\widetilde{Q}$ 위로 가는 동형이다. 따라서
(\hyperref[II.2.6.5.1-ko]{2.6.5.1})에 의해 결론이 따른다.
\end{proof}

\subsection{함자적 성질.}
\label{subsection:II.2.8-ko}

\begin{env}[2.8.1]
\phantomsection
\label{II.2.8.1-ko}
$S$, $S'$을 음이 아닌 차수로 등급이 주어진 두 환이라 하고,
$\varphi:S'\to S$를 등급환의 준동형이라 하자. $X=\operatorname{Proj}(S)$에서
$V_+(\varphi(S'_+))$의 여집합인 열린부분집합, 또는 같은 말로 $S'_+$의
동차 원소 $f'$ 전체에 대한 $D_+(\varphi(f'))$들의 합집합을
$G(\varphi)$로 나타내겠다. $\operatorname{Spec}(S)$에서
$\operatorname{Spec}(S')$로 가는 연속사상 ${}^a\varphi$
(\hyperref[I.1.2.1-ko]{I, 1.2.1})를 $G(\varphi)$에 제한하면
$G(\varphi)$에서 $\operatorname{Proj}(S')$로 가는 연속사상이 된다. 표기의
남용으로 이 사상도 ${}^a\varphi$로 나타내겠다. $f'\in S'_+$가 동차
원소이면
\[
\label{II.2.8.1.1-ko}
  {}^a\varphi^{-1}(D_+(f'))=D_+(\varphi(f'))
\tag{2.8.1.1}
\]
이다. 이는 ${}^a\varphi$가 $G(\varphi)$를 $\operatorname{Proj}(S')$로
보낸다는 사실과 (\hyperref[I.1.2.2.2-ko]{I, 1.2.2.2})에서 따른다. 한편
같은 표기 아래에서 준동형 $\varphi$는 등급환의 준동형
$S'_{f'}\to S_f$를 표준적으로 정의하며, 이를 차수 $0$인 원소들에
제한하면 다음 준동형을 얻는다.

\oldpage[II]{42}
$S'_{(f')}\to S_{(f)}$; 이를 $\varphi_{(f)}$로 나타내겠다. 이에 대응하여
(\hyperref[I.1.6.1-ko]{I, 1.6.1}) 아핀 스킴의 사상
$({}^a\varphi_{(f)},\widetilde{\varphi}_{(f)}):
\operatorname{Spec}(S_{(f)})\to\operatorname{Spec}(S'_{(f')})$을 얻는다.
$\operatorname{Spec}(S_{(f)})$를 $D_+(f)$ 위에서
$\operatorname{Proj}(S)$로부터 유도된 스킴과 표준적으로 동일시하면
(\hyperref[II.2.3.6-ko]{2.3.6}), 이로써 사상
$\Phi_f:D_+(f)\to D_+(f')$를 정의하였고, ${}^a\varphi_{(f)}$는
${}^a\varphi$를 $D_+(f)$에 제한한 사상과 동일시된다. 한편 $g'$이
$S'_+$의 두 번째 동차 원소이고 $g=\varphi(g')$이면 다음 도식이
가환임을 곧 알 수 있다.
\[
  \xymatrix{
    D_+(f)\ar[r]^{\Phi_f} & D_+(f')\\
    D_+(fg)\ar[u]\ar[r]_{\Phi_{fg}} & D_+(f'g')\ar[u]
  }
\]
이는 다음 도식이 가환하기 때문이다.
\[
  \xymatrix{
    S'_{(f')}\ar[r]^{\varphi_{(f)}}\ar[d]_{\omega_{f'g',f'}}
    & S_{(f)}\ar[d]^{\omega_{fg,f}}\\
    S'_{(f'g')}\ar[r]_{\varphi_{(fg)}} & S_{(fg)}
  }
\]
$G(\varphi)$의 정의와 (\hyperref[II.2.3.3.2-ko]{2.3.3.2})을 고려하면
따라서 다음을 알 수 있다.
\end{env}

\begin{proposition}[2.8.2]
\phantomsection
\label{II.2.8.2-ko}
등급환 준동형 $\varphi:S'\to S$가 주어지면, 유도 준스킴
$G(\varphi)$에서 $\operatorname{Proj}(S')$로 가는 유일한 사상
$({}^a\varphi,\widetilde{\varphi})$가 존재한다. 이 사상을
\emph{$\varphi$에 대응하는 사상}이라 하고 $\operatorname{Proj}(\varphi)$로
나타낸다. 모든 동차 원소 $f'\in S'_+$에 대하여, 이 사상을
$D_+(\varphi(f'))$에 제한한 사상은 $\varphi$가 정하는 준동형
$S'_{(f')}\to S_{(\varphi(f'))}$에 대응하는 사상과 일치한다.
\end{proposition}

\begin{proof}
앞의 표기에서 $f'\in S'_d$이면 다음 도식이 가환함에 유의하자.
\[
\label{II.2.8.2.1-ko}
  \xymatrix{
    S'_{(f')}\ar[r]^{\varphi_{(f)}}\ar[d]_{\sim}
    & S_{(f)}\ar[d]^{\sim}\\
    {S'}^{(d)}/(f'-1){S'}^{(d)}\ar[r]
    & S^{(d)}/(f-1)S^{(d)}
  }
\tag{2.8.2.1}
\]
여기서 세로 화살표들은
(\hyperref[II.2.2.5-ko]{2.2.5})의 동형들이다.
\end{proof}

\begin{corollary}[2.8.3]
\phantomsection
\label{II.2.8.3-ko}
\begin{enumerate}
  \item[\rm(i)] 사상 $\operatorname{Proj}(\varphi)$는 아핀이다.
  \item[\rm(ii)] $\operatorname{Ker}(\varphi)$가 멱영이면(특히
    $\varphi$가 단사이면), 사상 $\operatorname{Proj}(\varphi)$는
    우세 사상이다.
\end{enumerate}
\end{corollary}

\begin{proof}
(i)은 (\hyperref[II.2.8.2-ko]{2.8.2})와
(\hyperref[II.2.8.1.1-ko]{2.8.1.1})에서 곧바로 따른다. 한편
$\operatorname{Ker}(\varphi)$가 멱영이면, 모든 동차 원소
$f'\in S'_+$에 대하여 $f=\varphi(f')$로 둘 때
$\operatorname{Ker}(\varphi_f)$가 멱영임을 곧바로 확인할 수 있고,
따라서 $\operatorname{Ker}(\varphi_{(f)})$도 멱영이다. 그러므로 결론은
(\hyperref[II.2.8.2-ko]{2.8.2})와
(\hyperref[I.1.2.7-ko]{I, 1.2.7})에서 따른다.
\end{proof}

\oldpage[II]{43}
일반적으로 $\operatorname{Proj}(S)$에서 $\operatorname{Proj}(S')$로 가는
사상들 가운데 아핀이 아닌 것이 있으며, 따라서 이러한 사상은 등급환 준동형
$S'\to S$에서 나오지 않는다. 예를 들어 $A$가 체일 때의 구조 사상
$\operatorname{Proj}(S)\to\operatorname{Spec}(A)$이 그러하다. 여기서는
$\operatorname{Spec}(A)$를 $\operatorname{Proj}(A[T])$와 동일시한다
(\agref{II.3.1.7-ko}{3.1.7} 참조). 실제로 이는
(\hyperref[I.2.3.2-ko]{I, 2.3.2})에서 따른다.

\begin{env}[2.8.4]
\phantomsection
\label{II.2.8.4-ko}
$S''$을 음이 아닌 차수로 등급이 주어진 세 번째 환이라 하고,
$\varphi':S''\to S'$을 등급환 준동형이라 하며,
$\varphi''=\varphi\circ\varphi'$으로 놓자.
(\hyperref[II.2.8.1.1-ko]{2.8.1.1})과 공식
${}^a\varphi''=({}^a\varphi')\circ({}^a\varphi)$에 따라
$G(\varphi'')\subset G(\varphi)$임을 곧바로 확인할 수 있다. 또한
$\Phi$, $\Phi'$, $\Phi''$이 각각 $\varphi$, $\varphi'$, $\varphi''$에
대응하는 사상이면
$\Phi''=\Phi'\circ(\Phi|G(\varphi''))$이다.
\end{env}

\begin{env}[2.8.5]
\phantomsection
\label{II.2.8.5-ko}
$S$가 등급 $A$-대수이고 $S'$이 등급 $A'$-대수라고 하자. 또
$\psi:A'\to A$를 다음 도식이 가환하도록 하는 환 준동형이라 하자.
\[
\label{II.2.8.5.1-ko}
  \xymatrix{
    A'\ar[r]^{\psi}\ar[d] & A\ar[d]\\
    S'\ar[r]_{\varphi} & S
  }
\tag{2.8.5.1}
\]
그러면 $G(\varphi)$와 $\operatorname{Proj}(S')$을 각각
$\operatorname{Spec}(A)$와 $\operatorname{Spec}(A')$ 위의 스킴으로 볼 수
있다. $\Phi$와 $\Psi$가 각각 $\varphi$와 $\psi$에 대응하는 사상이면
다음 도식은 가환한다.
\[
  \xymatrix{
    G(\varphi)\ar[r]^{\Phi}\ar[d] & \operatorname{Proj}(S')\ar[d]\\
    \operatorname{Spec}(A)\ar[r]_{\Psi} & \operatorname{Spec}(A')
  }
\]
실제로 $f=\varphi(f')$이고 $f'$이 $S'_+$의 동차 원소일 때, $\Phi$를
$D_+(f)$에 제한한 사상에 대해 이를 보이면 충분하다. 이 경우에는 다음
도식의 가환성에서 결론이 따른다.
\[
  \xymatrix{
    A'\ar[r]^{\psi}\ar[d] & A\ar[d]\\
    S'_{(f')}\ar[r]_{\varphi_{(f)}} & S_{(f)}
  }
\]
\end{env}

\begin{env}[2.8.6]
\phantomsection
\label{II.2.8.6-ko}
이제 $M$을 등급 $S$-가군이라 하고, 자명하게 등급이 매겨지는
$S'$-가군 $M_{[\varphi]}$를 생각하자. $f'$을 $S'_+$의 동차 원소라 하고
$f=\varphi(f')$이라 하자. 표준 동형
$(M_{[\varphi]})_{f'}\xrightarrow{\sim}(M_f)_{[\varphi_f]}$이 존재함을
알고 있으며(\hyperref[0.1.5.2-ko]{0, 1.5.2}), 이 동형이 차수를
보존함은 자명하다. 따라서 표준 동형
$(M_{[\varphi]})_{(f')}\xrightarrow{\sim}
(M_{(f)})_{[\varphi_{(f)}]}$이 존재한다. 이 동형에 표준적으로 대응하는
층의 동형
$\widetilde{M_{[\varphi]}}|D_+(f')
\xrightarrow{\sim}(\Phi_f)_*(\widetilde{M}|D_+(f))$
이 존재한다
(\hyperref[II.2.5.2-ko]{2.5.2} 및
\hyperref[I.1.6.3-ko]{I, 1.6.3}). 더욱이,

\oldpage[II]{44}
$g'$가 $S'_+$의 두 번째 동차 원소이고 $g=\varphi(g')$이면 다음 도식은
가환한다.
\[
  \xymatrix{
    (M_{[\varphi]})_{(f')}\ar[r]^{\sim}\ar[d]
    & (M_{(f)})_{[\varphi_{(f)}]}\ar[d]\\
    (M_{[\varphi]})_{(f'g')}\ar[r]^{\sim}
    & (M_{(fg)})_{[\varphi_{(fg)}]}
  }
\]
따라서 동형
\[
  \widetilde{M_{[\varphi]}}|D_+(f'g')
  \xrightarrow{\sim}(\Phi_{fg})_*(\widetilde{M}|D_+(fg))
\]
은 동형
$\widetilde{M_{[\varphi]}}|D_+(f')
\xrightarrow{\sim}(\Phi_f)_*(\widetilde{M}|D_+(f))$을
$D_+(f'g')$로 제한한 것임을 곧바로 알 수 있다. $\Phi_f$는 사상
$\Phi$를 $D_+(f)$로 제한한 것이므로,
(\hyperref[II.2.8.1.1-ko]{2.8.1.1})을 고려하고
$X'=\operatorname{Proj}(S')$으로 두면 다음을 얻는다.
\end{env}

\begin{proposition}[2.8.7]
\phantomsection
\label{II.2.8.7-ko}
$\mathcal{O}_{X'}$-가군 $\widetilde{M_{[\varphi]}}$에서
$\mathcal{O}_{X'}$-가군 $\Phi_*(\widetilde{M}|G(\varphi))$로 가는
함자적인 표준 동형이 존재한다.
\end{proposition}

이로부터 등급 $S'$-가군 $M'$에서 등급 $S$-가군 $M$으로 가는
$\varphi$-사상들의 집합에서 $\Phi$-사상
$\widetilde{M'}\to\widetilde{M}|G(\varphi)$들의 집합으로 가는 함자적인
표준 사상을 곧바로 얻는다. (\hyperref[II.2.8.4-ko]{2.8.4})의 표기
아래에서 $M''$가 등급 $S''$-가군이면, $\varphi$-사상 $M'\to M$과
$\varphi'$-사상 $M''\to M'$의 합성에
$\widetilde{M'}|G(\varphi')\to\widetilde{M}|G(\varphi'')$와
$\widetilde{M''}\to\widetilde{M'}|G(\varphi')$의 합성이 표준적으로
대응한다.

\begin{proposition}[2.8.8]
\phantomsection
\label{II.2.8.8-ko}
(\hyperref[II.2.8.1-ko]{2.8.1})의 가정 아래에서 $M'$을 등급
$S'$-가군이라 하자. $(\mathcal{O}_X|G(\varphi))$-가군
$\Phi^*(\widetilde{M'})$에서 $(\mathcal{O}_X|G(\varphi))$-가군
$\widetilde{M'\otimes_{S'}S}|G(\varphi)$로 가는 함자적인 표준 준동형
$\nu$가 존재한다. 아이디얼 $S'_+$가 $S'_1$에 의해 생성되면
$\nu$는 동형이다.
\end{proposition}

\begin{proof}
실제로 $f'\in S'_d$ ($d>0$)에 대하여 다음과 같은
$S_{(f)}$-가군의 함자적인 표준 준동형을 정의한다. 여기서
$f=\varphi(f')$이다.
\[
\label{II.2.8.8.1-ko}
  \nu_f:M'_{(f')}\otimes_{S'_{(f')}}S_{(f)}
  \longrightarrow(M'\otimes_{S'}S)_{(f)}
\tag{2.8.8.1}
\]
이는 준동형
$M'_{(f')}\otimes_{S'_{(f')}}S_{(f)}\to
M'_{f'}\otimes_{S'_{f'}}S_f$와 표준 동형
$M'_{f'}\otimes_{S'_{f'}}S_f\xrightarrow{\sim}(M'\otimes_{S'}S)_f$
(\hyperref[0.1.5.4-ko]{0, 1.5.4})을 합성하여 정의하며, 뒤의 동형이
차수를 보존한다는 점을 사용한다. 두 번째 원소 $g'\in S'_+$와
$g=\varphi(g')$에 대하여 $D_+(f)$에서 $D_+(fg)$로 가는 제한
준동형과 $\nu_f$들이 양립함을 곧바로 확인할 수 있다. 따라서
(\hyperref[I.1.6.5-ko]{I, 1.6.5})을 고려하면 준동형
\[
  \nu:\Phi^*(\widetilde{M'})
  \longrightarrow\widetilde{M'\otimes_{S'}S}|G(\varphi)
\]
가 정의된다. 두 번째 주장을 증명하려면 모든 $f'\in S'_1$에 대하여
$\nu_f$가 동형임을 보이면 충분하다. 실제로 이 경우 $G(\varphi)$는
$D_+(\varphi(f'))$들의 합집합이다. 먼저 $\mathbf{Z}$-쌍선형 사상
$M'_m\times S_n\to M'_{(f')}\otimes_{S'_{(f')}}S_{(f)}$을 정의하자.
이 사상은 $(x',s)$에 원소
$(x'/{f'}^m)\otimes(s/f^n)$을 대응시킨다. 여기서 $m<0$일 때에는
$x'/{f'}^m$이 ${f'}^{-m}x'/1$이라는 약속을 쓴다.

\oldpage[II]{45}
(\hyperref[II.2.5.13-ko]{2.5.13})의 증명에서와 같이 이 사상으로부터
가군의 쌍준동형
\[
  \eta_f:M'\otimes_{S'}S
  \longrightarrow M'_{(f')}\otimes_{S'_{(f')}}S_{(f)}.
\]
을 얻는다. 또한 어떤 $r>0$에 대하여
$f^r\sum_i(x'_i\otimes s_i)=0$이면, 이를
$\sum_i({f'}^rx'_i\otimes s_i)=0$으로도 쓸 수 있다. 따라서
(\hyperref[0.1.5.4-ko]{0, 1.5.4})에 의해
$\sum_i({f'}^rx'_i/{f'}^{m_i+r})\otimes(s_i/f^{n_i})=0$이다. 즉
$\eta_f(\sum_i x'_i\otimes s_i)=0$이므로 $\eta_f$는 다음과 같이
인수분해된다.
\[
  M'\otimes_{S'}S\longrightarrow(M'\otimes_{S'}S)_f
  \xrightarrow{\eta'_f}M'_{(f')}\otimes_{S'_{(f')}}S_{(f)}~;
\]
끝으로 $\eta'_f$와 $\nu_f$가 서로 역인 동형임을 확인할 수 있다.

특히 (\hyperref[II.2.1.2.1-ko]{2.1.2.1})에서 모든
$n\in\mathbf{Z}$에 대하여 표준 준동형
\[
\label{II.2.8.8.2-ko}
  \Phi^*(\mathcal{O}_{X'}(n))\xrightarrow{\sim}
  \mathcal{O}_X(n)|G(\varphi)
\tag{2.8.8.2}
\]
을 얻는다.
\end{proof}

\begin{env}[2.8.9]
\phantomsection
\label{II.2.8.9-ko}
$A$, $A'$을 두 환이라 하고, $\psi:A'\to A$를 환 준동형이라 하자. 이
준동형은 사상
$\Psi:\operatorname{Spec}(A)\to\operatorname{Spec}(A')$를 정의한다.
$S'$를 음이 아닌 차수로 등급이 주어진 $A'$-대수라 하고
$S=S'\otimes_{A'}A$로 두자. 이는 자명하게 $S'_n\otimes_{A'}A$들을
동차 성분으로 갖는 등급 $A$-대수이다. 그러면 사상
$\varphi:s'\mapsto s'\otimes1$은 도식
(\hyperref[II.2.8.5.1-ko]{2.8.5.1})을 가환하게 하는 등급환
준동형이다. 여기서는 $S_+$가 $\varphi(S'_+)$에 의해 생성되는
$A$-가군이므로 $G(\varphi)=\operatorname{Proj}(S)=X$이다. 따라서
$X'=\operatorname{Proj}(S')$로 두면 다음 가환 도식을 얻는다.
\[
\label{II.2.8.9.1-ko}
  \xymatrix{
    X\ar[r]^{\Phi}\ar[d]_p & X'\ar[d]\\
    Y\ar[r]_{\Psi} & Y'
  }
\tag{2.8.9.1}
\]

이제 $M'$을 등급 $S'$-가군이라 하고
$M=M'\otimes_{A'}A=M'\otimes_{S'}S$로 두자. 이 조건 아래에서는
다음이 성립한다.
\end{env}

\begin{proposition}[2.8.10]
\phantomsection
\label{II.2.8.10-ko}
도식 (\hyperref[II.2.8.9.1-ko]{2.8.9.1})은 스킴 $X$를 곱
$X'\times_{Y'}Y$와 동일시한다. 또한 표준 준동형
$\nu:\Phi^*(\widetilde{M'})\to\widetilde{M}$
(\hyperref[II.2.8.8-ko]{2.8.8})은 동형이다.
\end{proposition}

\begin{proof}
첫째 주장은 $S'_+$의 모든 동차 원소 $f'$에 대하여
$f=\varphi(f')$로 둘 때, $\Phi$와 $p$를 $D_+(f)$에 제한한 사상들이
이 스킴을 곱 $D_+(f')\times_{Y'}Y$와 동일시함을 보이면 증명된다
(\hyperref[I.3.2.6.2-ko]{I, 3.2.6.2}). 다시 말해
(\hyperref[I.3.2.2-ko]{I, 3.2.2})에 따라 $S_{(f)}$가
$S'_{(f')}\otimes_{A'}A$와 표준적으로 동일시됨을 보이면 충분하다.
이는 차수를 보존하는 표준 동형
$S_f\xrightarrow{\sim}S'_{f'}\otimes_{A'}A$가 존재하므로 곧바로
따른다(\hyperref[0.1.5.4-ko]{0, 1.5.4}). 둘째 주장은 앞에서 본 바에
따라 $M'_{(f')}\otimes_{S'_{(f')}}S_{(f)}$가
$M'_{(f')}\otimes_{A'}A$와 동일시되고, 후자가 $M_{(f)}$와 동일시되는
사실에서 따른다. 실제로 $M_f$는 차수를 보존하는 동형에 의해
$M'_{f'}\otimes_{A'}A$와 표준적으로 동일시된다.
\end{proof}

\begin{corollary}[2.8.11]
\phantomsection
\label{II.2.8.11-ko}
모든 정수 $n\in\mathbf{Z}$에 대하여 $\widetilde{M(n)}$은
$\Phi^*(\widetilde{M'(n)})
=\widetilde{M'(n)}\otimes_{Y'}\mathcal{O}_Y$와 동일시된다. 특히
$\mathcal{O}_X(n)=\Phi^*(\mathcal{O}_{X'}(n))
=\mathcal{O}_{X'}(n)\otimes_{Y'}\mathcal{O}_Y$이다.
\end{corollary}

\begin{proof}
이는 (\hyperref[II.2.8.10-ko]{2.8.10})과
(\hyperref[II.2.5.15-ko]{2.5.15})에서 따른다.
\end{proof}

\oldpage[II]{46}

\begin{env}[2.8.12]
\phantomsection
\label{II.2.8.12-ko}
(\hyperref[II.2.8.9-ko]{2.8.9})의 가정 아래에서
$f'\in S'_d$ ($d>0$) 및 $f=\varphi(f')$에 대하여 다음 도식은
가환한다.
\[
\label{II.2.8.12.1-ko}
  \xymatrix{
    M'_{(f')}\ar[r]^{\sim}\ar[d]
    & M'^{(d)}/(f'-1)M'^{(d)}\ar[d]\\
    M_{(f)}\ar[r]^{\sim}
    & M^{(d)}/(f-1)M^{(d)}
  }
\tag{2.8.12.1}
\]
(\hyperref[II.2.2.5-ko]{2.2.5} 참조).
\end{env}

\begin{env}[2.8.13]
\phantomsection
\label{II.2.8.13-ko}
(\hyperref[II.2.8.9-ko]{2.8.9})의 표기와 가정을 그대로 두고,
$\mathcal{F}'$를 $\mathcal{O}_{X'}$-가군이라 하자.
$\mathcal{F}=\Phi^*(\mathcal{F}')$로 두면,
(\hyperref[II.2.8.11-ko]{2.8.11})과
(\hyperref[0.4.3.3-ko]{0, 4.3.3})에 의해 모든 $n\in\mathbf{Z}$에
대하여 $\mathcal{F}(n)=\Phi^*(\mathcal{F}'(n))$이다. 따라서
(\hyperref[0.3.7.1-ko]{0, 3.7.1})에 의해 표준 준동형
\[
  \Gamma(\rho):\Gamma(X',\mathcal{F}'(n))
  \longrightarrow\Gamma(X,\mathcal{F}(n))
\]
이 존재하며, 이로부터 등급가군의 표준 쌍준동형
\[
  \Gamma_\bullet(\mathcal{F}')\longrightarrow
  \Gamma_\bullet(\mathcal{F}).
\]
을 얻는다.

아이디얼 $S_+$가 $S_1$에 의해 생성된다고 가정하고,
$\mathcal{F}'=\widetilde{M'}$라 하자. 그러면
$M=M'\otimes_{A'}A$로 둘 때 $\mathcal{F}=\widetilde{M}$이다.
$f'$가 $S'_+$의 동차 원소이고 $f=\varphi(f')$이면,
$M_{(f)}=M'_{(f')}\otimes_{A'}A$임을 이미 보았으며 다음 도식은
가환한다.
\[
  \xymatrix{
    M'_0\ar[r]\ar[d]
    & M'_{(f')}\ar[d]
    & =\Gamma(D_+(f'),\widetilde{M'})\\
    M_0\ar[r]
    & M_{(f)}
    & =\Gamma(D_+(f),\widetilde{M})
  }
\]
이 관찰과 준동형 $\alpha$의 정의
(\hyperref[II.2.6.2-ko]{2.6.2})에서 다음 도식이 가환함을 곧바로
얻는다.
\[
\label{II.2.8.13.1-ko}
  \xymatrix{
    M'\ar[r]^{\alpha_{M'}}\ar[d]
    & \Gamma_\bullet(\widetilde{M'})\ar[d]\\
    M\ar[r]_{\alpha_M}
    & \Gamma_\bullet(\widetilde{M})
  }
\tag{2.8.13.1}
\]
마찬가지로 다음 도식도 가환한다.
\[
\label{II.2.8.13.2-ko}
  \xymatrix{
    \widetilde{\Gamma_\bullet(\mathcal{F}')}\ar[r]^{\beta_{\mathcal{F}'}}
      \ar[d]
    & \mathcal{F}'\ar[d]\\
    \widetilde{\Gamma_\bullet(\mathcal{F})}\ar[r]_{\beta_{\mathcal{F}}}
    & \mathcal{F}
  }
\tag{2.8.13.2}
\]
(오른쪽 세로 화살표는 표준 $\Phi$-사상
$\mathcal{F}'\to\Phi^*(\mathcal{F}')=\mathcal{F}$이다.)
\end{env}

\oldpage[II]{47}

\begin{env}[2.8.14]
\phantomsection
\label{II.2.8.14-ko}
계속해서 (\hyperref[II.2.8.9-ko]{2.8.9})의 가정과 표기법을 유지하고,
$N'$을 또 하나의 등급 $S'$-가군이라 하며
$N=N'\otimes_{A'}A$라 하자. 표준 쌍준동형 $M'\to M$, $N'\to N$은
(표준 환 준동형 $S'\to S$에 상대적인) 쌍준동형
$M'\otimes_{S'}N'\to M\otimes_S N$을 주며, 따라서 차수~$0$의
$S$-준동형
\[
  (M'\otimes_{S'}N')\otimes_{A'}A\longrightarrow M\otimes_S N,
\]
도 준다는 것은 곧바로 알 수 있다. 이에
(\hyperref[II.2.8.10-ko]{2.8.10})을 고려하여 다음과 같은
$\mathcal{O}_X$-가군의 준동형이 대응한다.
\[
\label{II.2.8.14.1-ko}
  \Phi^*(\widetilde{M'\otimes_{S'}N'})
  \longrightarrow\widetilde{M\otimes_S N}.
\tag{2.8.14.1}
\]

더 나아가 다음 도식이 가환함을 곧바로 확인할 수 있다.
\[
\label{II.2.8.14.2-ko}
  \xymatrix{
    \Phi^*(\widetilde{M'}\otimes_{\mathcal{O}_{X'}}\widetilde{N'})
      \ar[r]^{\sim}\ar[d]_{\Phi^*(\lambda)}
    & \widetilde{M}\otimes_{\mathcal{O}_X}\widetilde{N}
      \ar[d]^{\lambda}
    & =\Phi^*(\widetilde{M'})\otimes_{\mathcal{O}_X}
      \Phi^*(\widetilde{N'})\\
    \Phi^*(\widetilde{M'\otimes_{S'}N'})\ar[r]
    & \widetilde{M\otimes_S N}
  }
\tag{2.8.14.2}
\]
여기서 첫째 가로줄은 표준 동형
(\hyperref[0.4.3.3-ko]{0, 4.3.3})이다. 아이디얼 $S'_+$가
$S'_1$에 의해 생성되면, $S_+$가 $S_1$에 의해 생성됨은 명백하다.
이 경우 (\hyperref[II.2.8.14.2-ko]{2.8.14.2})의 두 세로 화살표는
동형이며(\hyperref[II.2.5.13-ko]{2.5.13}), 따라서
(\hyperref[II.2.8.14.1-ko]{2.8.14.1})도 동형이다.

마찬가지로, 차수 $k$의 준동형 $u'$에 역시 차수 $k$인 준동형
$u'\otimes1$을 대응시키면 표준 쌍준동형
$\operatorname{Hom}_{S'}(M',N')\to\operatorname{Hom}_S(M,N)$을 얻는다.
이로부터 다시 차수~$0$의 $S$-준동형
\[
  (\operatorname{Hom}_{S'}(M',N'))\otimes_{A'}A
  \longrightarrow\operatorname{Hom}_S(M,N)
\]
을 얻으며, 따라서 $\mathcal{O}_X$-가군의 준동형
\[
\label{II.2.8.14.3-ko}
  \Phi^*(\widetilde{\operatorname{Hom}_{S'}(M',N')})
  \longrightarrow\widetilde{\operatorname{Hom}_S(M,N)}.
\tag{2.8.14.3}
\]
을 얻는다.

더 나아가 다음 도식은 가환한다.
\[
  \xymatrix{
    \Phi^*(\widetilde{\operatorname{Hom}_{S'}(M',N')})\ar[r]
      \ar[d]_{\Phi^*(\mu)}
    & \widetilde{\operatorname{Hom}_S(M,N)}\ar[d]^{\mu}\\
    \Phi^*(\mathcal{H}om_{\mathcal{O}_{X'}}
      (\widetilde{M'},\widetilde{N'}))\ar[r]
    & \mathcal{H}om_{\mathcal{O}_X}(\widetilde{M},\widetilde{N})
  }
\]
여기서 둘째 가로줄은 표준 준동형
(\hyperref[0.4.4.6-ko]{0, 4.4.6})이다.
\end{env}

\oldpage[II]{48}

\begin{env}[2.8.15]
\phantomsection
\label{II.2.8.15-ko}
표기법과 가정이 (\hyperref[II.2.8.1-ko]{2.8.1})의 것이라 하자.
(\hyperref[II.2.4.7-ko]{2.4.7})에 의하면, $S_0$과 $S'_0$을
$\mathbf{Z}$로 바꾸고 $\varphi_0$을 항등사상으로 바꾼 뒤, 각각 $S$와
$S'$을 $S^{(d)}$와 ${S'}^{(d)}$로 바꾸고($d>0$) $\varphi$를
$S^{(d)}$에 대한 제한 $\varphi^{(d)}$로 바꾸어도 사상 $\Phi$는 동형까지
변하지 않는다.
\end{env}

\subsection{$\operatorname{Proj}(S)$ 스킴의 닫힌 부분스킴.}
\label{subsection:II.2.9-ko}

\begin{env}[2.9.1]
\phantomsection
\label{II.2.9.1-ko}
$\varphi:S\to S'$이 등급환의 준동형이라고 하자. 어떤 정수 $n$이 존재하여
$k\geq n$일 때 $\varphi_k:S_k\to S'_k$가 전사이면(각각 단사이면,
전단사이면), $\varphi$를 (TN)-\emph{전사}(각각 (TN)-\emph{단사},
(TN)-\emph{전단사})라고 한다. 그러면
(\hyperref[II.2.8.15-ko]{2.8.15})에 따라 $\Phi$의 연구는 $\varphi$가
\emph{전사}(각각 \emph{단사}, \emph{전단사})인 경우로 환원된다.
$\varphi$가 (TN)-전단사라고 하는 대신, 이때 이를
(TN)-\emph{동형사상}이라고도 한다.
\end{env}

\begin{proposition}[2.9.2]
\phantomsection
\label{II.2.9.2-ko}
$S$를 음이 아닌 차수로 등급이 주어진 환이라 하고
$X=\operatorname{Proj}(S)$라 하자.
\begin{enumerate}
  \item[\rm(i)] $\varphi:S\to S'$이 등급환의 (TN)-전사 준동형이면, 이에
    대응하는 사상 $\Phi$ (\hyperref[II.2.8.1-ko]{2.8.1})는
    $\operatorname{Proj}(S')$ 전체에서 정의되며
    $\operatorname{Proj}(S')$에서 $X$로 가는 닫힌 몰입 사상이다.
    $\mathfrak{J}$가 $\varphi$의 핵이면, $\Phi$에 결부된 $X$의 닫힌
    부분스킴은 $\mathcal{O}_X$의 준연접 아이디얼층
    $\widetilde{\mathfrak{J}}$로 정의된다.
  \item[\rm(ii)] 더 나아가 아이디얼 $S_+$가 차수~$1$인 유한 개의 동차
    원소로 생성된다고 가정하자. $X'$을 $\mathcal{O}_X$의 준연접
    아이디얼층 $\mathcal{J}$로 정의되는 $X$의 닫힌 부분스킴이라 하자.
    $\mathfrak{J}$를 표준 준동형
    $\alpha:S\to\Gamma_\bullet(\mathcal{O}_X)$
    (\hyperref[II.2.6.2-ko]{2.6.2})에 의한
    $\Gamma_\bullet(\mathcal{J})$의 역상인 $S$의 등급 아이디얼이라 하고,
    $S'=S/\mathfrak{J}$라 놓자. 그러면 $X'$은 등급환의 표준 준동형
    $S\to S'$에 대응하는 닫힌 몰입 사상
    $\operatorname{Proj}(S')\to X$에 결부된 부분스킴이다.
\end{enumerate}
\end{proposition}

\begin{proof}
\begin{enumerate}
  \item[\rm(i)] $\varphi$가 전사라고 가정해도 된다
    (\hyperref[II.2.9.1-ko]{2.9.1}). 가정에 따라 $\varphi(S_+)$가
    $S'_+$를 생성하므로 $G(\varphi)=\operatorname{Proj}(S')$이다.
    한편 둘째 주장은 $X$ 위에서 국소적으로 확인할 수 있다. 따라서
    $f$를 $S_+$의 동차 원소라 하고 $f'=\varphi(f)$라 놓자.
    $\varphi$는 등급환의 전사 준동형이므로
    $\varphi_{(f')}:S_{(f)}\to S'_{(f')}$가 전사이고 그 핵이
    $\mathfrak{J}_{(f)}$임을 곧바로 확인할 수 있다. 이로써 (i)가
    증명된다 (\hyperref[I.4.2.3-ko]{I, 4.2.3}).
  \item[\rm(ii)] (i)에 따라, 표준 단사 준동형
    $j:\mathfrak{J}\to S$에서 유도되는 준동형
    $\widetilde{j}:\widetilde{\mathfrak{J}}\to\mathcal{O}_X$가
    $\widetilde{\mathfrak{J}}$에서 $\mathcal{J}$로 가는 동형임을
    확인하면 충분하다. 이는 (\hyperref[II.2.7.11-ko]{2.7.11})에서
    따른다.
\end{enumerate}
\end{proof}

$\mathfrak{J}$는
$\widetilde{j}(\widetilde{\mathfrak{J}'})=\mathcal{J}$를 만족하는
$S$의 등급 아이디얼 $\mathfrak{J}'$들 가운데 \emph{가장 큰} 것이다.
실제로 정의로 거슬러 올라가 곧바로 확인하면
(\hyperref[II.2.6.2-ko]{2.6.2}), 이 관계에서
$\alpha(\mathfrak{J}')\subset\Gamma_\bullet(\mathcal{J})$가 따른다.

\begin{corollary}[2.9.3]
\phantomsection
\label{II.2.9.3-ko}
(\hyperref[II.2.9.2-ko]{2.9.2}, (i))의 가정이 성립하고, 더 나아가
아이디얼 $S_+$가 $S_1$로 생성된다고 가정하자. 그러면 모든
$n\in\mathbf{Z}$에 대하여 $\Phi^*(\widetilde{S(n)})$은
$\widetilde{S'(n)}$과 표준적으로 동형이고, 따라서 모든
$\mathcal{O}_X$-가군 $\mathcal{F}$에 대하여
$\Phi^*(\mathcal{F}(n))$은 $\Phi^*(\mathcal{F})(n)$과 표준적으로
동형이다.
\end{corollary}

\begin{proof}
이는 (\hyperref[II.2.8.8-ko]{2.8.8})의 특수한 경우이며,
(\hyperref[II.2.5.10.2-ko]{2.5.10.2})를 고려한다.
\end{proof}

\begin{corollary}[2.9.4]
\phantomsection
\label{II.2.9.4-ko}
(\hyperref[II.2.9.2-ko]{2.9.2}, (ii))의 가정이 성립한다고 하자.
$X$의 닫힌 부분준스킴 $X'$이 정역 스킴이기 위한 필요충분조건은 등급
아이디얼 $\mathfrak{J}$가 $S$의 소아이디얼인 것이다.
\end{corollary}

\oldpage[II]{49}

\begin{proof}
$X'$은 $\operatorname{Proj}(S/\mathfrak{J})$와 동형이므로,
(\hyperref[II.2.4.4-ko]{2.4.4})에 따라 조건은 충분하다. 필요성을
보이기 위하여 완전열
$0\to\mathcal{J}\to\mathcal{O}_X\to\mathcal{O}_X/\mathcal{J}$를
생각하자. 이로부터 완전열
\[
  0\longrightarrow\Gamma_\bullet(\mathcal{J})
  \longrightarrow\Gamma_\bullet(\mathcal{O}_X)
  \longrightarrow\Gamma_\bullet(\mathcal{O}_X/\mathcal{J}).
\]
을 얻는다. $f\in S_m$, $g\in S_n$이고
$\Gamma_\bullet(\mathcal{O}_X/\mathcal{J})$에서
$\alpha_{n+m}(fg)$의 상이 영일 때, $\alpha_m(f)$와
$\alpha_n(g)$의 상 가운데 하나가 영임을 보이면 충분하다. 그런데
정의에 따라 이 상들은 정역 스킴 $X'$ 위의 가역
$(\mathcal{O}_X/\mathcal{J})$-가군
$\mathcal{L}=(\mathcal{O}_X/\mathcal{J})(m)$과
$\mathcal{L}'=(\mathcal{O}_X/\mathcal{J})(n)$의 단면들이다. 가정에
의하면 이 두 단면의 곱은 $\mathcal{L}\otimes\mathcal{L}'$에서 영이다
((\hyperref[II.2.9.3-ko]{2.9.3}) 및
(\hyperref[II.2.5.14.1-ko]{2.5.14.1})). 따라서
(\hyperref[I.7.4.4-ko]{I, 7.4.4})에 따라 둘 중 하나가 영이다.
\end{proof}

\begin{corollary}[2.9.5]
\phantomsection
\label{II.2.9.5-ko}
$A$를 환, $M$을 $A$-가군, $S$를 차수~$1$인 동차 원소들의 집합
$S_1$로 생성되는 등급 $A$-대수라 하자. $u:M\to S_1$을 $A$-가군의
전사 준동형이라 하고, $\overline{u}:\mathbf{S}(M)\to S$를 $u$를
확장하는 $M$의 대칭대수 $\mathbf{S}(M)$에서 $S$로 가는
($A$-대수의) 준동형이라 하자. 그러면 $\overline{u}$에 대응하는 사상은
$\operatorname{Proj}(S)$에서 $\operatorname{Proj}(\mathbf{S}(M))$로
가는 닫힌 몰입 사상이다.
\end{corollary}

\begin{proof}
가정에 따라 $\overline{u}$는 전사이므로
(\hyperref[II.2.9.2-ko]{2.9.2})를 적용하면 된다.
\end{proof}
